\pdfinfoomitdate=1
\pdftrailerid{}
\pdfsuppressptexinfo=15
\newcommand{\DoNotLoadEpstopdf}{}
\documentclass[11pt]{article}
\usepackage[T1]{fontenc}
\usepackage[utf8]{inputenc}
\usepackage{lmodern,amsmath,amssymb,amsthm,mathtools,booktabs,array,longtable}
\usepackage[margin=1in]{geometry}
\usepackage{microtype}
\usepackage[hidelinks]{hyperref}
\usepackage{xurl}
\usepackage{enumitem}
\setlist{itemsep=3pt,topsep=5pt}
\allowdisplaybreaks[2]
\emergencystretch=2em
\setcounter{tocdepth}{1}
\newtheorem{theorem}{Theorem}[section]
\newtheorem{lemma}[theorem]{Lemma}
\newtheorem{proposition}[theorem]{Proposition}
\newtheorem{corollary}[theorem]{Corollary}
\theoremstyle{remark}\newtheorem{remark}[theorem]{Remark}
% Macros of the two sources. Names defined by both (\OO, \R, \Z, \Ai, \Bi, \E,
% \Prb, \ceil, \floor) have the same definition in each.
\newcommand{\OO}{\mathcal O}
\newcommand{\R}{\mathbb R}
\newcommand{\Z}{\mathbb Z}
\newcommand{\Ai}{\operatorname{Ai}}
\newcommand{\Bi}{\operatorname{Bi}}
\newcommand{\E}{\mathbb E}
\newcommand{\Prb}{\mathbb P}
\newcommand{\ceil}[1]{\left\lceil #1\right\rceil}
\newcommand{\floor}[1]{\left\lfloor #1\right\rfloor}
% Report 193 (Part II)
\newcommand{\FP}{\operatorname{FP}}
\newcommand{\ii}{\mathrm i}
\newcommand{\FF}{\mathcal F}
\newcommand{\QQ}{\mathcal Q}
\newcommand{\TT}{\mathcal T}
\newcommand{\inv}{\mathcal I}
% Added in the merge (batch 110).
\newcommand{\file}[1]{\nolinkurl{#1}}
\newcommand{\vol}{\mathrm{vol}}
\newenvironment{writenote}[1][]{\par\smallskip\begingroup\small
 \noindent\textbf{[write]}\if\relax\detokenize{#1}\relax\else~\textbf{#1.}\fi\ }%
 {\par\endgroup\smallskip}
\newenvironment{partabstract}{\begin{quote}\small\noindent\textbf{Abstract of this Part (as delivered).}\ }{\end{quote}}
\newcommand{\partsource}[1]{\par\noindent\begin{minipage}{\textwidth}\small #1\end{minipage}\par\medskip}
\hypersetup{pdftitle={Exponential towers (OEIS A096537 and A096542)},pdfauthor={Research report},pdfsubject={Merged report a096537-exponential-towers (batch 110): Airy limits and the corrected asymptotic of A096537 (Report 192); fractional corrections and Gamma laws (Report 193)}}
\title{\textbf{Exponential towers\\ (OEIS A096537 and A096542)}\\[7pt]
\large Airy limits and the corrected asymptotic of A096537;\\ fractional corrections and Gamma laws for the shifted family}
\author{Two research reports of one external research session\\(bundle Reports 192 and 193), merged into one ProveIt report}
\date{Both Parts: 4 October 2026\quad merged: 7 October 2026}
\begin{document}
\hypersetup{pageanchor=false}
\maketitle
\begin{abstract}
Let $F_1(x)=\exp(x\exp(2x\exp(3x\cdots)))=\sum a_nx^n/n!$ (OEIS A096537), $F_y=\exp(xyF_{y+1})$ the shifted family whose coefficients $P_n(y)$ are the rows of A096542, and $\rho=1-e^{-1}$. This report prints two manuscripts, a theorem and its sequel. Part~I proves $a_n/(n!)^2\sim\sqrt{2/(\pi\rho)}\,\rho^{-n}n^{-1/2}$, that is $a_n\sim2\sqrt{2\pi/\rho}\,n^{2n+1/2}e^{-n}(e-1)^{-n}$, through the critical finite-height towers $H_h(1/(eh))\sim Kh^{1/6}e^{-(e-1)h}$, an Airy size law, positive-factorization minor arcs and a uniform lattice local limit; it also proves a Gaussian height law for depth-weighted labelled trees and a refined Lambert inverse. Part~II proves the first correction $a_n\rho^n\sqrt n/(n!)^2=C+c_1n^{-1/3}+o(n^{-1/3})$ with $C=\sqrt{2/(\pi\rho)}$ and $-0.61<c_1<-0.25$, $c_1$ an absolutely convergent Airy integral whose sign is certified by a real reduction and directed interval quadrature, and the Gamma prefactor $P_n(y)/(n!)^2\sim C\rho^{1-y}\rho^{-n}n^{y-3/2}/\Gamma(y)$ for every fixed $y>0$, which answers two of Part~I's open questions. The write records two refuted OEIS formulas, with proofs: the amplitude $2\pi$ conjectured by V\'aclav Kot\v{e}\v{s}ovec in A096537 is wrong (the correct amplitude is larger by the factor $\sqrt{2/(\pi\rho)}=1.00355\ldots$), and the A096542 formula $T(n,1)=n\,a_n$ should read $T(n,1)=n\,a_{n-1}$. Further coefficients, effective constants and onsets, extreme shifts and a complex Borel continuation remain open. The report is unrefereed, and nothing in it is formalized.
\end{abstract}
\hypersetup{pageanchor=true}
{\small\setlength{\parskip}{0pt}\tableofcontents}
\newpage

\section*{Guide to this report}\phantomsection\label{xtw:sec:guide}
\addcontentsline{toc}{section}{Guide to this report}
The two Parts are two manuscripts of the session bundle of Reports 1--243
(arrival commit \texttt{60f54ea06}), placed in this directory by commit
\texttt{8622ca7e5} (batch~110, cluster 110-TOWER).
\begin{center}
\small
\begin{tabular}{@{}l l >{\raggedright\arraybackslash}p{0.50\textwidth} l@{}}
\toprule
Part & Bundle report & Delivered title & Labels\\
\midrule
I & 192 (base) & Airy Limits and the Corrected Asymptotic of A096537: Critical exponential towers and depth weighted labeled trees (4~October 2026) & \texttt{xtw:airy:}\\
II & 193 & Fractional Corrections and Gamma Laws for Exponential Towers: A sequel to the leading Airy asymptotic (4~October 2026) & \texttt{xtw:frac:}\\
\bottomrule
\end{tabular}
\end{center}

\emph{One object, a theorem and its sequel.} Both Parts study Hanna's
continued exponential $F_1(x)=\exp(x\exp(2x\exp(3x\cdots)))$, the
coefficients $a_n=n![x^n]F_1$ of OEIS A096537, and the shifted family
$F_y=\exp(xyF_{y+1})$ whose coefficients $P_n(y)$ are the rows of A096542,
with the same constants $\rho=1-e^{-1}$, $\alpha=e-1$ and the same weighted
labelled trees. Part~I proves the leading equivalent
$a_n/(n!)^2\sim\sqrt{2/(\pi\rho)}\,\rho^{-n}n^{-1/2}$ through critical
finite-height towers, an Airy size law and a uniform lattice local limit; this
\textbf{refutes the amplitude $2\pi$ conjectured in A096537}
(Remark~\ref{xtw:airy:rem:oeis537}). Part~II is an explicit sequel. It cites
Part~I for the leading theorem, restates the leading facts it uses
(Section~\ref{xtw:frac:sec:input}), and answers two items of Part~I's list of
open questions: it proves the first correction,
$b_n\rho^n\sqrt n=C+c_1n^{-1/3}+o(n^{-1/3})$ with $-0.61<c_1<-0.25$
(Theorem~\ref{xtw:frac:thm:correction}), and the Gamma prefactor for every
fixed positive shift (Theorem~\ref{xtw:frac:thm:family}). At $y=1$ its
Theorem~\ref{xtw:frac:thm:family} is Part~I's Theorem~\ref{xtw:airy:thm:main}
with the correction added, and its inverse \eqref{xtw:frac:eq:y1inverse} is
Part~I's $r(Y)$ of \eqref{xtw:airy:eq:Lambertintro} with one further term.

\emph{Arrangement.} The base of the merge is Report~192; its
\file{Report192.tex} was staged as this \file{article.tex}, and it is printed
first, as Part~I, because Report~193 depends on it. Report~193 is printed
second, as Part~II, in full: its proofs of the leading facts are part of a
quantitative argument that also gives the correction, so nothing is removed
as a duplicate; where it reproves a statement of Part~I, a dated note says so.

\emph{Numbering.} Sections are numbered continuously; statements are numbered
within sections and equations continuously through the report, as in both
manuscripts. Part~I keeps Report~192's Sections~1--14, its
equations~(1)--(109) and all its other numbers; the write's two remarks
(Remarks~\ref{xtw:airy:rem:oeis537} and~\ref{xtw:airy:rem:oeis542}) are placed
after the last numbered statement of their sections, so that no delivered
number moves. Report~193's Section~$k$ is Section~$k+14$ here
(Sections~15--28), its statements move with their sections (Report~193's
Theorem~1.1 is Theorem~\ref{xtw:frac:thm:correction}), and its
equation~($k$) is equation~($k+109$). Section~\ref{xtw:sec:further}, which
gathers the open questions of both Parts under Vladimir's standing rule of
4~October 2026, was added in the write.

Notes added in the write are set in small type and tagged \textbf{[write]},
with their date (7~October 2026). No delivered statement was changed and no
delivered symbol was renamed.

\section*{What is proved, and where}\phantomsection\label{xtw:sec:status}
\addcontentsline{toc}{section}{What is proved, and where}
\begin{center}
\small
\begin{tabular}{@{}>{\raggedright\arraybackslash}p{0.58\textwidth} >{\raggedright\arraybackslash}p{0.37\textwidth}@{}}
\toprule
Statement & Where\\
\midrule
\multicolumn{2}{@{}l}{\emph{The leading theory (Part I)}}\\
Exact-height coefficients and depth-weighted labelled trees; $a_n\ge n^2a_{n-1}$ & Section~\ref{xtw:airy:sec:exact}, Lemma~\ref{xtw:airy:lem:monotonicity}\\
$b_n=a_n/(n!)^2\sim\sqrt{2/(\pi\rho)}\,\rho^{-n}n^{-1/2}$, i.e.\ $a_n\sim2\sqrt{2\pi/\rho}\,n^{2n+1/2}e^{-n}(e-1)^{-n}$; \textbf{the OEIS amplitude $2\pi$ is refuted} & Theorem~\ref{xtw:airy:thm:main}, Remark~\ref{xtw:airy:rem:oeis537}\\
Critical tower $H_h(1/(eh))\sim Kh^{1/6}e^{-\alpha h}$; uniform real Airy window & Theorems~\ref{xtw:airy:thm:critical} and~\ref{xtw:airy:thm:window}\\
Airy size law with $p(0)=\Ai(0)^2/\beta$; uniform minor arcs; uniform lattice local limit & Theorems~\ref{xtw:airy:thm:limitsintro}, \ref{xtw:airy:thm:minor}, Section~\ref{xtw:airy:sec:local}\\
Gaussian height limit of depth-weighted trees & Section~\ref{xtw:airy:sec:height}\\
Refined Lambert inverse and shrinking integer sandwich & Theorem~\ref{xtw:airy:thm:inverseintro}, Section~\ref{xtw:airy:sec:inverse}\\
$T(n,1)=n\,a_{n-1}$ in A096542; \textbf{the entry's formula $T(n,1)=n\,a_n$ is refuted} & Section~\ref{xtw:airy:sec:verification}, Remark~\ref{xtw:airy:rem:oeis542}\\
\midrule
\multicolumn{2}{@{}l}{\emph{The sequel (Part II)}}\\
$b_n\rho^n\sqrt n=C+c_1n^{-1/3}+o(n^{-1/3})$, $-0.61<c_1<-0.25$, $c_1$ an absolutely convergent Airy integral & Theorems~\ref{xtw:frac:thm:correction}, \ref{xtw:frac:thm:realintegral}, Section~\ref{xtw:frac:sec:certificate}\\
$b_n(y)\sim C\rho^{1-y}\rho^{-n}n^{y-3/2}/\Gamma(y)$ with the same relative correction, uniformly on compact $Y\subset(0,\infty)$; normal limit of root-ancestor choices & Theorem~\ref{xtw:frac:thm:family}, Section~\ref{xtw:frac:sec:ancestor}\\
Local Gamma correction; positive real Borel boundary; threshold inverses at the fractional scale & Theorem~\ref{xtw:frac:thm:modpoisson}, Sections~\ref{xtw:frac:sec:borel} and~\ref{xtw:frac:sec:inverse}\\
\midrule
\multicolumn{2}{@{}l}{\emph{Open}}\\
\multicolumn{2}{@{}p{0.96\textwidth}@{}}{The next coefficient scale; effective rates, onsets and certified inversion; sharper certified constants; shifts near $0$ or growing; larger tilt domains and rare heights; a complex Borel continuation; the full comparison with Bender--Vinson and Poland. See Section~\ref{xtw:sec:further}.}\\
\bottomrule
\end{tabular}
\end{center}
The finite computations (exact integers, rational checks, and Part~II's
directed interval quadrature for the sign of $c_1$) check identities,
normalizations and one finite enclosure; no asymptotic statement rests on
them, and the interval certificate rests on the analytic reduction and tail
bound proved in Part~II.

\section*{The OEIS entries}\phantomsection\label{xtw:sec:oeis}
\addcontentsline{toc}{section}{The OEIS entries}
The write read both entries on 7~October 2026. \textbf{A096537}
\cite{oeis537} (revision \#8 of 17~December 2019), by Paul D. Hanna
(24~June 2004), has the formula field ``Conjecture: a(n) $\sim$ 2 * Pi *
n\^{}(2*n + 1/2) / (exp(n) * (exp(1) - 1)\^{}n). - Vaclav Kotesovec, Dec 16
2019''. This conjecture is \textbf{false}: Part~I's
Theorem~\ref{xtw:airy:thm:main} gives the same expression with amplitude
$2\sqrt{2\pi/\rho}$, and the two amplitudes differ because $\pi\rho<2$
(Remark~\ref{xtw:airy:rem:oeis537}). The entry's other formula, ``a(n) =
A096542(n+1, 1)/(n+1)'', is right. \textbf{A096542} \cite{oeis542} (revision
\#11 of 19~December 2017), by Hanna (25~June 2004), has the formula line
``T(n, 1) = n*A096537(n)''. This line is \textbf{false} at every $n\ge2$; the
correct statement is $T(n,1)=n\,\mathrm{A096537}(n-1)$, and $n=2$ is a
counterexample: the entry's own row~2 is $0,2,3$, so $T(2,1)=2=2\cdot
\mathrm{A096537}(1)$, while $2\cdot\mathrm{A096537}(2)=10$
(Remark~\ref{xtw:airy:rem:oeis542}). Both manuscripts record this index
mismatch (Sections~\ref{xtw:airy:sec:verification}
and~\ref{xtw:frac:sec:reproducibility}). Nothing was submitted to the OEIS;
whether to propose the corrections is a decision for a human editor. The
OEIS terms in \file{data/192-airy-oeis_fixtures.json} are OEIS content under
CC~BY-SA~4.0; \file{data/192-airy-generated_coefficients_48.json} is
generated by the package, not taken from the OEIS.

\section*{The inverses and the transseries volume}\phantomsection\label{xtw:sec:inverses}
\addcontentsline{toc}{section}{The inverses and the transseries volume}
The repository's transseries volume \cite{tsvol} has a Part~0 apparatus for
inverting rapidly growing functions; its symbols are subscripted ``vol''
here. Statement by statement:
\begin{itemize}
\item Part~I's \eqref{xtw:airy:eq:ceilingexact}, $N(Y)=\lceil z(Y)\rceil$ with
$z$ the inverse of the piecewise-linear interpolation $g$ of $\log a_n$, and
Part~II's \eqref{xtw:frac:eq:exactceiling}, the same for $\log P_n(y)$: instances
of item~(1) of \file{p0:thm:staircase}, with $A_n=\log a_n$ (resp.\ $\log
P_n(y)$) on its eventually increasing part and the admissible interpolation
$g$ (resp.\ $g_y$). The sandwiches \eqref{xtw:airy:eq:sandwichintro} and
\eqref{xtw:frac:eq:ceilingsandwich} follow from item~(1) and the monotonicity
of the ceiling, as item~(2) does; the real-inverse estimates that feed them
are the Parts' own.
\item The Lambert base $x=L/(2W_0(L/(2e\sqrt\rho)))$ of Part~I
\eqref{xtw:airy:eq:Lambertintro} and Part~II \eqref{xtw:frac:eq:lambert} (where it
is $X$), the root of $f(t)=2t\log(t/(e\sqrt\rho))=L$: an \emph{instance} of
the factorial core \file{p0:prop:factorial-core} with $\kappa_\vol=2$,
$d_\vol=-(2+\log\rho)$ and target $L$, which gives
$v_\vol=W_0(L/(2e\sqrt\rho))$ and $X_\vol=e\sqrt\rho\,e^{v_\vol}=L/(2v_\vol)$.
\item The refined inverses \eqref{xtw:airy:eq:inversefinal},
\eqref{xtw:frac:eq:exactinverse}, \eqref{xtw:frac:eq:lambertinverse} and
\eqref{xtw:frac:eq:y1inverse}: first-order (Newton) terms of the volume's
perturbation-operator series \file{p0:eq:operator-series} about the core $f$,
with error bounds proved by the Parts; \emph{analogues}, not instances of a
theorem of the volume.
\end{itemize}
Neither Part claims novelty for the inversion mechanics.

\section*{Notation across the two Parts}\phantomsection\label{xtw:sec:notation}
\addcontentsline{toc}{section}{Notation across the two Parts}
Common to both Parts with the same meaning: $F_y$, $P_n(y)$, $a_n$, $b_n$
($=b_n(1)$ in Part~II), $\rho$, $\alpha$, $c_{n,h}$, $H_h$, $\Ai$, $\Bi$, and
the Riccati objects $w$, $U$ (Part~II's $w_\lambda$, $U_\lambda$ reduce to
Part~I's at $\lambda=0$). No symbol was renamed. The letters below change
meaning:
\begingroup\small
\renewcommand{\theHtable}{notation.\arabic{table}}
\begin{longtable}{@{}>{\raggedright\arraybackslash}p{0.13\textwidth} >{\raggedright\arraybackslash}p{0.83\textwidth}@{}}
\toprule
Symbol & Meanings\\
\midrule
\endhead
$\beta$, $a$ & Part~I: $\beta=2^{-1/3}$. Part~II: $a=2^{-1/3}$ is that constant, and $\beta=y-1$ is the shift offset (Section~\ref{xtw:frac:sec:input}). \emph{False reading:} Part~II's $\beta$ is not $2^{-1/3}$.\\
$K$, $\kappa$ & Part~I: $K=2^{1/6}\sqrt e/(\sqrt\pi\Ai(0)^2)$ the critical constant. Part~II: $\kappa=2^{1/6}\sqrt{e/\pi}$, so $K=\kappa/\Ai(0)^2=\FF(0)$; $K$ an integer cutoff \eqref{xtw:frac:eq:cutoffs}; $K_n$ the root-ancestor statistic.\\
$C$, $c_1$, $d$ & Part~I writes $\sqrt{2/(\pi\rho)}$ without a name. Part~II: $C=\sqrt{2/(\pi\rho)}$, $c_1$ the first correction, $d=c_1/C$; $d_y$ a Newton displacement. Part~I: no $C$.\\
$\varepsilon$, $m$ & Part~I: $\varepsilon=1/(eh)$ the mesh, $m=h^{1/3}$. Part~II: $\epsilon=h^{-1/3}=1/m$. \emph{False reading:} the two epsilons are different scales.\\
$p$, $f$ & Part~I: $p$ the Airy size density \eqref{xtw:airy:eq:pdensity}; $f(t)=2t\log(t/(e\sqrt\rho))$ in Section~\ref{xtw:airy:sec:inverse}. Part~II: $f(z)$ in Section~\ref{xtw:frac:sec:input} is $K\,p(z)$, Part~I's density times the critical constant, since $\FF(\ii t)=K\psi(t)$; so $f(0)=Kp(0)=\kappa/a$; $f$ is the same core as Part~I's in Section~\ref{xtw:frac:sec:inverse}. \emph{False reading:} Part~II's $f$ is not a probability density.\\
$x$, $X$, $r$, $z$ & Part~I: $x$ the Lambert base, $r(Y)$ the refined inverse, $z(Y)$ the interpolated inverse. Part~II: $X$ is that base, $r_y$ a Newton point, $r(\lambda)=\Bi/\Ai$, $z_y(L)$ the interpolated inverse.\\
$L_n$, $N$ & Part~I: $L_n$ the height of the weighted tree; $N_h$ the size variable; $N(Y)$ the threshold. Part~II: $L_n=\log(n/\rho)$; $N_y(Y_0)$ the threshold.\\
$T$, $D$, $M$, $J$, $Q$ & Part~I: $T(z)$ the tree function, $D_i$ differences or factors, $J$ a divided difference or a compact set, $Q_j$ approximants. Part~II: $T=\epsilon M$ and $M$ cutoffs, $D_c(t)$, $D=2\log(X/\sqrt\rho)$, $J_c(\lambda)$, $\QQ(\lambda)$ the correction integrand, $M_y(t)$ the smooth log model.\\
\bottomrule
\end{longtable}
\addtocounter{table}{-1}% the uncaptioned longtable must not consume a table number
\endgroup

\section*{Provenance and merge decisions}\phantomsection\label{xtw:sec:provenance}
\addcontentsline{toc}{section}{Provenance and merge decisions}
Two manuscripts of the session bundle of Reports 1--243, both dated
4~October 2026. Their title blocks read ``Report192'' and ``Report193'', with
the subtitles above in the author slot; both PDF author fields read
``Research report''; neither names an author, a tool or an addressee,
neither pins a repository commit, and neither carries ``prepared for private
review'' wording. They arrived in commit \texttt{60f54ea06} and were placed by
\texttt{8622ca7e5} (batch~110, cluster 110-TOWER).
\begin{itemize}
\item \textbf{Part~I}, bundle Report~192, archive
\file{A096537_Corrected_Asymptotic_and_Airy_Limits_Source.zip} (662,864 bytes;
23 files; \file{Report192.tex}, 1,502 lines, 24 pages); the base.
\item \textbf{Part~II}, bundle Report~193, archive
\file{A096537_Fractional_Corrections_and_Gamma_Laws_Source.zip} (1,008,394
bytes; 37 files; \file{Report193.tex}, 1,965 lines, 29 pages). It cites
Report~192 as \texttt{report192}; that bibliography entry now points to
Part~I. Three of its generated files,
\file{generated/certificate/coefficient.json}, \file{core.json} and
\file{panels.json}, are byte copies of its fixtures
\file{code/fixtures/*_reference.json} and are not shipped.
\end{itemize}
Where the merge had to choose:
\begin{enumerate}
\item \emph{Base and order.} Report~192 first (see the Guide).
\item \emph{Repeated arguments.} Part~II re-derives the leading facts it
uses inside a quantitative argument (the parabolic entrance with the forced
logarithm, the positive factorization on minor arcs, the height averaging);
each proof is needed for its own theorem, so both are printed in full, with
dated notes in Part~II naming the Part~I counterparts.
\item \emph{Front matter, numbering and bibliography.} The two title blocks
and abstracts are printed at the heads of their Parts, with each manuscript's
opening paragraph; one table of contents replaces two. The two
bibliographies are merged: the seven shared keys (\texttt{oeis537},
\texttt{oeis542}, \texttt{bv}, \texttt{poland}, \texttt{dlmf2},
\texttt{dlmf7}, \texttt{dlmf9}) are single entries keeping both Parts'
details, \texttt{report192} points to Part~I, and the write added the entry
for the transseries volume.
\item \emph{Write additions.} This front matter, Remarks~\ref{xtw:airy:rem:oeis537}
and~\ref{xtw:airy:rem:oeis542}, Section~\ref{xtw:sec:further}, the dated notes
and the label prefixes. Everything else is delivered text.
\end{enumerate}

\section*{What was checked, and what was not}\phantomsection\label{xtw:sec:trust}
\addcontentsline{toc}{section}{What was checked, and what was not}
The report is unrefereed and not formalized. At intake (dossier of batch~110,
6~October 2026) both manuscripts were read and no mathematical error was
found. An independent computation of the full-height tower coefficients
reproduced $1,1,5,61,1377,49721,2625313,190735749,18246616321$ and gave
$b_n\rho^n\sqrt n=0.9144$, $0.9462$, $0.9609$, $0.9695$ at $n=100$, $400$,
$1000$, $2000$; least-squares fits $c_0+\sum_jc_jn^{-j/3}$ on the ranges
$n\ge200$, $400$, $800$ gave $c_0$ between $1.00356$ and $1.00363$, against
$\sqrt{2/(\pi\rho)}=1.0035525\ldots$ (the conjectured amplitude corresponds to
$c_0=1$), and $c_1\approx-0.437$, inside Part~II's enclosure; exact rows of
A096542 through $n=8$ confirmed $T(n,1)=n\,a_{n-1}$. These are numerical
observations, not proofs. The delivered suites were rerun on copies (results
in the README). The write (7~October 2026) read both OEIS entries, checked the
inequality $\pi\rho<2$ and the rows of A096542 quoted in
Remark~\ref{xtw:airy:rem:oeis542}, checked the identifications with the
transseries volume above, and reran Part~I's exact checker and Part~II's exact
checker and interval tests on copies. What was \emph{not} done: the analytic
proofs were read and their displayed algebra spot-checked, but they are the
manuscripts' proofs, not an independent verification; Bender--Vinson, Poland
and Chapoton were not read by the write; Part~II's 562-panel interval
certificate was rerun by the intake but not by the write.

\section*{Relation to neighbouring reports and formal projects}\phantomsection\label{xtw:sec:neighbours}
\addcontentsline{toc}{section}{Relation to neighbouring reports and formal projects}
No other placed report treats A096537, A096542, continued exponentials or
depth-weighted labelled trees (repository search, 6--7~October 2026). The
transseries volume \cite{tsvol} is related to the inverses as stated above.
Placement in the collection confers no formal status: no Lean or Rocq file in
the repository treats these sequences, and no statement of this report is
formalized.
\clearpage

\part{Airy limits and the corrected asymptotic of A096537: critical exponential towers and depth weighted labeled trees}\label{xtw:airy:part}
\partsource{Bundle Report~192, archive \file{A096537_Corrected_Asymptotic_and_Airy_Limits_Source.zip}, delivered as \file{Report192.tex} (24 pages), the base of the merge. Delivered title block ``Report192 / Airy Limits and the Corrected Asymptotic of A096537 / Critical exponential towers and depth weighted labeled trees / 4 October 2026''. Printed as delivered, with \textbf{[write]} notes and two \textbf{[write]} remarks; section, statement and equation numbers are those of the delivery. The opening paragraph ``Attribution and scope'' is the manuscript's own.}
\begin{partabstract}
For the formal exponential generating function
$F_1(x)=\exp(x\exp(2x\exp(3x\cdots)))=\sum a_nx^n/n!$, we prove
\[
 \frac{a_n}{(n!)^2}\sim\sqrt{\frac{2}{\pi\rho}}\,
 \rho^{-n}n^{-1/2},\qquad \rho=1-e^{-1}.
\]
Thus the amplitude $2\pi$ in the current OEIS conjecture must be replaced by
$2\sqrt{2\pi/\rho}$. The ratio is $1.00355251532641107955\ldots$.
The proof derives the critical finite-height tower equivalent, including its
$h^{1/6}$ amplitude, from an entrance Riccati equation and a normalized Airy
solution. A uniform real tilt window yields an Airy size law; a separate
positive factorization controls every Fourier minor arc and gives a uniform
lattice local limit. An all-height bound then justifies the final Laplace sum.
We also prove a Gaussian height limit for the associated depth weighted labeled
trees and a refined Lambert-function inverse with a shrinking integer sandwich.
The report supplies self-contained proofs of these steps and exact reproducible
finite checks. No effective asymptotic error, all-orders expansion, or claim of
worldwide priority is made.
\end{partabstract}

\paragraph{Attribution and scope.}
The sequence and the entire shifted functional equation are prior constructions
of Paul D. Hanna recorded in OEIS A096537 and A096542 \cite{oeis537,oeis542}.
General continued exponentials and their weighted tree interpretation are part
of the earlier framework of Bender and Vinson \cite{bv}. The exact theorem scope
of that paper and of Poland's continuation \cite{poland} remains unresolved in
the available full-text comparison. We prove the formulas stated here; we do not
infer their novelty from unsuccessful searches. Section~\ref{xtw:airy:sec:literature}
records the precise limits of the comparison.

\section{Results and proof structure}\label{xtw:airy:sec:results}
All logarithms are natural. Set
\begin{equation}\label{xtw:airy:eq:constants}
 \alpha=e-1,\qquad \rho=\alpha/e=1-e^{-1},\qquad
 \beta=2^{-1/3},\qquad
 K=\frac{2^{1/6}\sqrt e}{\sqrt\pi\,\Ai(0)^2}.
\end{equation}
The Airy functions have their usual normalization as solutions of $y''=zy$.
In particular $\Ai(0)=1/(3^{2/3}\Gamma(2/3))$.
The functional equations
\begin{equation}\label{xtw:airy:eq:shiftfamily}
 F_y(x)=\exp\bigl(xyF_{y+1}(x)\bigr),\qquad F_y(0)=1,
\end{equation}
define a unique family of formal series: the coefficient of degree $n$ on the
right uses coefficients only through degree $n-1$ of the next series. We put
\begin{equation}\label{xtw:airy:eq:anbn}
 a_n=n![x^n]F_1(x),\qquad b_n=\frac{a_n}{(n!)^2}.
\end{equation}
No analytic convergence of the infinite tower is assumed or used. Every
analytic evaluation in the proof is of a finite entire tower.

\begin{theorem}[Corrected coefficient equivalent]\label{xtw:airy:thm:main}
As $n\to\infty$,
\begin{align}
 b_n&\sim\sqrt{\frac{2}{\pi\rho}}\,\rho^{-n}n^{-1/2},\label{xtw:airy:eq:bnmain}\\
 a_n&\sim 2\sqrt{\frac{2\pi}{\rho}}\,
 \frac{n^{2n+1/2}}{e^n(e-1)^n}.\label{xtw:airy:eq:anmain}
\end{align}
Equivalently,
\begin{equation}\label{xtw:airy:eq:logmain}
 \log a_n=2n\log n-(1+\log(e-1))n
       +\frac12\log n+\log\!\left(2\sqrt{\frac{2\pi}{\rho}}\right)+o(1).
\end{equation}
\end{theorem}

The retrieved A096537 entry attributes to V\'aclav Kote\v{s}ovec a conjecture,
dated 16 December 2019, with the same expression as \eqref{xtw:airy:eq:anmain} but
amplitude $2\pi$ \cite{oeis537}. The corrected-to-conjectured ratio is
\begin{equation}\label{xtw:airy:eq:ratio}
 \sqrt{\frac{2}{\pi\rho}}=1.00355251532641107955\ldots>1.
\end{equation}
The strict inequality does not depend on decimal computation: $e<11/4$ gives
$\rho<7/11$, whereas $\pi<22/7$, so $\pi\rho<2$. Thus
Theorem~\ref{xtw:airy:thm:main} disproves that particular amplitude conjecture.

For $h\ge1$, define
\begin{equation}\label{xtw:airy:eq:towers}
 A_{h+1,h}(x)=1,\qquad A_{i,h}(x)=\exp\bigl(ixA_{i+1,h}(x)\bigr)
 \quad(1\le i\le h),
\end{equation}
and write $A_{1,0}=1$ and $H_h=A_{1,h}-A_{1,h-1}$. These entire functions
have nonnegative Taylor coefficients, with $H_h$ starting in degree $h$.
Theorem~\ref{xtw:airy:thm:window} below gives, uniformly for $\lambda$ in each compact
subset of $(-\infty,1]$,
\begin{equation}\label{xtw:airy:eq:windowintro}
 H_h\!\left(\frac{1+\lambda h^{-2/3}}{eh}\right)
 \sim \frac{2^{1/6}\sqrt e}{\sqrt\pi\,\Ai(-\beta\lambda)^2}
       h^{1/6}e^{-\alpha h+\alpha\lambda h^{1/3}}.
\end{equation}
The upper endpoint $1$ is a convenient sufficient restriction, not an assertion
about the largest possible parameter domain.

\begin{theorem}[Two different fluctuation laws]\label{xtw:airy:thm:limitsintro}
Define $c_{n,h}=[x^n]H_h(x)$ and an integer variable $N_h$ by
\begin{equation}\label{xtw:airy:eq:sizeprobintro}
 \Prb(N_h=n)=\frac{c_{n,h}(eh)^{-n}}{H_h(1/(eh))}.
\end{equation}
Then $(N_h-\alpha h)/h^{2/3}$ converges to a law with bounded continuous density
$p$ and characteristic function
\begin{equation}\label{xtw:airy:eq:psiintro}
 \psi(t)=\frac{\Ai(0)^2}{\Ai(-\beta\mathrm i t)^2},\qquad
 p(0)=\frac{\Ai(0)^2}{\beta}.
\end{equation}
Moreover,
\begin{equation}\label{xtw:airy:eq:lltintro}
 \sup_{n\in\Z}\left|h^{2/3}\Prb(N_h=n)
 -p\!\left(\frac{n-\alpha h}{h^{2/3}}\right)\right|\longrightarrow0.
\end{equation}
Separately, let $L_n$ be the height of a labeled tree on $\{0,\ldots,n\}$,
rooted at $0$, chosen with probability proportional to the product of all
nonroot depths. Then
\begin{equation}\label{xtw:airy:eq:cltintro}
 \frac{L_n-n/\alpha}{\sqrt n/\alpha}\ \xrightarrow{d}\ \mathcal N(0,1).
\end{equation}
\end{theorem}

The finite-height amplitude $h^{1/6}$, the size fluctuation scale $h^{2/3}$,
and the fixed-size height fluctuation scale $\sqrt n$ describe different objects.
They must not be identified with one another.

\begin{theorem}[Refined threshold inverse]\label{xtw:airy:thm:inverseintro}
Let $Y\to\infty$, $L=\log Y$, and
\begin{equation}\label{xtw:airy:eq:Lambertintro}
 x=\frac{L}{2W_0\!\left(L/(2e\sqrt\rho)\right)},\qquad
 r(Y)=x-\frac14-\frac{\log(8\pi/\sqrt\rho)}{4\log(x/\sqrt\rho)}.
\end{equation}
Here $W_0$ is the positive real Lambert function. The inverse $z(Y)$ of the
continuous piecewise-linear interpolation of $\log a_n$, on its eventually
strictly increasing portion, satisfies
\begin{equation}\label{xtw:airy:eq:zintro}
 z(Y)=r(Y)+o(1/\log x).
\end{equation}
For $N(Y)=\min\{n:a_n\ge Y\}$ there exists $\epsilon(Y)>0$ tending to zero such
that, for all sufficiently large $Y$,
\begin{equation}\label{xtw:airy:eq:sandwichintro}
 \ceil{r(Y)-\frac{\epsilon(Y)}{\log x}}
 \le N(Y)\le
 \ceil{r(Y)+\frac{\epsilon(Y)}{\log x}}.
\end{equation}
This is an existence statement for a vanishing error envelope, not an effective
finite-target bound or a universally valid formula $N(Y)=\ceil{r(Y)}$.
\end{theorem}

\subsection{Dependencies and the role of each argument}
The proof has the following order.
\begin{enumerate}
\item Section~\ref{xtw:airy:sec:exact} identifies the exact-height coefficients and the
weighted-tree normalization, including the two factorials in $b_n$.
\item Sections~\ref{xtw:airy:sec:entrance}--\ref{xtw:airy:sec:critical} derive the critical
entrance and its sensitivity, prove bulk product matching, and obtain $K$.
\item Sections~\ref{xtw:airy:sec:window}--\ref{xtw:airy:sec:airyprob} establish the uniform real
window, complex characteristic-function limit by analytic continuation, and
$p(0)$ by a Wronskian integral.
\item Section~\ref{xtw:airy:sec:minor} proves Fourier-tail bounds independently of the
Airy approximation. Section~\ref{xtw:airy:sec:local} then proves the uniform lattice
local limit. Weak convergence alone would not suffice here.
\item Section~\ref{xtw:airy:sec:sum} obtains an all-height coefficient majorant and
performs the height sum. This is where the Airy constants cancel.
\item Sections~\ref{xtw:airy:sec:height} and \ref{xtw:airy:sec:inverse} derive the height central
limit theorem and the refined inverse.
\end{enumerate}
Every limiting matching argument first takes $h\to\infty$ at a fixed positive
matching time, then sends that time to infinity. No unstated uniformity along a
growing matching time is used. Exact computations in Section~\ref{xtw:airy:sec:verification}
check identities and normalizations; they do not replace any asymptotic proof.

\begin{remark}[\textsf{[write]} The amplitude conjectured in A096537 is false]\label{xtw:airy:rem:oeis537}
[Added 7 October 2026, batch~110.] The formula field of OEIS A096537
(revision \#8, 17~December 2019, live on 7~October 2026) reads ``Conjecture:
a(n) $\sim$ 2 * Pi * n\^{}(2*n + 1/2) / (exp(n) * (exp(1) - 1)\^{}n). -
Vaclav Kotesovec, Dec 16 2019''. This conjecture is false. By
Theorem~\ref{xtw:airy:thm:main},
\[
 \frac{a_n\,e^n(e-1)^n}{n^{2n+1/2}}\longrightarrow2\sqrt{\frac{2\pi}{\rho}}
 =2\pi\sqrt{\frac{2}{\pi\rho}},
\]
whereas the conjecture asserts the limit $2\pi$. The two limits differ because
$\pi\rho<2$: from $e<11/4$, $\rho=1-1/e<1-4/11=7/11$, and $\pi<22/7$, so
$\pi\rho<2$, as shown after \eqref{xtw:airy:eq:ratio}. So the conjectured form
$n^{2n+1/2}e^{-n}(e-1)^{-n}$ is right and only its amplitude is wrong, by the
factor $\sqrt{2/(\pi\rho)}=1.00355251532641107954\ldots$ (the abstract and
\eqref{xtw:airy:eq:ratio} print $\ldots07955\ldots$, a rounding; the
truncation is $\ldots07954$). The intake's independent computation of $a_n$
to $n=2000$, with fits in powers of $n^{-1/3}$, extrapolated
$a_n\rho^n\sqrt n/(n!)^2$ to $1.0036$, not to $1$. The approach is slow because
of the first correction $c_1n^{-1/3}$ proved in Part~II
(Theorem~\ref{xtw:frac:thm:correction}): with Part~II's diagnostic value
$c_1\approx-0.4365$, the normalized value at $n=2000$ is about $0.969$, about
$3.5\,\%$ below its limit, which is why a short-range fit could not have
distinguished the two amplitudes. No correction has been submitted to the
OEIS.
\end{remark}

\section{Exact coefficients and weighted trees}\label{xtw:airy:sec:exact}
Let $\mathcal T_n$ be the trees with vertex set $\{0,1,\ldots,n\}$ and root $0$.
For $T\in\mathcal T_n$, let $d_T(v)$ be the distance from $0$ to $v$ and set
\begin{equation}\label{xtw:airy:eq:treeweight}
 w_y(T)=\prod_{v=1}^n(y+d_T(v)-1),\qquad w(T)=w_1(T).
\end{equation}
The root has no weight. Decomposing the unordered collection of root branches
and using the labeled exponential formula gives
\begin{equation}\label{xtw:airy:eq:treeidentity}
 F_y(x)=\sum_{n\ge0}\frac{x^n}{n!}\sum_{T\in\mathcal T_n}w_y(T).
\end{equation}
Indeed a branch has a root carrying $xy$ and descendants described by
$F_{y+1}$; the unordered set of branches gives $\exp(xyF_{y+1})$.
The same argument with maximum depth at most $h$ gives $A_{1,h}$ at $y=1$.
Consequently
\begin{equation}\label{xtw:airy:eq:heightidentity}
 n!c_{n,h}=\sum_{\substack{T\in\mathcal T_n\\\operatorname{height}(T)=h}}w(T),
 \qquad a_n=\sum_{h=1}^n n!c_{n,h}\quad(n\ge1).
\end{equation}
In particular
\begin{equation}\label{xtw:airy:eq:bdecomp}
 b_n=\sum_{h=1}^n\frac{c_{n,h}}{n!}.
\end{equation}

These conclusions can also be proved without combinatorics. At the bottom,
$A_{h,h}-1=e^{hx}-1$ has positive coefficients and least degree one. If the
successive difference has nonnegative coefficients and least degree $j$,
then the identity
$e^U-e^V=e^V(e^{U-V}-1)$ shows that the difference one level above has
nonnegative coefficients and least degree $j+1$. Thus $H_h$ starts in degree
$h$ and the telescoping sum is finite in every coefficient. This proves the
formal identity even though $F_1$ has zero ordinary radius of convergence:
a chain alone gives $[x^n]F_1\ge n!$.

For computational purposes, if
$B(x)=\sum B_nx^n/n!$ and $A(x)=\exp(ixB(x))=\sum A_nx^n/n!$, differentiating
$A$ yields the exact integer recurrence
\begin{equation}\label{xtw:airy:eq:coefficientrec}
 A_0=1,\qquad
 A_n=\sum_{k=1}^n\binom{n-1}{k-1}kiB_{k-1}A_{n-k}.
\end{equation}
Starting with $B_0=1$ and $B_n=0$ for $n>0$, apply this from level $h$ down
to $1$. Height $n$ already gives the stabilized coefficient $a_n$.
The initial values are
\[
 1,\ 1,\ 5,\ 61,\ 1377,\ 49721,\ 2625313,\ 190735749.
\]

\begin{lemma}[Monotonicity]\label{xtw:airy:lem:monotonicity}
For $n\ge1$, $a_n\ge n^2a_{n-1}$, and hence $b_n\ge b_{n-1}$.
In particular $a_n$ is strictly increasing for $n\ge2$.
\end{lemma}
\begin{proof}
Write $P_j(y)=j![x^j]F_y$. Restricting the root to one child gives
$a_n\ge nP_{n-1}(2)$. In any tree with $n-1$ nonroot vertices, arrange its
nonroot depths increasingly as $d_1\le\cdots\le d_{n-1}$. Then $d_j\le j$:
a vertex at depth $d$ has ancestors at each depth $1,\ldots,d-1$. Thus
\[
 \prod_{j=1}^{n-1}\frac{d_j+1}{d_j}
 \ge\prod_{j=1}^{n-1}\frac{j+1}{j}=n.
\]
Apply this to every tree in $P_{n-1}(2)$ to obtain
$P_{n-1}(2)\ge nP_{n-1}(1)$. Empty products cover $n=1$.
\end{proof}

\section{Critical entrance and sensitivity}\label{xtw:airy:sec:entrance}
Throughout Sections~\ref{xtw:airy:sec:entrance}--\ref{xtw:airy:sec:critical}, put
$x_h=1/(eh)$ and $m=h^{1/3}$. At this real argument let
\[
 q_i^A=\log A_{i,h}(x_h),\qquad q_i^B=\log A_{i,h-1}(x_h),
\]
where $q_h^A=1/e$ and $q_h^B=0$. Define
$u_k^A=1-q_{h-k}^A$ and $u_k^B=1-q_{h-k}^B$, $0\le k\le h-1$.
Both obey the exact recurrence
\begin{equation}\label{xtw:airy:eq:urec}
 u_k=1-(1-k/h)e^{-u_{k-1}},\qquad
 u_0^A=\rho,\quad u_0^B=1.
\end{equation}
Here $0<u_k\le1$. With $\phi(u)=1-e^{-u}$, the terminal states have the
important relation $u_0^A=\phi(u_0^B)$. It will fix the sensitivity constant.

Let $w$ solve
\begin{equation}\label{xtw:airy:eq:wode}
 w''(t)=\tfrac12tw(t),\qquad w(0)=0,\quad w'(0)=1,
\end{equation}
and put $U=2w'/w$. Positive initial slope and $w''\ge0$ while $w\ge0$ imply
$w,w'>0$ for $t>0$. Direct differentiation and a Taylor series at zero give
\begin{equation}\label{xtw:airy:eq:riccati}
 U'=t-\tfrac12U^2,\qquad U(t)=2/t+\OO(t^2)\quad(t\downarrow0).
\end{equation}

\begin{lemma}[Entrance with normalized sensitivity]\label{xtw:airy:lem:entrance}
Uniformly on each compact interval $0<d\le t\le L<\infty$,
\begin{align}
 mu_{\floor{mt}}^A&\longrightarrow U(t),&
 mu_{\floor{mt}}^B&\longrightarrow U(t),\label{xtw:airy:eq:entrance}\\
 m^2\bigl(u_{\floor{mt}}^B-u_{\floor{mt}}^A\bigr)
 &\longrightarrow\frac{2}{w(t)^2}.&&\label{xtw:airy:eq:sensitivity}
\end{align}
\end{lemma}
\begin{proof}
For fixed $a\in(0,1]$, let $v_k=\phi^k(a)$. This sequence decreases to zero and
\[
 v_{k+1}=v_k-\tfrac12v_k^2+\OO(v_k^3),\qquad
 \frac1{v_{k+1}}-\frac1{v_k}\longrightarrow\frac12.
\]
Ces\`aro summation gives $kv_k\to2$. Since $\phi$ is increasing and
$1$-Lipschitz on $[0,1]$, comparison with \eqref{xtw:airy:eq:urec}, using the same initial
state, gives
\begin{equation}\label{xtw:airy:eq:autocompare}
 0\le u_k-v_k\le\sum_{j=1}^k j/h=\frac{k(k+1)}{2h}.
\end{equation}
For a small fixed $d>0$, this implies
\begin{equation}\label{xtw:airy:eq:entrysmall}
 mu_{\floor{dm}}=2/d+\OO(d^2)+o_h(1).
\end{equation}
On $d\le k/m\le L$, the scaled values $V_k=mu_k$ are bounded by the same
comparison. Taylor expansion of the recurrence is uniform there and gives
\begin{equation}\label{xtw:airy:eq:euler}
 V_k-V_{k-1}=m^{-1}\bigl(k/m-V_{k-1}^2/2\bigr)+\OO_{d,L}(m^{-2}).
\end{equation}
For fixed $d$, discrete Gronwall therefore compares this Euler scheme uniformly
on $[d,L]$ to the exact ODE flow from its initial value. Two nonnegative
Riccati solutions have difference $D$ satisfying
$D'=-(U_1+U_2)D/2$, so the absolute difference contracts. Compare the flow
started in \eqref{xtw:airy:eq:entrysmall} to \eqref{xtw:airy:eq:riccati}; take $h\to\infty$ and
then $d\downarrow0$. The discrepancy is $\OO(d^2)$ and tends to zero. This
proves \eqref{xtw:airy:eq:entrance}, including uniformity and harmless grid rounding.

For sensitivity use $v_k=\phi^k(1)$ and the autonomous pair $v_k,v_{k+1}$.
Set $\Delta_k=u_k^B-u_k^A$ and $\delta_k=v_k-v_{k+1}$. Their initial
differences agree. The exact divided difference
\begin{equation}\label{xtw:airy:eq:divdiff}
 J(a,b)=\frac{\phi(b)-\phi(a)}{b-a}
       =\int_0^1e^{-(1-r)a-rb}\,dr
\end{equation}
has the following property: if both endpoints change by at most $E$, the
logarithm of $J$ changes by at most $E$. Indeed each integrand changes by a
factor between $e^{-E}$ and $e^E$. Iterating the differences gives
\[
 \Delta_k=\Delta_0\prod_{j=1}^k(1-j/h)J(u_{j-1}^A,u_{j-1}^B),\qquad
 \delta_k=\delta_0\prod_{j=1}^kJ(v_j,v_{j-1}).
\]
Use \eqref{xtw:airy:eq:autocompare} on each endpoint. Uniformly for $k\le dm$ and large
$h$,
\begin{equation}\label{xtw:airy:eq:deltaentry}
 \left|\log\frac{\Delta_k}{\delta_k}\right|
 \le C(k^2/h+k^3/h)=\OO(d^3)+o_h(1).
\end{equation}
Since $kv_k\to2$, Taylor expansion of $v_k-\phi(v_k)$ gives
$k^2\delta_k\to2$, so
\begin{equation}\label{xtw:airy:eq:smallDelta}
 m^2\Delta_{\floor{dm}}=\frac{2}{d^2}
       \exp\{\OO(d^3)+o_h(1)\}.
\end{equation}
On $d\le j/m\le L$, \eqref{xtw:airy:eq:entrance} and \eqref{xtw:airy:eq:divdiff} give
\[
 \log J(u_{j-1}^A,u_{j-1}^B)=-U(j/m)/m+o_h(1/m)
\]
uniformly. Also $\sum_{j\le Lm}|\log(1-j/h)|=\OO_L(m^2/h)=o(1)$.
Hence, for fixed $d\le t\le L$,
\begin{equation}\label{xtw:airy:eq:sensitivityprop}
 \frac{\Delta_{\floor{tm}}}{\Delta_{\floor{dm}}}
 \longrightarrow\exp\!\left(-\int_d^tU(r)\,dr\right)
 =\frac{w(d)^2}{w(t)^2}.
\end{equation}
Now $w(d)=d+\OO(d^4)$. Let $d\downarrow0$ in
\eqref{xtw:airy:eq:smallDelta}--\eqref{xtw:airy:eq:sensitivityprop} to prove
\eqref{xtw:airy:eq:sensitivity}. The local logarithmic error used here is
$o_h(1/m)$, which is sufficient after summing; an unproved $\OO(m^{-2})$
rate of entrance convergence is not needed.
\end{proof}

\section{Bulk matching and the endpoint product}\label{xtw:airy:sec:bulk}
Let $T(z)$ be the small real solution of $Te^{-T}=z$, for $0\le z\le1/e$.
In particular $T(0)=0$ and $T(1/e)=1$. For $z<1/e$ define
\begin{equation}\label{xtw:airy:eq:Rdef}
 R(z)=\frac{T(z)e^{T(z)}}{(1-T(z))^2}.
\end{equation}
We use an approximate solution with its first adiabatic correction; retaining
this correction is necessary to recover the amplitude.

\begin{lemma}[Bulk product matching]\label{xtw:airy:lem:bulk}
For $i=h-\floor{tm}$, write $\varepsilon=1/(eh)$,
$z_j=j\varepsilon$, $T_j=T(z_j)$, and $v_i=1-T_i$. For either of the two
critical towers,
\begin{equation}\label{xtw:airy:eq:bulk}
 \limsup_{h\to\infty}\left|
 \log\!\left(v_i\prod_{j=1}^{i-1}\frac{q_j}{T_j}\right)\right|
 =\OO(t^{-3/2})\qquad(t\to\infty).
\end{equation}
\end{lemma}
\begin{proof}
Put $Q_j=T_j+\varepsilon R(z_j)$. Differentiation gives
\begin{align*}
 T'&=e^T/(1-T),&T''&=e^{2T}(2-T)/(1-T)^3,\\
 R'&=e^{2T}(1+2T-T^2)/(1-T)^4,&R&=T(T'+R).
\end{align*}
The last identity cancels the first-order residual. Taylor's theorem yields
\begin{equation}\label{xtw:airy:eq:residual}
 E_j=Q_j-z_je^{Q_{j+1}},\qquad
 |E_j|\le C\varepsilon^2\frac{T_j}{(1-T_j)^4}\quad(j<i),
\end{equation}
provided $h v_i^3$ exceeds an absolute constant. Here is the relevant uniformity.
Writing $v_j=1-T_j$, neighboring gaps satisfy
$v_j/v_{j+1}\le1+C\varepsilon/v_i^2$. Taylor bounds on the neighboring interval
may therefore use $v_j$. The expansion of $Q_{j+1}$ is
$T_j+\varepsilon(T'_j+R_j)+\OO(\varepsilon^2v_j^{-4})$;
$z_je^{T_j}=T_j$, and the remaining exponential quadratic term is also
$\OO(\varepsilon^2T_jv_j^{-4})$. This proves \eqref{xtw:airy:eq:residual}, including
its factor $T_j$ at the left endpoint.

The Airy function and derivative asymptotics, written explicitly in
Section~\ref{xtw:airy:sec:critical}, imply
\begin{equation}\label{xtw:airy:eq:Ularge}
 U(t)=\sqrt{2t}-\frac1{2t}+\OO(t^{-5/2}).
\end{equation}
The turning-point expansion of $Te^{-T}=z$ and Lemma~\ref{xtw:airy:lem:entrance} then give
\begin{equation}\label{xtw:airy:eq:boundaryQ}
 mv_i\longrightarrow\sqrt{2t},\qquad
 m(q_i-Q_i)=\OO(t^{-5/2})+o_h(1).
\end{equation}
For fixed sufficiently large $t$, all $Q_j$ with $j\le i$ are at most
$1-cv_i$: both $T$ and $R$ increase and
$\varepsilon R(z_i)/v_i=\OO(t^{-3/2})$.
For the true tower use $M=\max(q_i,T_i)<1$.
On $[0,1]$ the function $u e^{-u}$ increases, so $z_i e^M\le M$; descending
induction gives $q_j\le M\le1-cv_i$ for $j\le i$.
The same bounds for both arguments in the mean-value theorem show
\[
 |q_j-Q_j|\le r|q_{j+1}-Q_{j+1}|+|E_j|,\qquad r\le1-c'v_i.
\]
For example the derivative is at most $z_i e^{M'}$ with
$M'=\max(q_i,Q_i,T_i)\le1-cv_i$; its logarithm is at most $-cv_i$.
Summing the geometric propagation therefore gives
\begin{equation}\label{xtw:airy:eq:sumdefect}
 \sum_{j=1}^i|q_j-Q_j|
 \le\frac C{v_i}\left(|q_i-Q_i|+\sum_{j=1}^{i-1}|E_j|\right).
\end{equation}
The monotonicity of the summand in \eqref{xtw:airy:eq:residual}, followed by the
substitution $dz=e^{-T}(1-T)dT$, gives
\[
 \sum_{j<i}|E_j|\le C\bigl(h^{-1}v_i^{-2}+h^{-2}v_i^{-4}\bigr).
\]
Indeed its integral is bounded by $C v_i^{-2}$ and the last mesh cell costs
$C\varepsilon^2v_i^{-4}$. Equations \eqref{xtw:airy:eq:boundaryQ} and
\eqref{xtw:airy:eq:sumdefect} now give
\begin{equation}\label{xtw:airy:eq:sumdefectlimit}
 \limsup_{h\to\infty}\sum_{j=1}^i|q_j-Q_j|=\OO(t^{-3/2}).
\end{equation}

To convert this additive control to a product without dividing by small $T_j$,
use the exact identity
\[
 \log(q_j/T_j)=q_{j+1}-T_j.
\]
Thus
\[
 \sum_{j=1}^{i-1}\log(q_j/T_j)
 =\sum_{j=2}^i(q_j-T_j)+T_i-T_1.
\]
Replace $q_j-T_j$ by $\varepsilon R(z_j)$ using
\eqref{xtw:airy:eq:sumdefectlimit}. The Riemann-sum error is
$\OO(\varepsilon R(z_i))=o_h(1)$ for fixed $t$, and direct substitution gives
\begin{equation}\label{xtw:airy:eq:Rintegral}
 \int_0^{z_i}R(z)\,dz=-T_i-\log(1-T_i).
\end{equation}
As $T_1\to0$, these formulas prove \eqref{xtw:airy:eq:bulk}.
\end{proof}

\begin{lemma}[Endpoint Riemann formula]\label{xtw:airy:lem:endpoint}
If $f$ is absolutely continuous on $[0,1]$ with $f'\in L^1[0,1]$, then
\begin{equation}\label{xtw:airy:eq:riemann}
 \sum_{j=1}^h f(j/h)-h\int_0^1f(t)\,dt
 \longrightarrow\frac{f(1)-f(0)}2.
\end{equation}
\end{lemma}
\begin{proof}
Integrating over each mesh cell makes the difference exactly
$\int_0^1 f'(t)\{ht\}\,dt$, where braces denote fractional part; grid points
have measure zero. The bounded functions $\{ht\}$ converge weakly against
$L^1$ functions to $1/2$. Prove this first for interval indicators by summing
whole mesh cells, then for step functions, and finally use $L^1$
approximation and $0\le\{ht\}\le1$.
\end{proof}

\begin{lemma}[Critical equilibrium product]\label{xtw:airy:lem:product}
With $T_j=T(j/(eh))$,
\begin{equation}\label{xtw:airy:eq:fullproduct}
 \prod_{j=1}^hT_j\sim\sqrt{2\pi e h}\,e^{-\alpha h}.
\end{equation}
For fixed $t>0$ and $i=h-\floor{tm}$,
\begin{equation}\label{xtw:airy:eq:tailproduct}
 \prod_{j=i}^hT_j\longrightarrow
 \exp\!\left(-\frac{2\sqrt2}{3}t^{3/2}\right).
\end{equation}
\end{lemma}
\begin{proof}
Since $T_j=z_je^{T_j}$, the product of $z_j$ is $h!/(eh)^h$ and is treated
exactly by Stirling. Apply Lemma~\ref{xtw:airy:lem:endpoint} to $f(t)=T(t/e)$.
Its derivative is integrable, as follows either by substituting $t/e=Te^{-T}$
or by its square-root endpoint expansion, so
\[
 \sum_{j=1}^hT(j/(eh))=h\int_0^1T(t/e)\,dt+\frac12+o(1).
\]
The identity
\begin{equation}\label{xtw:airy:eq:logTintegral}
 \int_0^1\log T(t/e)\,dt=-\alpha
\end{equation}
follows by substituting $t=eTe^{-T}$, or from the antiderivative
$z\log T(z)+e^{-T(z)}$ of $\log T(z)$. The factorial and this integral give
\eqref{xtw:airy:eq:fullproduct}; the $1/2$ endpoint term accounts for $\sqrt e$.
Finally
\[
 \log T((1-y)/e)=-\sqrt{2y}+\OO(y)\quad(y\downarrow0).
\]
Summing for $y=(h-j)/h\le t/m^2$ gives \eqref{xtw:airy:eq:tailproduct} by a Riemann
sum. The total $\OO(y)$ remainder is $\OO(t^2/m)=o_h(1)$.
\end{proof}

\section{Critical finite-height amplitude}\label{xtw:airy:sec:critical}
The normalized solution of \eqref{xtw:airy:eq:wode} is exactly
\begin{equation}\label{xtw:airy:eq:airyw}
 w(t)=\frac\pi\beta\left[\Ai(0)\Bi(\beta t)-\Bi(0)\Ai(\beta t)\right].
\end{equation}
The Airy Wronskian $\Ai\Bi'-\Ai'\Bi=1/\pi$ verifies both initial conditions.
The positive-real function and derivative expansions \cite{dlmf2,dlmf7} give
\begin{align}
 w(t)&\sim\sqrt\pi\,\Ai(0)\beta^{-5/4}t^{-1/4}
         e^{(\sqrt2/3)t^{3/2}},\label{xtw:airy:eq:wlarge}\\
 \frac{2w'(t)}{w(t)}&=\sqrt{2t}-\frac1{2t}+\OO(t^{-5/2}).\label{xtw:airy:eq:wderivative}
\end{align}
Equation \eqref{xtw:airy:eq:wderivative} uses the separate expansion for $\Bi'$ as well
as for $\Bi$; it is not obtained by differentiating a bare asymptotic
equivalence. This supplies the input \eqref{xtw:airy:eq:Ularge} in the bulk proof
without using a finite-tower conclusion circularly.

\begin{theorem}[Critical tower equivalent]\label{xtw:airy:thm:critical}
As $h\to\infty$,
\begin{equation}\label{xtw:airy:eq:critical}
 H_h(1/(eh))\sim K h^{1/6}e^{-\alpha h},\qquad
 K=\frac{2^{1/6}\sqrt e}{\sqrt\pi\,\Ai(0)^2}.
\end{equation}
\end{theorem}
\begin{proof}
Let $D_i=q_i^A-q_i^B$. Descending the exact $q$ recurrence and applying the
mean-value theorem at each level gives
\[
 D_1=D_i\prod_{j=1}^{i-1}\widetilde q_j,
 \qquad q_j^B\le\widetilde q_j\le q_j^A.
\]
By squeezing, Lemma~\ref{xtw:airy:lem:bulk} applies to the product of these intermediate
factors. Also $q_1^A,q_1^B\to0$: for fixed large $t$, the bulk bounds keep
$q_2$ bounded and $q_1=(eh)^{-1}e^{q_2}$. Therefore
$H_h/D_1\to1$. Lemma~\ref{xtw:airy:lem:entrance} gives
$D_i\sim2/(m^2w(t)^2)$ at $i=h-\floor{tm}$.
Using the bulk product, \eqref{xtw:airy:eq:fullproduct}, \eqref{xtw:airy:eq:tailproduct}, and
$mv_i\to\sqrt{2t}$, we find
\begin{equation}\label{xtw:airy:eq:criticalmatch}
 \limsup_{h\to\infty}\left|
 \log\frac{H_h(1/(eh))e^{\alpha h}h^{-1/6}}
 {2\sqrt{\pi e}\,t^{-1/2}e^{(2\sqrt2/3)t^{3/2}}/w(t)^2}
 \right|=\OO(t^{-3/2}).
\end{equation}
The power is $h^{1/2}m^{-1}=h^{1/6}$, coming from the equilibrium product,
sensitivity, and inverse gap. By \eqref{xtw:airy:eq:wlarge}, the denominator in
\eqref{xtw:airy:eq:criticalmatch} tends to
\[
 \frac{2\sqrt{e/\pi}\,\beta^{5/2}}{\Ai(0)^2}
 =\frac{2^{1/6}\sqrt e}{\sqrt\pi\,\Ai(0)^2}=K.
\]
First choose $t$ large so that the logarithmic matching error is small, then
let $h$ grow. Sending $t\to\infty$ proves the limit, rather than only a
subsequential relation.
\end{proof}

\section{A uniform real critical window}\label{xtw:airy:sec:window}
For $\lambda$ in a fixed compact set $J\subset(-\infty,1]$, put
\begin{equation}\label{xtw:airy:eq:windowx}
 x_h(\lambda)=\frac{1+\lambda/m^2}{eh},\qquad m=h^{1/3}.
\end{equation}
This is positive for all sufficiently large $h$, uniformly on $J$.
All uniform statements in this section refer to such a fixed compact set.

\begin{theorem}[Real critical window]\label{xtw:airy:thm:window}
Uniformly for $\lambda\in J$,
\begin{equation}\label{xtw:airy:eq:window}
 H_h(x_h(\lambda))\sim K(\lambda)h^{1/6}
             e^{-\alpha h+\alpha\lambda m},\qquad
 K(\lambda)=\frac{2^{1/6}\sqrt e}{\sqrt\pi\,\Ai(-\beta\lambda)^2}.
\end{equation}
\end{theorem}

\subsection{Parameterized entrance and positive continuation}
Define $q^A,q^B,u^A,u^B$ exactly as before but at $x_h(\lambda)$. Their
recurrence is now
\begin{equation}\label{xtw:airy:eq:urectilt}
 \begin{split}
 u_k&=1-(1+\lambda/m^2)(1-k/h)e^{-u_{k-1}},\\
 u_0^A&=1-(1+\lambda/m^2)/e,\qquad u_0^B=1.
 \end{split}
\end{equation}
Let
\begin{equation}\label{xtw:airy:eq:wlambda}
 w_\lambda''=\tfrac12(t-\lambda)w_\lambda,\qquad
 w_\lambda(0)=0,\quad w_\lambda'(0)=1,
 \qquad U_\lambda=2w_\lambda'/w_\lambda.
\end{equation}
We first justify positivity on the whole range being used. For $\lambda\le0$,
convexity gives $w_\lambda,w_\lambda'>0$ on $t>0$. For $0<\lambda\le1$,
suppose $w_\lambda'$ first vanishes at $t\le\lambda$. Prior to that point,
$0<w_\lambda(r)\le r$, since $w_\lambda''\le0$, and
\[
 w_\lambda'(t)
 =1-\tfrac12\int_0^t(\lambda-r)w_\lambda(r)\,dr
 \ge1-\lambda t^2/4\ge3/4,
\]
a contradiction. Beyond $\lambda$, convexity preserves positivity. Therefore
$U_\lambda$ is positive on $t>0$ and
\begin{equation}\label{xtw:airy:eq:Ulambdasmall}
 U_\lambda'=t-\lambda-U_\lambda^2/2,\qquad
 U_\lambda(t)=2/t-(\lambda/3)t+\OO_J(t^2).
\end{equation}

We give the perturbation and continuation details because assuming discrete
positivity in advance would be circular. Compare \eqref{xtw:airy:eq:urectilt} up to
$k\le dm$ with the autonomous pair $\phi^{k+1}(1),\phi^k(1)$.
Stop at the first nonpositive discrete state. Before that time the derivative
of the recurrence with respect to its state is at most $1+C_J/m^2$.
The initial perturbation is $\OO_J(m^{-2})$; the step perturbation is
$\OO_J(m^{-2}+k/h)$. Discrete Gronwall gives the endpoint error
\begin{equation}\label{xtw:airy:eq:tiltperturb}
 \OO_J\!\left(m^{-2}+k/m^2+k^2/h\right).
\end{equation}
Uniformly for starting values in a compact subset of $(0,1]$, the autonomous
iterate is at least $c/(k+1)$; this follows from $k\phi^k(a)\to2$, with
finitely many initial indices absorbed in $c$. For $k\le dm$, the error in
\eqref{xtw:airy:eq:tiltperturb} is at most $C_J(m^{-2}+(d+d^2)/m)$ whereas the
smallest autonomous value is at least $c'/(dm)$. Choose fixed $d>0$ small
enough. This rules out the first crossing, including the initially bounded
indices, so \eqref{xtw:airy:eq:tiltperturb} holds throughout the entrance segment.
It follows that
\begin{equation}\label{xtw:airy:eq:tiltentry}
 mu_{\floor{dm}}=2/d+\OO_J(d+d^2)+o_h(1).
\end{equation}

For bounded $V=mu$, the recurrence Taylor expansion is
\[
 V_k-V_{k-1}=m^{-1}\bigl(k/m-\lambda-V_{k-1}^2/2\bigr)
                 +\OO_{d,L,J}(m^{-2}).
\]
Fix $J$ and $L$. The functions $U_\lambda$ have a common positive lower bound
on $(0,L]$, since they diverge uniformly at zero and are positive on compact
subintervals. Choose $d$ so that the error in \eqref{xtw:airy:eq:tiltentry} relative
to \eqref{xtw:airy:eq:Ulambdasmall} is less than one quarter of that lower bound.
The exact Riccati flow from the discrete entrance value cannot reach zero:
until a hypothetical first zero, its difference from $U_\lambda$ contracts.
It also stays in a bounded set on $[d,L]$. Standard bounded-interval Euler
stability now keeps the discrete trajectory close to that flow. Positivity
and boundedness of the discrete trajectory are thus conclusions of a stopped
comparison argument. First take $h\to\infty$, then $d\downarrow0$. Uniformly
on compact positive $t$ intervals and on $J$,
\begin{equation}\label{xtw:airy:eq:tiltentrance}
 mu_{\floor{mt}}^A\longrightarrow U_\lambda(t),\qquad
 mu_{\floor{mt}}^B\longrightarrow U_\lambda(t).
\end{equation}

For the sensitivity, use the same integral divided difference
\eqref{xtw:airy:eq:divdiff}. Summing \eqref{xtw:airy:eq:tiltperturb} over $k\le dm$ changes its
logarithmic product by $\OO_J(d^2+d^3)+o_h(1)$. The additional recurrence
factor $1+\lambda/m^2$ contributes only $\OO_J(k/m^2)=o_h(1)$, and the
initial difference ratio is $1+\lambda/m^2$. The autonomous normalization is
still $2$. Hence
\[
 m^2\Delta_{\floor{dm}}=\frac2{d^2}
       \exp\{\OO_J(d^2+d^3)+o_h(1)\}.
\]
Propagate from $d$ to $t$ using \eqref{xtw:airy:eq:tiltentrance}, with the uniform
$o_h(1/m)$ logarithmic error as before, and let $d\downarrow0$ to obtain
\begin{equation}\label{xtw:airy:eq:tiltsensitivity}
 m^2\bigl(u_{\floor{mt}}^B-u_{\floor{mt}}^A\bigr)
       \longrightarrow\frac2{w_\lambda(t)^2}.
\end{equation}

\subsection{Uniform bulk matching}
Set $\varepsilon=x_h(\lambda)$, $z_j=j\varepsilon$,
$i=h-\floor{tm}$, and $T_j=T(z_j)$ only for $j\le i$.
Choose a fixed $t>1+\max(0,\sup J)$, so $z_i<1/e$ for large $h$.
The same approximation $Q_j=T_j+\varepsilon R(z_j)$ has the uniform residual
\[
 |Q_j-z_je^{Q_{j+1}}|\le C_J\varepsilon^2T_j/(1-T_j)^4.
\]
Writing $v_i=1-T_i$, the entrance estimates give
\begin{align}
 mv_i&\longrightarrow\sqrt{2(t-\lambda)},\label{xtw:airy:eq:tiltgap}\\
 m(q_i-Q_i)&\longrightarrow\sqrt{2(t-\lambda)}-U_\lambda(t)
                         -\frac1{2(t-\lambda)}.\label{xtw:airy:eq:tiltboundary}
\end{align}
The normalized Airy representation is
\begin{equation}\label{xtw:airy:eq:shiftairy}
 w_\lambda(t)=\frac\pi\beta\bigl[
 \Ai(-\beta\lambda)\Bi(\beta(t-\lambda))
 -\Bi(-\beta\lambda)\Ai(\beta(t-\lambda))\bigr].
\end{equation}
The Wronskian checks its initial conditions. Its coefficient of the growing
$\Bi$ term is positive: a negative coefficient contradicts $w_\lambda>0$
at large $t$, and zero would make a multiple of the decaying $\Ai$,
contradicting $w_\lambda,w_\lambda'>0$. In particular
$\Ai(-\beta\lambda)>0$ for $\lambda\le1$, bounded away from zero on $J$.
The positive-real function and derivative expansions give, uniformly on $J$,
\begin{align}
 w_\lambda(t)&\sim\sqrt\pi\,\Ai(-\beta\lambda)\beta^{-5/4}
 (t-\lambda)^{-1/4}e^{(\sqrt2/3)(t-\lambda)^{3/2}},\label{xtw:airy:eq:shiftairylarge}\\
 U_\lambda(t)&=\sqrt{2(t-\lambda)}-\frac1{2(t-\lambda)}
                         +\OO_J((t-\lambda)^{-5/2}).\label{xtw:airy:eq:shiftUlarge}
\end{align}
Thus \eqref{xtw:airy:eq:tiltboundary} is $\OO_J((t-\lambda)^{-5/2})$.
Every residual, neighbor-gap, scalar contraction, and residual-sum estimate
in the proof of Lemma~\ref{xtw:airy:lem:bulk} applies on $[0,z_i]$, with uniform
constants on $J$. More explicitly, $v_i\asymp m^{-1}\sqrt{t-\lambda}$;
the summed residual after division by $v_i$ is
$\OO_J((t-\lambda)^{-3/2})+o_h(1)$ and the boundary error after that division
is $\OO_J((t-\lambda)^{-3})+o_h(1)$. The exact product identity and
\eqref{xtw:airy:eq:Rintegral} are unchanged. We obtain
\begin{equation}\label{xtw:airy:eq:tiltbulk}
 v_i\prod_{j=1}^{i-1}(q_j/T_j)
 =\exp\{\OO_J(t^{-3/2})+o_h(1)\}
\end{equation}
in the same iterated-limit sense as \eqref{xtw:airy:eq:bulk}. Squeezing also proves
this for the intermediate divided-difference factors. No value of the real
branch $T$ beyond $z_i$ is used.

\subsection{A truncated equilibrium product}
For $\lambda>0$ the real turning point occurs before level $h$, so we must
not extend the full product of Lemma~\ref{xtw:airy:lem:product} past that point.
Put
\begin{equation}\label{xtw:airy:eq:Hreal}
 \mathcal H=\frac h{1+\lambda/m^2}=h-\lambda m+\OO_J(m^{-1}),
 \qquad r=i-1,\qquad b=r/\mathcal H.
\end{equation}
Then $z_j=j/(e\mathcal H)$ and
$1-b=(t-\lambda)/m^2+\OO_J(m^{-3})$.
Factoring $T(z)=ze^{T(z)}$ and applying Stirling gives
\begin{equation}\label{xtw:airy:eq:truncatedlog}
 \log\prod_{j=1}^{r}T(j/(e\mathcal H))
 =\mathcal H\int_0^b\log T(v/e)\,dv
     +\frac12\log(2\pi r)+\frac12T(b/e)+o_h(1).
\end{equation}
Here the singular $\log(v/e)$ part is handled by the factorial, never by an
endpoint derivative estimate. For completeness, the error for the remaining
function $f(v)=T(v/e)$ is
$\int_0^b f'(v)\{\mathcal H v\}\,dv$. Extend $f'\mathbf1_{[0,b]}$ by zero
to $[0,1]$. These truncated derivatives tend to $f'$ in $L^1$, uniformly
on $J$; the bounded fractional parts have the same weak limit $1/2$ for
real $\mathcal H\to\infty$. This proves the $o_h(1)$ in
\eqref{xtw:airy:eq:truncatedlog}, including its moving endpoint.

By \eqref{xtw:airy:eq:logTintegral} and the square-root expansion at $1$,
\[
 \mathcal H\int_0^b\log T(v/e)\,dv
 =-\alpha\mathcal H+\frac{2\sqrt2}{3}(t-\lambda)^{3/2}+o_h(1).
\]
The $\OO(m^{-3})$ floor error costs $o_h(1)$ in this exponent.
As $r/h\to1$ and $T(b/e)\to1$, we have proved
\begin{equation}\label{xtw:airy:eq:truncatedproduct}
 \prod_{j=1}^{i-1}T(j\varepsilon)
 \sim\sqrt{2\pi e h}\,
 e^{-\alpha h+\alpha\lambda m+(2\sqrt2/3)(t-\lambda)^{3/2}},
\end{equation}
uniformly on $J$ for fixed sufficiently large $t$.

\subsection{Completion of the window theorem}
The exact difference product still gives
$H_h\sim D_i\prod_{j<i}\widetilde q_j$, because $q_1^A,q_1^B\to0$
uniformly on $J$. Combine \eqref{xtw:airy:eq:tiltsensitivity}, \eqref{xtw:airy:eq:tiltgap},
\eqref{xtw:airy:eq:tiltbulk}, and \eqref{xtw:airy:eq:truncatedproduct}. The ratio of
$H_h e^{\alpha h-\alpha\lambda m}h^{-1/6}$ to
\[
 \frac{2\sqrt{\pi e}\,(t-\lambda)^{-1/2}
 e^{(2\sqrt2/3)(t-\lambda)^{3/2}}}{w_\lambda(t)^2}
\]
has limiting logarithmic error $\OO_J(t^{-3/2})$. By
\eqref{xtw:airy:eq:shiftairylarge} this last display tends uniformly on $J$ to
$K(\lambda)$. Take $h\to\infty$ first and then $t\to\infty$.
This proves Theorem~\ref{xtw:airy:thm:window} with the asserted compact uniformity.

\section{The Airy size law and its density}\label{xtw:airy:sec:airyprob}
Define $N_h$ by \eqref{xtw:airy:eq:sizeprobintro}. Since $H_h$ is entire, positive at
$x_h=1/(eh)$, and has nonnegative coefficients, this is a probability law
with support in the integers $n\ge h$. Set
\[
 X_h=(N_h-\alpha h)/m^2.
\]
For real $\lambda$ near zero,
\begin{equation}\label{xtw:airy:eq:MGF}
 M_h(\lambda)=\E e^{\lambda X_h}
 =e^{-\alpha\lambda m}
 \frac{H_h(x_he^{\lambda/m^2})}{H_h(x_h)}.
\end{equation}
The effective parameter in Theorem~\ref{xtw:airy:thm:window} is
\[
 \lambda_{\rm eff}=m^2(e^{\lambda/m^2}-1)
 =\lambda+\OO_J(m^{-2}).
\]
Even in the exponential factor the discrepancy vanishes:
$\alpha m(\lambda_{\rm eff}-\lambda)=\OO_J(m^{-1})$.
Uniformity of the window theorem therefore gives
\begin{equation}\label{xtw:airy:eq:MGFlimit}
 M_h(\lambda)\longrightarrow
 M(\lambda)=\frac{\Ai(0)^2}{\Ai(-\beta\lambda)^2}
\end{equation}
on an open real interval containing zero. The moment-generating-function
continuity theorem implies weak convergence $X_h\xrightarrow d X$.

\subsection{Analytic continuation to every imaginary argument}
One still needs to identify the characteristic function for every real
Fourier argument, not just near zero. Choose real $a<0<b<1$ inside the
interval used in \eqref{xtw:airy:eq:MGFlimit}. The complex MGFs are entire and, for
$a\le\Re z\le b$,
\begin{equation}\label{xtw:airy:eq:vitalibound}
 |M_h(z)|\le M_h(\Re z)\le M_h(a)+M_h(b).
\end{equation}
The right-hand side is uniformly bounded. Vitali's theorem consequently gives
local analytic convergence throughout the vertical strip to the analytic
continuation of \eqref{xtw:airy:eq:MGFlimit}. To spell out the uniqueness step,
normal-family compactness supplies a convergent subsequence on each compact
subset; every such subsequential limit agrees on the real interval and hence,
by the identity theorem, on the entire connected strip. Thus all subsequences
have the same limit. The Airy function has only negative real zeros
\cite{dlmf9}, and the positivity proved after \eqref{xtw:airy:eq:shiftairy} excludes a
zero of $\Ai(-\beta z)$ on the chosen real interval. The displayed candidate
is therefore holomorphic on the strip. In particular
\begin{equation}\label{xtw:airy:eq:psi}
 \E e^{\mathrm itX}=\psi(t)=\frac{\Ai(0)^2}{\Ai(-\beta\mathrm it)^2}
 \quad(t\in\R).
\end{equation}

The complex Airy asymptotic on the rays with arguments $\pm\pi/2$
\cite{dlmf7} gives
\begin{equation}\label{xtw:airy:eq:psitail}
 |\Ai(-\beta\mathrm it)^{-2}|
 =\OO\!\left(|t|^{1/2}e^{-(2/3)|t|^{3/2}}\right)
 \quad(|t|\to\infty).
\end{equation}
Thus $\psi\in L^1(\R)$, and Fourier inversion gives the bounded continuous
density
\begin{equation}\label{xtw:airy:eq:pdensity}
 p(v)=\frac1{2\pi}\int_{\R}e^{-\mathrm itv}\psi(t)\,dt.
\end{equation}

\subsection{The exact density at zero}
The Wronskian and the connection identity \cite{dlmf2} imply
\begin{align}
 \left(\frac{\Bi}{\Ai}\right)'(z)&=\frac1{\pi\Ai(z)^2},\label{xtw:airy:eq:ratioDerivative}\\
 \Bi(z)&=\mathrm i\Ai(z)+2e^{-\mathrm i\pi/6}
                      \Ai(ze^{-2\pi\mathrm i/3}).\label{xtw:airy:eq:connection}
\end{align}
At $z=\mathrm iR$, the rotated Airy argument in \eqref{xtw:airy:eq:connection} has
angle $-\pi/6$. Its Airy function decays exponentially, while
$\Ai(\mathrm iR)$ grows exponentially. Hence
\[
 \frac{\Bi(\mathrm iR)}{\Ai(\mathrm iR)}\longrightarrow\mathrm i.
\]
Complex conjugation gives the limit $-\mathrm i$ at $z=-\mathrm iR$.
There are no Airy zeros on the imaginary axis. Integrating
\eqref{xtw:airy:eq:ratioDerivative} upward along it gives
\begin{equation}\label{xtw:airy:eq:verticalintegral}
 \int_{-\mathrm i\infty}^{\mathrm i\infty}\frac{dz}{\Ai(z)^2}
 =\pi\left[\frac{\Bi}{\Ai}\right]_{-\mathrm i\infty}^{\mathrm i\infty}
 =2\pi\mathrm i.
\end{equation}
The integral is absolutely convergent by \eqref{xtw:airy:eq:psitail}. Under
$z=-\beta\mathrm it$, the increasing real $t$ axis maps to the downward
imaginary axis. Since $dt=\mathrm i\,dz/\beta$, reversing the orientation
in \eqref{xtw:airy:eq:verticalintegral} proves
\begin{equation}\label{xtw:airy:eq:pzero}
 p(0)=\frac{\Ai(0)^2}{2\pi}\int_{\R}\frac{dt}{\Ai(-\beta\mathrm it)^2}
     =\frac{\Ai(0)^2}{\beta}>0.
\end{equation}
This calculation fixes the normalization exactly. It does not yet imply
anything about the individual lattice probabilities $\Prb(N_h=n)$.

\section{Fourier minor arcs from positive factorization}\label{xtw:airy:sec:minor}
Let
\begin{equation}\label{xtw:airy:eq:phih}
 \varphi_h(\theta)=\E e^{\mathrm i\theta N_h}
   =\frac{H_h(x_he^{\mathrm i\theta})}{H_h(x_h)},\qquad x_h=1/(eh).
\end{equation}
The argument in this section does not use the critical equivalent or the
Airy window theorem.

\begin{theorem}[Uniform minor-arc bound]\label{xtw:airy:thm:minor}
There exist constants $c>0$, $A>0$, and $h_0$ such that, for $h\ge h_0$,
\begin{equation}\label{xtw:airy:eq:minor}
 |\varphi_h(\theta)|\le
 \begin{cases}
 e^{-ch|\theta|^{3/2}},&Ah^{-2/3}\le|\theta|\le1/8,\\
 e^{-ch},&1/8\le|\theta|\le\pi.
 \end{cases}
\end{equation}
\end{theorem}

\subsection{An exact positive factorization}
For a complex variable $z$, abbreviate
$A_i(z)=A_{i,h}(z)$, $B_i(z)=A_{i,h-1}(z)$,
$\Delta_i(z)=A_i(z)-B_i(z)$, where $B_h=1$. Define
\[
 D_i(z)=\int_0^1\exp\bigl(izt\Delta_{i+1}(z)\bigr)\,dt
 \quad(1\le i<h).
\]
Here $i$ is the integer level, not the imaginary unit $\mathrm i$.
The exact difference identity is
\[
 \Delta_i(z)=izB_i(z)\Delta_{i+1}(z)D_i(z),\qquad
 \Delta_h(z)=e^{hz}-1.
\]
Iteration gives an identity of entire functions,
\begin{equation}\label{xtw:airy:eq:factorization}
 H_h(z)=(e^{hz}-1)(h-1)!z^{h-1}
                 \prod_{i=1}^{h-1}B_i(z)D_i(z).
\end{equation}
Each factor has nonnegative Taylor coefficients. For $D_i$ this follows from
the nonnegative coefficients and zero constant term of $\Delta_{i+1}$,
followed by coefficientwise integration over $0\le t\le1$.
Normalize each factor by its positive value at $x_h$ and apply the triangle
inequality separately. The monomial has normalized modulus one, and all
other normalized moduli are at most one. Dropping the terminal and $D_i$
factors gives
\begin{equation}\label{xtw:airy:eq:dropfactors}
 |\varphi_h(\theta)|\le
 \prod_{i=1}^{h-1}\frac{|B_i(x_he^{\mathrm i\theta})|}{B_i(x_h)}.
\end{equation}
This remains true if a dropped factor vanishes on the circle.

\subsection{Homogeneous tree comparison}
Write $Q_i(z)=izB_{i+1}(z)$, the globally defined entire exponent of $B_i$;
no choice of logarithm is involved. All its coefficients are nonnegative.
Define
\begin{equation}\label{xtw:airy:eq:homogeneous}
 T_0(z)=0,\qquad T_{j+1}(z)=ze^{T_j(z)},\qquad
 T(z)=\sum_{r\ge1}\frac{r^{r-1}}{r!}z^r.
\end{equation}
The final series follows from Lagrange inversion and is the same small
solution $Te^{-T}=z$ used earlier. Coefficientwise,
\begin{equation}\label{xtw:airy:eq:homogeneouscompare}
 Q_i(z)\succeq T_{h-i}(iz),\qquad T_L(z)\preceq T(z).
\end{equation}
The first comparison starts at $i=h-1$ with equality
$Q_{h-1}(z)=(h-1)z=T_1((h-1)z)$. Descending induction uses that every deeper
level parameter is at least $i$ and exponentiation preserves nonnegative
coefficient inequalities.
Put $z_i=ix_h=i/(eh)$ and
\begin{equation}\label{xtw:airy:eq:damping}
 \mathcal D_L(z,\theta)
 =\sum_{r\ge1}[z^r]T_L(z)\,z^r(1-\cos(r\theta)).
\end{equation}
In \eqref{xtw:airy:eq:damping} the coefficient notation refers to the power series
before evaluating its argument; equivalently, if $T_L(z)=\sum t_{L,r}z^r$,
then $\mathcal D_L=\sum t_{L,r}z^r(1-\cos(r\theta))$.
Since $1-\cos(r\theta)\ge0$, applying \eqref{xtw:airy:eq:homogeneouscompare} to the
real-part loss of each exponent in \eqref{xtw:airy:eq:dropfactors} is legitimate and yields
\begin{equation}\label{xtw:airy:eq:sumdamping}
 |\varphi_h(\theta)|\le
 \exp\!\left(-\sum_{i=1}^{h-1}\mathcal D_{h-i}(z_i,\theta)\right).
\end{equation}
The comparison is on nonnegative coefficient sums, not on oscillatory
complex values.

\subsection{A finite-iteration tail bound}
For $0\le z<1/e$, set $\tau=T(z)$, $\epsilon=1-\tau$, and
$d_j=\tau-T_j(z)$. Then
\[
 d_{j+1}=\tau(1-e^{-d_j}),\qquad0\le d_0=\tau\le1.
\]
For $0\le u\le1$, Taylor's inequality gives
$1-e^{-u}\le u-u^2/3$. Consequently
\[
 d_{j+1}\le d_j-d_j^2/3,\qquad d_{j+1}\le\tau d_j.
\]
The first inequality gives $d_j\le3/(j+3)$: its reciprocal increases by at
least $1/3$, starting from $1/d_0\ge1$ (the zero case is immediate).
Apply the second inequality after $j=\floor{L/2}$ to obtain
\begin{equation}\label{xtw:airy:eq:truncationtail}
 d_L\le\frac6L e^{-\epsilon L/2},\qquad
 d_L\le6\epsilon e^{-\epsilon L/2}\quad\hbox{if }\epsilon L\ge1.
\end{equation}
Positivity of $T-T_L$ and $0\le1-\cos\le2$ imply
\begin{equation}\label{xtw:airy:eq:truncationdamp}
 \mathcal D_L(z,\theta)\ge\mathcal D_\infty(z,\theta)-2d_L,
 \quad
 \mathcal D_\infty(z,\theta)=
 \sum_{r\ge1}\frac{r^{r-1}z^r}{r!}(1-\cos(r\theta)).
\end{equation}
This uses only real subcritical contraction, with no assumed complex
fixed-point estimate.

\subsection{A block of near-terminal levels}
By symmetry suppose $0<\theta\le1/8$. Select integer levels with
\[
 k=h-i\in[h\theta,2h\theta],\quad L=k,\quad
 \delta=k/h,\quad z_i=(1-\delta)/e.
\]
Thus $\theta\le\delta\le2\theta\le1/4$.
For integers $r\in[1/\theta,2/\theta]$ we have $r\theta\in[1,2]$, so
$1-\cos(r\theta)\ge1-\cos1$.
The elementary Stirling bound $r!\le3\sqrt r(r/e)^r$, valid for $r\ge1$,
gives
\[
 \frac{r^{r-1}z_i^r}{r!}\ge\frac{(1-\delta)^r}{3r^{3/2}}
 \ge\frac{e^{-8}}{3r^{3/2}}.
\]
Indeed $\log(1-\delta)\ge-2\delta$ and $r\delta\le4$.
The degree interval contains at least $7/(8\theta)$ integers, each with
$r^{-3/2}\ge(\theta/2)^{3/2}$. Hence
\begin{equation}\label{xtw:airy:eq:dinflower}
 \mathcal D_\infty(z_i,\theta)\ge c_0\sqrt\theta,
 \qquad c_0=\frac{7(1-\cos1)e^{-8}}{24\,2^{3/2}}>0.
\end{equation}
Write $\epsilon_i=1-T(z_i)$. The exact gap identity is
\[
 \delta=1-(1-\epsilon_i)e^{\epsilon_i}
        =\int_0^{\epsilon_i}v e^v\,dv.
\]
As $0\le\epsilon_i\le1$, it follows that
\begin{equation}\label{xtw:airy:eq:gapbounds}
 \sqrt{2/e}\sqrt\theta\le\epsilon_i\le2\sqrt\theta,
 \qquad\epsilon_iL\ge\sqrt{2/e}\,h\theta^{3/2}.
\end{equation}
Put $X=h\theta^{3/2}$. For a fixed sufficiently large $M$ and $X\ge M$,
\eqref{xtw:airy:eq:truncationtail} and \eqref{xtw:airy:eq:gapbounds} yield
\[
 2d_L\le24\sqrt\theta\,e^{-\sqrt{2/e}X/2}
       \le(c_0/2)\sqrt\theta.
\]
Here choose $M$ so that $\sqrt{2/e}M\ge1$ and
$24e^{-\sqrt{2/e}M/2}\le c_0/2$.
Equations \eqref{xtw:airy:eq:truncationdamp} and \eqref{xtw:airy:eq:dinflower} then give
$\mathcal D_L(z_i,\theta)\ge(c_0/2)\sqrt\theta$ for every selected level.
Their number is at least $h\theta-1\ge h\theta/2$ after increasing $M$.
All indices lie in $1\le k\le h-1$ for large $h$, because $k\le h/4$ and
the lower cutoff makes $h\theta$ large. Substitution in
\eqref{xtw:airy:eq:sumdamping} gives
\begin{equation}\label{xtw:airy:eq:smallarcsfinal}
 |\varphi_h(\theta)|\le e^{-(c_0/4)h\theta^{3/2}}
 \quad\hbox{if }0<\theta\le1/8,\quad
 \theta\ge M^{2/3}h^{-2/3}.
\end{equation}

\subsection{Large arcs and lattice resonances}
For $1/8\le|\theta|\le\pi$, keep only the degree-one coefficient
$Q_i(z)\succeq iz$. At every level $i\in[h/4,h/2]$ its real-part loss is
at least
\[
 ix_h(1-\cos\theta)\ge(4e)^{-1}(1-\cos(1/8))>0.
\]
There are a positive constant times $h$ such levels. Equation
\eqref{xtw:airy:eq:dropfactors} gives $|\varphi_h(\theta)|\le e^{-ch}$ uniformly.
Together with \eqref{xtw:airy:eq:smallarcsfinal}, this proves
Theorem~\ref{xtw:airy:thm:minor}. In particular no lattice resonance away from zero
is left untreated.

\section{Uniform lattice transfer}\label{xtw:airy:sec:local}
We now combine the pointwise characteristic limit of Section~\ref{xtw:airy:sec:airyprob}
with the separately proved Fourier bounds. For integration define
\[
 \Phi_h(t)=e^{-\mathrm it\alpha h/h^{2/3}}
                \varphi_h(t/h^{2/3})
\]
on $|t|\le\pi h^{2/3}$ and extend it by zero outside that interval.
For each fixed $t$, $\Phi_h(t)\to\psi(t)$. On every fixed compact interval,
dominated convergence applies because $|\Phi_h|\le1$. On
$A\le|t|\le h^{2/3}/8$, Theorem~\ref{xtw:airy:thm:minor} gives
$|\Phi_h(t)|\le e^{-c|t|^{3/2}}$. On the remaining large arc the integral is
at most $2\pi h^{2/3}e^{-ch}\to0$. Since $\psi$ is integrable,
\begin{equation}\label{xtw:airy:eq:Lone}
 \int_\R|\Phi_h(t)-\psi(t)|\,dt\longrightarrow0.
\end{equation}
The lattice inversion formula is
\[
 h^{2/3}\Prb(N_h=n)=\frac1{2\pi}
 \int_{-\pi h^{2/3}}^{\pi h^{2/3}}
 e^{-\mathrm it(n-\alpha h)/h^{2/3}}\Phi_h(t)\,dt.
\]
Compare this with \eqref{xtw:airy:eq:pdensity}; the absolute error, uniformly in every
integer $n$, is at most $(2\pi)^{-1}$ times \eqref{xtw:airy:eq:Lone}.
We have proved the uniform local limit \eqref{xtw:airy:eq:lltintro}. In particular,
for some absolute $C$ and every integer $n$ and $h\ge1$,
\begin{equation}\label{xtw:airy:eq:atombound}
 \Prb(N_h=n)\le C h^{-2/3}.
\end{equation}
For sufficiently large $h$ this follows from the uniform local limit and
boundedness of $p$; enlarge $C$ over finitely many smaller heights. It can
also be obtained directly by integrating Theorem~\ref{xtw:airy:thm:minor}, bounding
the central arc by its width.

\section{All-height domination and the coefficient sum}\label{xtw:airy:sec:sum}
The exact probability definition gives
\begin{equation}\label{xtw:airy:eq:exactcoefficient}
 \frac{c_{n,h}}{n!}=\frac{(eh)^n}{n!}
                 H_h(1/(eh))\Prb(N_h=n).
\end{equation}
Use the critical equivalent and \eqref{xtw:airy:eq:atombound}, increasing a constant
for the finitely many small $h$, to get the global bound
\begin{equation}\label{xtw:airy:eq:allheight}
 \frac{c_{n,h}}{n!}\le
 C\frac{e^n}{n!}h^{n-1/2}e^{-\alpha h}
 \quad(h\ge1,\ n\ge0).
\end{equation}
The left side vanishes when $h>n$. This bound is available at every height;
it does not extend a compact-window asymptotic into the tails.

Let
\begin{equation}\label{xtw:airy:eq:heightwindow}
 \mathcal W_n=\{h\in\Z:|h-n/\alpha|\le n^{3/5}\}.
\end{equation}
On this window $h\sim n/\alpha$ and
$(n-\alpha h)/h^{2/3}\to0$ uniformly. Continuity and positivity of $p(0)$,
the uniform local limit, and the critical equivalent yield
\begin{equation}\label{xtw:airy:eq:centralcoeff}
 \frac{c_{n,h}}{n!}
 =(Kp(0)+o(1))\frac{e^n}{n!}h^{n-1/2}e^{-\alpha h}
 \quad\hbox{uniformly for }h\in\mathcal W_n.
\end{equation}
The window is wider than $\sqrt n$ and narrower than $n^{2/3}$; these two
inequalities are the reason for the convenient exponent $3/5$.
As $\alpha>1$, all its heights lie in $1\le h\le n$ for large $n$.

\begin{lemma}[Discrete Laplace sum with tails]\label{xtw:airy:lem:laplace}
For $f_n(t)=t^{n-1/2}e^{-\alpha t}$ on $t>0$,
\begin{equation}\label{xtw:airy:eq:laplace}
 \sum_{h\ge1}f_n(h)\sim\int_0^\infty f_n(t)\,dt
 =\frac{\Gamma(n+1/2)}{\alpha^{n+1/2}}.
\end{equation}
Restricting the sum to $\mathcal W_n$ preserves this equivalent.
\end{lemma}
\begin{proof}
A nonnegative unimodal function has sum-integral error at most twice its
maximum, by applying the left and right rectangle comparisons on its two
monotone pieces. Here the maximum is attained at
$t=(n-1/2)/\alpha$. Stirling gives
\[
 \frac{\int_0^\infty f_n(t)\,dt}{\max_{t>0}f_n(t)}
 \sim\frac{\sqrt{2\pi n}}\alpha.
\]
Thus the relative error is $\OO(n^{-1/2})$.
For the tails, write $t=(n/\alpha)(1+u)$. Relative to the central value,
\begin{equation}\label{xtw:airy:eq:laplaceexponent}
 \log\frac{f_n(t)}{f_n(n/\alpha)}
 =n\bigl(\log(1+u)-u\bigr)-\tfrac12\log(1+u).
\end{equation}
On $|u|\le1/2$, the first term is at most $-cnu^2$.
On either side outside this fixed central region the exponent has a
strictly negative gap, and toward $u=\infty$ it is eventually bounded above
by $-cnu$ after adjusting constants; toward $u=-1$ it tends to $-\infty$.
The boundaries of $\mathcal W_n$ have $|u|=\alpha n^{-2/5}$, so their
exponential loss is at least $c n^{1/5}$. Integrating these elementary
bounds shows that the omitted integrals are at most a polynomial in $n$
times $e^{-cn^{1/5}}$ times the full integral. Apply the rectangle
comparison separately to each monotone tail: its sum differs from its
integral by at most twice its boundary maximum, which satisfies the same
bound. This proves that the omitted discrete mass is negligible.
\end{proof}

We can now prove Theorem~\ref{xtw:airy:thm:main}. Sum \eqref{xtw:airy:eq:centralcoeff} on
$\mathcal W_n$ and use \eqref{xtw:airy:eq:allheight} with Lemma~\ref{xtw:airy:lem:laplace}
outside it. The latter contribution is negligible, so
\begin{align}
 b_n&\sim Kp(0)\frac{e^n}{n!}
              \frac{\Gamma(n+1/2)}{\alpha^{n+1/2}}\notag\\
 &\sim\frac{Kp(0)}{\sqrt\alpha}\,\rho^{-n}n^{-1/2}.
 \label{xtw:airy:eq:assembledconstant}
\end{align}
By \eqref{xtw:airy:eq:constants} and \eqref{xtw:airy:eq:pzero},
\begin{equation}\label{xtw:airy:eq:cancelconstant}
 \frac{Kp(0)}{\sqrt\alpha}
 =\frac{2^{1/6}\sqrt e}{\sqrt\pi\,\beta\sqrt\alpha}
 =\sqrt{\frac{2}{\pi\rho}}.
\end{equation}
The two factors $\Ai(0)^2$ cancel only after the lattice extraction and
height summation have both been justified. Multiplying by
$(n!)^2\sim2\pi n(n/e)^{2n}$ proves \eqref{xtw:airy:eq:anmain}, and taking logarithms
proves \eqref{xtw:airy:eq:logmain}. This completes the coefficient theorem.

\section{Gaussian height fluctuations}\label{xtw:airy:sec:height}
For the weighted-tree law of Theorem~\ref{xtw:airy:thm:limitsintro},
\begin{equation}\label{xtw:airy:eq:heightprob}
 \Prb(L_n=h)=\frac{n!c_{n,h}}{a_n}
            =\frac{c_{n,h}/n!}{b_n}.
\end{equation}
Equations \eqref{xtw:airy:eq:centralcoeff}, \eqref{xtw:airy:eq:assembledconstant}, and
Lemma~\ref{xtw:airy:lem:laplace} show that throughout $\mathcal W_n$ this law is
uniformly asymptotic to the normalized positive kernel
\[
 \frac{h^{n-1/2}e^{-\alpha h}}{\sum_{j\ge1}j^{n-1/2}e^{-\alpha j}}.
\]
Equation \eqref{xtw:airy:eq:allheight} and the tail part of Lemma~\ref{xtw:airy:lem:laplace}
show that both distributions give asymptotically vanishing mass to the
complement of $\mathcal W_n$. It follows in particular that their total
variation distance tends to zero.

For fixed bounded $s$, set $h=n/\alpha+s\sqrt n/\alpha$, with grid spacing
$\alpha/\sqrt n$ in the variable $s$. Taylor expansion of
\eqref{xtw:airy:eq:laplaceexponent} at $u=s/\sqrt n$ gives, uniformly on bounded
$s$ intervals,
\[
 \log\frac{f_n(h)}{f_n(n/\alpha)}=-s^2/2+o(1).
\]
The same quadratic bound used in Lemma~\ref{xtw:airy:lem:laplace} gives uniformly
small normalized tails as $|s|$ becomes large, with the farther tails
exponentially small. Riemann sums and the normalizing integral therefore
imply convergence of this kernel to the standard normal density in the
$s$ coordinate. Total variation comparison transfers the weak limit to
$L_n$, proving \eqref{xtw:airy:eq:cltintro}.

This statement centers at $n/\alpha$ and scales by $\sqrt n/\alpha$.
We do not identify these with exact finite-$n$ mean and standard deviation,
or infer moment expansions merely from weak convergence. Nor does this
argument say that the Airy law itself becomes Gaussian: that law describes
size under a fixed-height Boltzmann measure, whereas this one describes
height under the fixed-size tree measure.

\section{Refined inversion and the integer threshold}\label{xtw:airy:sec:inverse}
Put
\begin{equation}\label{xtw:airy:eq:inversef}
 A=2\sqrt{2\pi/\rho},\qquad
 f(t)=2t\log\!\left(\frac{t}{e\sqrt\rho}\right).
\end{equation}
Because $2+\log\rho=1+\log(e-1)$, Theorem~\ref{xtw:airy:thm:main} gives
\begin{equation}\label{xtw:airy:eq:loginverse}
 \log a_n=f(n)+\tfrac12\log n+\log A+o(1).
\end{equation}
Let $g$ be the continuous piecewise-linear interpolation through these
integer log values on the eventually increasing portion. The second
derivative of $f(t)+\tfrac12\log t$ is $\OO(1/t)$. Linear interpolation
therefore changes the smooth expression by $\OO(1/t)$, and the integer
$o(1)$ error remains uniformly $o(1)$ after interpolation. Thus
\begin{equation}\label{xtw:airy:eq:g}
 g(t)=f(t)+\tfrac12\log t+\log A+o(1)
 \quad(t\to\infty,\ t\in\R).
\end{equation}

The positive real solution of $f(x)=L=\log Y$ is
\[
 x=\frac{L}{2W_0(L/(2e\sqrt\rho))}.
\]
To verify this, set $w=\log(x/(e\sqrt\rho))$; then
$we^w=L/(2e\sqrt\rho)$ and $x=L/(2w)$.
Also
\begin{equation}\label{xtw:airy:eq:fprime}
 f'(x)=2\log(x/\sqrt\rho),\qquad f''(x)=2/x.
\end{equation}
Let $z(Y)=g^{-1}(L)$. First $z\sim x$ follows from \eqref{xtw:airy:eq:g} and the
leading term $2t\log t$. On the interval between $x$ and $z$, the derivative
in \eqref{xtw:airy:eq:fprime} is comparable to $\log x$. The correction in
\eqref{xtw:airy:eq:g} is $\OO(\log z)$, so the mean-value theorem improves this to
$z-x=\OO(1)$. Taylor expansion with that bounded displacement now yields
\begin{align}
 z&=x-\frac{\tfrac12\log x+\log A}{2\log(x/\sqrt\rho)}
          +o(1/\log x)\notag\\
  &=x-\frac14-
    \frac{\log(8\pi/\sqrt\rho)}{4\log(x/\sqrt\rho)}
          +o(1/\log x).\label{xtw:airy:eq:inversefinal}
\end{align}
For the algebra in the second line use $A^2=8\pi/\rho$.

By Lemma~\ref{xtw:airy:lem:monotonicity}, for all sufficiently large thresholds the
integer inverse is exactly
\begin{equation}\label{xtw:airy:eq:ceilingexact}
 N(Y)=\ceil{z(Y)}.
\end{equation}
Equation \eqref{xtw:airy:eq:inversefinal} implies that some positive
$\epsilon(Y)\to0$ bounds $|z(Y)-r(Y)|\log x$ from above. Monotonicity of the
ceiling proves \eqref{xtw:airy:eq:sandwichintro}, completing
Theorem~\ref{xtw:airy:thm:inverseintro}. The envelope is not explicitly estimated by
the proof. Thresholds may approach an integer crossing arbitrarily closely,
so replacing this two-sided sandwich by one ceiling of $r(Y)$ would not be
justified.
\begin{writenote}[7 October 2026]
In the transseries volume's terms (front matter, ``The inverses and the transseries volume''), \eqref{xtw:airy:eq:ceilingexact} is an instance of item~(1) of \file{p0:thm:staircase}, with the piecewise-linear $g$ as admissible interpolation, and the sandwich follows from it as item~(2) does; the Lambert base $x$ is an instance of \file{p0:prop:factorial-core}. Part~II carries $r(Y)$ one term further (\eqref{xtw:frac:eq:y1inverse}).
\end{writenote}

\section{Exact verification and reproducibility}\label{xtw:airy:sec:verification}
The accompanying archive contains this TeX source and PDF, authored Python
code, small attributed integer fixtures, generated verification receipts,
and a deterministic offline builder. The core arithmetic and package checks
use only the Python standard library; optional arbitrary-precision Airy
diagnostics are isolated from the exact verification.

\subsection{Three finite descriptions of the coefficients}
The primary coefficient computation uses \eqref{xtw:airy:eq:coefficientrec} with
exact integers. It recomputes all coefficients through $n=48$ and the
exact-height decomposition, rather than importing a floating approximation.
A triangular truncation of the same defining recurrence is a useful second
implementation but is not represented as an independent combinatorial model.

The first independent combinatorial check enumerates all Pr\"ufer strings on
$n+1$ labels, roots each resulting labeled tree at $0$, computes all depths,
and accumulates their product by maximum depth. For $n=1,\ldots,7$ the total
number of trees checked is
\[
 \sum_{n=1}^7(n+1)^{n-1}=280392.
\]
It verifies every one of the $28$ positive-height cells in that range against
\eqref{xtw:airy:eq:heightidentity}. The one-vertex root object is checked separately.

A further exact method sums by level populations. If the positive level
sizes are $n_1,\ldots,n_h$ with sum $n$, then the labels can be assigned in
$n!/\prod_jn_j!$ ways, and each vertex at level $j\ge2$ chooses one of the
$n_{j-1}$ parents at the preceding level. For shift $y$, the total contribution
of this profile is
\begin{equation}\label{xtw:airy:eq:profile}
 \frac{n!}{\prod_{j=1}^h n_j!}
 \prod_{j=2}^h n_{j-1}^{n_j}
 \prod_{j=1}^h(y+j-1)^{n_j}.
\end{equation}
Every such assignment is a rooted tree and has exactly the prescribed levels,
so summing \eqref{xtw:airy:eq:profile} over positive compositions gives another
independent finite formula. The package checks all height cells through
$n=12$ using this formula. These checks confirm the weighting and
normalization, not the limiting constants by numerical extrapolation.

\subsection{Provenance and deliberately bounded claims}
The official displayed A096537 entry supplies $15$ terms, indices $0$ through
$14$, which are checked exactly \cite{oeis537}. Its linked table advertises
indices through $200$, but that file was not retrieved or used in this
package. The coefficients from $15$ through $48$ in the regression fixture
are explicitly labeled generated data; they are not additional official
OEIS comparisons. Nine complete polynomial rows of A096542, through $n=8$,
are also checked at the rational shifts $1/2,1,2,3$ \cite{oeis542}.
The defining series and these rows imply
$T(n,1)=n a_{n-1}$, in agreement with the cross-reference on A096537; the
retrieved A096542 display $T(n,1)=na_n$ has an index mismatch. That display
is not used as an input identity.

\begin{remark}[\textsf{[write]} The A096542 formula for $T(n,1)$ is off by one index]\label{xtw:airy:rem:oeis542}
[Added 7 October 2026, batch~110.] OEIS A096542 (revision \#11,
19~December 2017, live on 7~October 2026), whose rows are $T(n,k)=[y^k]P_n(y)$,
has the formula line ``T(n, 1) = n*A096537(n)''. The correct identity is
\[
 T(n,1)=n\,a_{n-1}\qquad(n\ge1).
\]
Proof: differentiating $F_y=\exp(xyF_{y+1})$ in $y$ at $y=0$ gives
$[y]F_y=xF_1$, so $T(n,1)/n!=[x^n]\,xF_1(x)=a_{n-1}/(n-1)!$ (Part~II proves the
same in Section~\ref{xtw:frac:sec:reproducibility}). The entry's line is false
for every $n\ge2$, because $a_n\ge n^2a_{n-1}>a_{n-1}$ there
(Lemma~\ref{xtw:airy:lem:monotonicity}). The entry's own rows show it: row~2
is $0,2,3$, so $T(2,1)=2=2a_1$, while $2a_2=10$; likewise $T(3,1)=15=3\cdot5$,
$T(4,1)=244=4\cdot61$ and $T(5,1)=6885=5\cdot1377$, against $3\cdot61$,
$4\cdot1377$ and $5\cdot49721$. A096537's own line ``a(n) = A096542(n+1,
1)/(n+1)'' is the correct identity. No correction has been submitted to the
OEIS.
\end{remark}

These finite checks include coefficient nonnegativity, the all-path diagonal
$n!c_{n,n}=(n!)^2$, and domain/fixture rejection tests. Supporting shifted or
height-marked finite checks do not assert a shifted-family prefactor theorem
or an all-orders asymptotic extension of this report.

\subsection{Running and rebuilding}
From the extracted package, the mandatory replay is
\begin{quote}\ttfamily
 python3 -I -S -B reproduce.py
\end{quote}
It runs the exact checker both normally and under \texttt{-O} and requires
byte-identical canonical JSON. Mathematical validity checks use explicit
exceptions rather than removable Python assertions. The guard suite is
\texttt{test\_build.py}; the extracted release inventory and its hashes are
checked by \texttt{verify\_manifest.py}.

For the complete PDF and ZIP build, use \texttt{build.py --output} with a
fresh external directory whose parent already exists. The builder rejects
existing output paths, source/output overlap, traversal components, symlinks,
wrong types, unexpected inventory, changed sources, and malformed receipts.
It runs the guard suite and exact replay in both Python modes, then invokes
installed pdfLaTeX in an isolated environment for three passes with shell
escape disabled. The settled final log must have no warnings, unresolved
references, missing glyphs, or overfull/underfull boxes, and the final two
auxiliary states must agree.

The fixed epoch is 4 October 2026 at 00:00 UTC. Archive members have sorted
names, fixed timestamps and permissions, and stored compression. A closed
manifest records every distributed non-manifest file with its byte count
and SHA-256. The separately delivered TeX and PDF are byte-identical to the
archive copies. With the same Python/TeX stack, clean extractions rebuild
byte-identical artifacts; different TeX or font versions may change PDF
bytes. Hashes detect accidental modification, not replacement of both the
software and its hashes by an adversary. Full command instructions and the
precise manifest schema are in \texttt{README\_REPRODUCIBILITY.md}.
\begin{writenote}[7 October 2026]
In this report the package's programs are \file{code/192-airy-*.py}, its fixtures, generated coefficients and recorded outputs \file{data/192-airy-*}, and its \file{README_REPRODUCIBILITY.md}, \file{DATA_SOURCES.md}, \file{code/README.md} and \file{data/README.md} are \file{192-airy-*.md}. The delivered PDF, the package \file{README.md} (replaced by this report's) and \file{SHA256SUMS.json} are not shipped; the replay and build scripts check the closed delivered inventory, so they run only in a re-extracted archive. The report's README gives a route that runs the exact checker on a copy.
\end{writenote}

The optional \texttt{diagnose\_mpmath.py} evaluates finite towers and Airy
expressions at chosen working precision. It supplies illustrations only,
not interval certificates, certified decimal digits, effective remainders,
or proof of any limiting assertion. In particular a short-range fit cannot
establish the small amplitude difference in \eqref{xtw:airy:eq:ratio}.

\section{Prior sources and remaining questions}\label{xtw:airy:sec:literature}
\subsection{What is already part of the literature}
Hanna's A096537 entry defines the linear continued exponential, and A096542
explicitly gives the entire shifted equation \eqref{xtw:airy:eq:shiftfamily}, not
merely its specialization at $y=1$ \cite{oeis537,oeis542}. Neither the sequence
nor that family is introduced here as new.

The primary metadata and abstract of Bender and Vinson's 1996 paper describe
a general continued-exponential representation with coefficients at each
level and a tree-graph interpretation carrying the corresponding vertex
amplitudes \cite{bv}. Thus the broad tower/tree framework deserves prior
credit. The specialization of its general displayed tower to level
coefficients $1,2,3,\ldots$ is immediate algebra; it does not reveal the scope
of results in the unavailable article body. A full-text theorem comparison
with that paper has not been completed here.

Poland's 1998 paper treats continued exponentials in statistical mechanics
\cite{poland}. Publisher metadata identifies it as open-archive, but its full
text could not be retrieved in this environment. This is an access limitation
of this comparison, not evidence that the paper is closed or irrelevant.
In particular, the exact overlap with linear unbounded weights, coefficient
asymptotics, or Airy finite-height behavior remains unresolved for both papers.

Chapoton's inspected paper concerns hyperplane arrangements associated with
fixed labeled rooted trees. Its depth-product characteristic polynomial and
chamber-count formulas are relevant related enumerative structure
\cite{chapoton}. Its later lattice enumeration is attached to a fixed tree;
no summed coefficient asymptotic of the kind proved here was found in the
inspected text. This is a statement about that paper, not its surrounding
literature. Similarly named sequential depth-weighted attachment models
should not be identified with the static probability proportional to
\eqref{xtw:airy:eq:treeweight} without a separate equivalence argument.

A bounded search using the sequence identifiers, the proposed corrected
constant and decimal variants, and continued-exponential/Airy/depth-weighted
terms found no exact prior occurrence of the correction or the specific
finite-height theorem. Formula indexing and full-text access are incomplete.
That negative search result is not a priority certificate. The defensible
conclusion here is the proved correction to the currently displayed OEIS
amplitude, with prior-method attribution and explicit limits on literature
coverage. The primary special-function inputs are the standard Airy facts
listed in \cite{dlmf2,dlmf7,dlmf9}; the finite-tower, Fourier-tail, and
coefficient-transfer arguments are given in full in this report.

\subsection{Questions not settled by the present estimates}
The following are promising extensions, not conclusions of this report.
\begin{enumerate}
\item \textbf{Rates and further coefficients.} The entrance and matching
arguments establish uniform equivalents in the iterated-limit form required
here. They do not give an effective rate for the final coefficient equivalent.
A further term would require quantitative entrance and sensitivity errors,
a refined real or complex window, a corresponding quantitative Fourier
$L^1$ remainder, and propagation through the height sum. Formal expansions
or fitted powers alone do not supply those estimates.
\item \textbf{Fixed positive shifts.} The shifted functional equation is
already prior. Determining a complete $y$-dependent polynomial exponent and
amplitude needs a separate normalization and uniform analysis; checks of
finite polynomial rows do not prove it.
\item \textbf{Larger tilt domains and rare heights.} The compact real window
used here stops at the convenient upper bound $1$. Extending it toward Airy
zeros, or obtaining full height large-deviation equivalents away from the
central range, requires additional analysis.
\item \textbf{Effective inversion.} The shrinking sandwich is an asymptotic
existence result. A usable certified threshold algorithm with a proved error
bar would need explicit constants and an onset bound for the coefficient
error, followed by a crossing-aware integer comparison.
\item \textbf{Full primary-source comparison.} Authorized full-text copies of
Bender--Vinson and Poland would permit a page-level comparison of variable
weights, divergent formal series, critical finite depth, and coefficient
asymptotics. Until then no absolute novelty assertion is warranted.
\end{enumerate}
The central distinction is between a formal Airy balance and a coefficient
theorem. Here the normalized entrance, real-window continuation, separate
minor arcs, uniform lattice inversion, and all-height domination provide
the complete bridge. Further-order work must preserve all of those controls.
\begin{writenote}[7 October 2026]
Part~II (Report~193) answers the first-coefficient part of item~1, with $b_n\rho^n\sqrt n=C+c_1n^{-1/3}+o(n^{-1/3})$, $-0.61<c_1<-0.25$ (Theorem~\ref{xtw:frac:thm:correction}), and item~2 for every fixed positive shift, with the Gamma prefactor uniformly on compact sets of $y>0$ (Theorem~\ref{xtw:frac:thm:family}). Effective rates, further coefficients, shifts near $0$ or growing, and items~3--5 remain open; all are gathered in Section~\ref{xtw:sec:further}.
\end{writenote}

\part{Fractional corrections and Gamma laws for exponential towers: a sequel to the leading Airy asymptotic}\label{xtw:frac:part}
\partsource{Bundle Report~193, archive \file{A096537_Fractional_Corrections_and_Gamma_Laws_Source.zip}, delivered as \file{Report193.tex} (29 pages). Delivered title block ``Report193 / Fractional Corrections and Gamma Laws for Exponential Towers / A sequel to the leading Airy asymptotic / 4 October 2026''. Printed as delivered, with \textbf{[write]} notes; Report~193's Section~$k$ is Section~$k+14$ here and its statements move with their sections; its equation~($k$) is equation~($k+109$) here. The opening paragraph ``Prior constructions and scope'' is the manuscript's own. Its \texttt{Report192} is Part~I.}
\begin{partabstract}
For the known formal family $F_y(x)=\exp(xyF_{y+1}(x))$, put
$P_n(y)=n![x^n]F_y(x)$ and $\rho=1-e^{-1}$. We prove that
$P_n(1)\rho^n\sqrt n/(n!)^2=C+c_1n^{-1/3}+o(n^{-1/3})$, where
$C=\sqrt{2/(\pi\rho)}$ and $-0.61<c_1<-0.25$.
The coefficient is defined by an absolutely convergent Airy perturbation
integral. Its strict sign follows from a real-Airy reduction, a proved
positive-real tail bound, and reproducible directed interval quadrature.
The discrete proof retains the parabolic entrance logarithm, controls a
growing complex window, and cancels the larger overlap terms before Fourier
inversion and height averaging. We also prove the same relative fractional correction uniformly on compact
positive real $y$ sets, together with the Gamma prefactor
$P_n(y)/(n!)^2\sim C\rho^{1-y}\rho^{-n}n^{y-3/2}/\Gamma(y)$,
a normal limit for root-ancestor choices, and refined threshold inverses.
The first correction gives eventual approach from below, but no effective
finite-index onset or monotonicity of the normalized sequence. The report
makes no all-orders, complex Borel-sector, or absolute-priority claim.
\end{partabstract}

\paragraph{Prior constructions and scope.}
The shifted family and polynomial triangle are due to Paul D. Hanna as
recorded in OEIS A096537 and A096542 \cite{oeis537,oeis542}. General continued
exponentials and their tree interpretation belong to the earlier framework
of Bender and Vinson \cite{bv}. A complete full-text theorem comparison with
that paper and Poland's continuation \cite{poland} remains unresolved.
The formulas proved here are not presented as a worldwide novelty claim.
Report192 established the leading unshifted Airy coefficient theorem;
Section~\ref{xtw:frac:sec:input} states exactly the leading input retained here.
The new error estimates, sign reduction, and offset normalization are proved
in the present text rather than inferred from fitted coefficients.

\section{Statements and normalizations}\label{xtw:frac:sec:statements}
All logarithms are natural. Set
\begin{equation}\label{xtw:frac:eq:constants}
 \alpha=e-1,\quad \rho=\alpha/e,\quad a=2^{-1/3},\quad
 C=\sqrt{\frac2{\pi\rho}},\quad
 \kappa=2^{1/6}\sqrt{\frac e\pi}.
\end{equation}
The equations
\begin{equation}\label{xtw:frac:eq:family}
 F_y(x)=\exp\bigl(xyF_{y+1}(x)\bigr),\qquad F_y(0)=1
\end{equation}
define a unique formal family: degree $n$ on the right needs degrees only
through $n-1$ at the next level. We write
\[
 P_n(y)=n![x^n]F_y(x),\qquad b_n(y)=\frac{P_n(y)}{(n!)^2},
 \qquad a_n=P_n(1),\quad b_n=b_n(1).
\]
No analytic convergence of an infinite positive-real tower is used.
Every analytic tower evaluation below is finite and entire.

\begin{theorem}[The first fractional correction]\label{xtw:frac:thm:correction}
As $n\to\infty$,
\begin{equation}\label{xtw:frac:eq:main}
 b_n\rho^n\sqrt n=C+c_1n^{-1/3}+o(n^{-1/3}),
 \qquad -0.61<c_1<-0.25.
\end{equation}
The constant $c_1$ is the absolutely convergent integral in
\eqref{xtw:frac:eq:c1definition}, equivalently the real integral in
\eqref{xtw:frac:eq:realconstant}. Thus $b_n\rho^n\sqrt n<C$ for all sufficiently
large $n$, with a gap of order $n^{-1/3}$.
\end{theorem}
The bracket is a rigorous enclosure under the stated arithmetic trust
boundary, not a certified decimal expansion. The ordinary numerical value
$-0.43652529$ is only a diagnostic. No finite onset of the strict inequality,
no monotonicity of $b_n\rho^n\sqrt n$, and no explicit numerical error bound at a specified finite $n$
is asserted.

\begin{theorem}[The Gamma prefactor]\label{xtw:frac:thm:family}
For every compact $Y\subset(0,\infty)$,
\begin{equation}\label{xtw:frac:eq:gammaprefactor}
 \begin{gathered}
 b_n(y)=C_y\rho^{-n}n^{y-3/2}\left[1+d n^{-1/3}+o_Y(n^{-1/3})\right],
 \qquad y\in Y,\\
 C_y=\frac{C\rho^{1-y}}{\Gamma(y)},\qquad d=\frac{c_1}{C}.
 \end{gathered}
\end{equation}
The relative coefficient $d$ is independent of $y$. The same critical Airy
size law applies to every fixed positive shift.
For the root-ancestor statistic $K_n$ defined in
Section~\ref{xtw:frac:sec:ancestor}, with $L_n=\log(n/\rho)$,
\begin{equation}\label{xtw:frac:eq:pgfclt}
 \frac{P_n(y)}{P_n(1)}=\frac{(n/\rho)^{y-1}}{\Gamma(y)}
                   [1+o_Y(n^{-1/3})],
 \qquad
 \frac{K_n-L_n}{\sqrt{L_n}}\xrightarrow{d}\mathcal N(0,1).
\end{equation}
The ratio equivalent is uniform on every such $Y$.
\end{theorem}
Uniformity near $y=0$ and shifts growing with $n$ are outside the theorem.
A separate quantitative local complex argument in Section~\ref{xtw:frac:sec:cumulants}
justifies mean and variance constants; these are not obtained by
differentiating a merely real asymptotic.
\begin{writenote}[7 October 2026]
At $y=1$, $C_y=C$ and Theorem~\ref{xtw:frac:thm:family} is Part~I's Theorem~\ref{xtw:airy:thm:main} with the correction of Theorem~\ref{xtw:frac:thm:correction} added. The theorem answers Part~I's question ``Fixed positive shifts'' (Section~\ref{xtw:airy:sec:literature}) for every fixed $y>0$.
\end{writenote}

\subsection{Definition of the correction without numerical fitting}
For $\lambda\in\ii\R$ let
\begin{align}
 w_\lambda''(t)&=\tfrac12(t-\lambda)w_\lambda(t),
 &w_\lambda(0)&=0,&w_\lambda'(0)&=1,\label{xtw:frac:eq:wode}\\
 U_\lambda(t)&=2w_\lambda'(t)/w_\lambda(t),
 &G_\lambda(t)&=\tfrac12-\tfrac12(t-\lambda)U_\lambda(t)
                              -\tfrac1{12}U_\lambda(t)^3.
\end{align}
The nonvanishing needed here is proved in Section~\ref{xtw:frac:sec:complex}.
For any constant $c$, define
\begin{equation}\label{xtw:frac:eq:Vdefinition}
 V_c(t)=\frac1{w_\lambda(t)^2}
 \left\{c-\frac23\log t+
 \int_0^t\left[w_\lambda(r)^2G_\lambda(r)+\frac2{3r}\right]dr\right\}.
\end{equation}
The integral is convergent, and $V_c'+U_\lambda V_c=G_\lambda$.
Let $I_c'=V_c$ have the uniquely fixed finite-part normalization
\begin{equation}\label{xtw:frac:eq:Idefinition}
 I_c(t)-\frac{\frac23(\log t+1)-c}{t}\longrightarrow0
 \quad(t\downarrow0),
\end{equation}
and set, with the principal square root,
\begin{equation}\label{xtw:frac:eq:Jdefinition}
 J_c(\lambda)=\lim_{T\to\infty}
 \left[I_c(T)+\frac{(T-\lambda)^2}{3}
                    -\frac43\sqrt{2(T-\lambda)}\right].
\end{equation}
Existence and polynomial growth on the imaginary axis are established below.
Put
\begin{equation}\label{xtw:frac:eq:FQ}
 \FF(\lambda)=\frac{\kappa}{\Ai(-a\lambda)^2},\qquad
 \QQ(\lambda)=\FF(\lambda)\left[\frac{\lambda^2}{2}-J_0(\lambda)\right].
\end{equation}
Then the exact first correction is
\begin{equation}\label{xtw:frac:eq:c1definition}
 c_1=\frac{\alpha^{-1/6}}{2\pi}\int_{\R}\QQ(\ii\tau)\,d\tau.
\end{equation}
Conjugation makes this integral real. The definition contains neither fitted
finite-sequence data nor an arbitrary normalization constant.

\section{Exact heights and the leading input}\label{xtw:frac:sec:input}
For $y>0$, write $\beta=y-1$, $\ell_j=j+\beta$, and define
\begin{equation}\label{xtw:frac:eq:finite}
 A_{h+1,h;y}=1,\qquad
 A_{j,h;y}(x)=\exp\bigl(\ell_jxA_{j+1,h;y}(x)\bigr),\qquad
 E_{h,y}=A_{1,h;y}-A_{1,h-1;y}.
\end{equation}
The second tower is $1$ at level $h$; $A_{1,0;y}=1$.
Every $E_{h,y}$ has nonnegative coefficients and least degree $h$:
at the bottom the difference is $e^{\ell_hx}-1$, and each upward step
multiplies the difference by $x$ before applying
$e^u-e^v=e^v(e^{u-v}-1)$. Thus, with $c_{n,h;y}=[x^n]E_{h,y}$,
\begin{equation}\label{xtw:frac:eq:heightsum}
 b_n(y)=\sum_{h=1}^n\frac{c_{n,h;y}}{n!}\quad(n\ge1).
\end{equation}
For $y=1$ we abbreviate $E_{h,1}=H_h$ and $c_{n,h;1}=c_{n,h}$.

\paragraph{Leading-input recap.}
The unshifted leading result from Report192 is
\begin{align}
 H_h((eh)^{-1})&\sim\FF(0)h^{1/6}e^{-\alpha h},\label{xtw:frac:eq:leadingheight}\\
 c_{n,h}/n!&\le B\frac{e^n}{n!}h^{n-1/2}e^{-\alpha h}
                      \qquad(h\ge1,n\ge0),\label{xtw:frac:eq:allheight}
\end{align}
for an absolute $B$. The centered finite-height characteristic transform
has a bound $B e^{-b|t|^{3/2}}$ beyond a fixed central interval under
$\theta=h^{-2/3}t$, until the remote arc, whose contribution is exponentially
small in $h$. The Airy inverse
\[
 f(z)=\frac1{2\pi}\int_\R e^{-\ii tz}\FF(\ii t)\,dt
\]
has bounded derivatives of every fixed order and
$f(0)=\kappa/a=\sqrt{2e/\pi}$. This recap is a leading theorem, not a
first-order remainder assumption. The present argument also recovers these
leading facts from its quantitative window and the independent positive
factorization below, so they can be read as orientation rather than
separate assumptions. Sections~\ref{xtw:frac:sec:entrance}--\ref{xtw:frac:sec:transfer}
prove the additional quantitative estimates needed for
Theorem~\ref{xtw:frac:thm:correction}. For completeness, the positive-factor minor-arc
argument and its shifted version are given in Section~\ref{xtw:frac:sec:minor};
Sections~\ref{xtw:frac:sec:offset}--\ref{xtw:frac:sec:familytransfer} give the full normalization
and height sum underlying Theorem~\ref{xtw:frac:thm:family}.
\begin{writenote}[7 October 2026]
Report192 is Part~I. The recap is Part~I's Theorem~\ref{xtw:airy:thm:critical} (with $\FF(0)=\kappa/\Ai(0)^2=K$), its all-height bound \eqref{xtw:airy:eq:allheight}, its minor-arc Theorem~\ref{xtw:airy:thm:minor} and its density, here as $f=K\,p$ with $p$ of \eqref{xtw:airy:eq:pdensity}. Where this Part re-derives these facts inside its quantitative argument, the proofs are printed in full, as their estimates are needed for the correction.
\end{writenote}

\paragraph{Airy facts used.}
The Wronskian and connection identities \cite{dlmf2} give
\begin{equation}\label{xtw:frac:eq:wairy}
 w_\lambda(t)=\frac\pi a\left[
 \Ai(-a\lambda)\Bi(a(t-\lambda))-
 \Bi(-a\lambda)\Ai(a(t-\lambda))\right].
\end{equation}
We use the separate function and derivative sector expansions, with
remainders, from \cite{dlmf7}; asymptotic equivalents are never
differentiated without those derivative expansions. The zeros of $\Ai$
are negative real \cite{dlmf9}. On the imaginary axis the function
$\FF$, its fixed-order derivatives, and its products with fixed polynomials
are bounded by a polynomial times $e^{-b|\tau|^{3/2}}$.
Finally, if $r(\lambda)=\Bi(-a\lambda)/\Ai(-a\lambda)$, then
\[
 r'(\lambda)=-\frac a{\pi\kappa}\FF(\lambda),\quad
 r(\ii\infty)=-\ii,\quad r(-\ii\infty)=\ii.
\]
Integrating the derivative upward proves $f(0)=\kappa/a$ and also shows
that $r$ is bounded on the imaginary axis.

\section{The normalized parabolic entrance}\label{xtw:frac:sec:entrance}
Set $\epsilon=h^{-1/3}$ and evaluate at
$x=e^{\lambda\epsilon^2}/(eh)$, where $\lambda=\ii\tau$.
Use exponent variables $q_j^A=jxA_{j+1,h}(x)$ and
$q_j^B=jxA_{j+1,h-1}(x)$, so no complex logarithm is needed to define them.
With $u_k^p=1-q_{h-k}^p$, the exact branches are
\begin{equation}\label{xtw:frac:eq:recurrence}
 u_k^p=1-e^{\lambda\epsilon^2}(1-k/h)e^{-u_{k-1}^p},\quad
 u_0^A=1-e^{\lambda\epsilon^2}/e,\quad u_0^B=1.
\end{equation}
Write $\bar u_k=(u_k^A+u_k^B)/2$ and
$\Delta_k=u_k^B-u_k^A$.

\begin{lemma}[Autonomous normalization]\label{xtw:frac:lem:fatou}
Let $v_0=1$ and $v_{k+1}=1-e^{-v_k}$. There is a finite constant $c_*$ such that
\begin{align}
 v_k&=\frac2k-\frac{\frac23\log k+2c_*}{k^2}
                       +\OO(\log^2k/k^3),\label{xtw:frac:eq:fatou}\\
 \frac{v_k+v_{k+1}}2&=\frac2k-
        \frac{\frac23\log k+2c_*+1}{k^2}+\OO(\log^2k/k^3),\\
 v_k-v_{k+1}&=\frac2{k^2}-
        \frac{\frac43\log k+4c_*+\frac43}{k^3}
                     +\OO(\log^2k/k^4).
\end{align}
\end{lemma}
\begin{proof}
The decreasing positive sequence tends to zero. For $r_k=2/v_k$, Taylor's
theorem gives $r_{k+1}=r_k+1+1/(3r_k)+\OO(r_k^{-3})$.
First $r_k\sim k$, then summation gives $r_k=k+\OO(\log k)$.
The increments of $r_k-k-\frac13\log k$ are therefore
$\OO(\log k/k^2)$, an absolutely summable sequence. Its tail is
$\OO(\log k/k)$, proving convergence to $c_*$ and \eqref{xtw:frac:eq:fatou}.
For the difference use
$v_k-v_{k+1}=v_k^2/2-v_k^3/6+\OO(v_k^4)$ rather than subtracting error
terms. The midpoint follows by subtracting half this difference.
\end{proof}

Take the quantitative cutoffs
\begin{equation}\label{xtw:frac:eq:cutoffs}
 K=\ceil{\epsilon^{-3/5}},\quad t_0=\epsilon K,\quad
 M=\ceil{\epsilon^{-11/8}},\quad T=\epsilon M,\quad i=h-M,
\end{equation}
and the growing Fourier window
\begin{equation}\label{xtw:frac:eq:window}
 |\tau|\le L_\epsilon=A\bigl(\log(1/\epsilon)\bigr)^{2/3}.
\end{equation}
Here $A$ is any fixed positive constant, chosen later large enough for the
Fourier tails. In the remainder estimates below, a factor called
$\operatorname{polylog}(1/\epsilon)$ means a fixed positive polynomial in
$1+\log(1/\epsilon)$, with degree independent of $h,\lambda$.
Every parameter-dependent constant is polynomial in $L_\epsilon$; no
uncontrolled exponential of $L_\epsilon$ is allowed.

\begin{lemma}[Early perturbation and separation]\label{xtw:frac:lem:early}
Uniformly on \eqref{xtw:frac:eq:window},
\begin{align}
 \bar u_K&=\frac{v_K+v_{K+1}}2-\lambda\epsilon^2K/3
       +\OO(L_\epsilon\epsilon^2\log K+\epsilon^3K^2
                         +L_\epsilon^2\epsilon^4K^3),\label{xtw:frac:eq:earlyu}\\
 \log\Delta_K&=\log2-2\log K
       -\frac{\frac23\log K+2c_*+\frac23}{K}
       +\frac{\lambda\epsilon^2K^2}{6}
       +\OO(\epsilon^{6/5}\operatorname{polylog}(1/\epsilon)).
                                                        \label{xtw:frac:eq:earlyD}
\end{align}
The logarithm is continued from the positive real value at $\lambda=0$.
\end{lemma}
\begin{proof}
Put $v_k^B=v_k$, $v_k^A=v_{k+1}$, $d_k^p=u_k^p-v_k^p$, and
$b_k^p=e^{-v_{k-1}^p}$. Lemma~\ref{xtw:frac:lem:fatou} gives the summable kernel
\begin{equation}\label{xtw:frac:eq:earlykernel}
 \prod_{r=j+1}^k b_r^p\le B\left(\frac{j+1}{k+1}\right)^2.
\end{equation}
Writing $a_k=e^{\lambda\epsilon^2}(1-k\epsilon^3)$, Taylor's theorem yields
\[
 d_k=b_k(d_{k-1}-(a_k-1))
          +\OO(|a_k-1||d_{k-1}|+|d_{k-1}|^2).
\]
The kernel proves by bootstrap
$|d_k|\le B(|\lambda|\epsilon^2(k+1)+\epsilon^3(k+1)^2)$ for $k\le K$:
the nonlinear contribution relative to this bound is
$\OO((1+|\lambda|)\epsilon^2K^2+\epsilon^3K^3)=o(1)$.
The initial errors are $d_0^B=0$ and
$d_0^A=-\lambda\epsilon^2/e+\OO(|\lambda|^2\epsilon^4)$.
The linear response $l_k=b_k(l_{k-1}-1)$, with $l_0^B=0,l_0^A=-1/e$,
satisfies $l_k=-k/3+\OO(\log(k+2))$: substituting $-k/3$ gives residual
$\OO(\log(k+2)/(k+1))$, and \eqref{xtw:frac:eq:earlykernel} sums it.
Subtract $\lambda\epsilon^2l_k$ and sum the remaining forcing to obtain
\[
 d_k=-\lambda\epsilon^2k/3+
 \OO(|\lambda|\epsilon^2\log(k+2)+\epsilon^3(k+1)^2
                  +|\lambda|^2\epsilon^4(k+1)^3).
\]
In particular $\sum_{k<K}\bar d_k=-\lambda\epsilon^2K(K-1)/6+
\OO(\epsilon^{6/5}\operatorname{polylog})$.

For $r(z)=\log(\sinh(z/2)/(z/2))$ the exact identity is
\begin{equation}\label{xtw:frac:eq:exactD}
 \log\Delta_K=\log\Delta_0+\sum_{j=1}^K\log a_j
                   -\sum_{k=0}^{K-1}\bar u_k+\sum_{k=0}^{K-1}r(\Delta_k).
\end{equation}
Let $\Delta_k^*=v_k-v_{k+1}$. The same bootstrap gives
$|\log(\Delta_k/\Delta_k^*)|\le
 B(|\lambda|\epsilon^2(k+1)^2+\epsilon^3(k+1)^3)$.
Since $r'(z)=\OO(|z|)$ in a fixed neighborhood of this trajectory and
$\Delta_k^*=\OO((k+1)^{-2})$, the difference of the $r$ sums is at most
$B(|\lambda|\epsilon^2+\epsilon^3\log(K+2))$.
Subtract the autonomous identity in \eqref{xtw:frac:eq:exactD}. The initial
logarithm and $\sum\log a_j$ cost only
$\OO(|\lambda|\epsilon^2K+\epsilon^3K^2)$.
The midpoint sum provides the displayed tilt; expanding
$\log\Delta_K^*$ with Lemma~\ref{xtw:frac:lem:fatou} proves \eqref{xtw:frac:eq:earlyD}.
All errors stated here have at least the $\epsilon^{6/5}$ margin up to
fixed logarithmic factors.
\end{proof}

The forced entrance parameter is
\begin{equation}\label{xtw:frac:eq:ch}
 c_h=-\frac23\log(1/\epsilon)-2c_*-1.
\end{equation}
It is essential at finite height. Its eventual disappearance from the
coefficient correction is proved by Fourier inversion, not assumed.

\section{Complex stability and quantitative inner propagation}\label{xtw:frac:sec:complex}
\begin{lemma}[Energy stability on the whole imaginary axis]\label{xtw:frac:lem:energy}
For $\lambda=\ii\tau$, $\tau\in\R$, the solution $w_\lambda(t)$ is nonzero
for $t>0$, and
\begin{align}
 \Re U_\lambda(t)&\ge2/t,&
 \Re U_\lambda(t)&\ge c\sqrt t\quad(t\ge1),&
 \Re U_\lambda(t)&\ge U_0(t),\label{xtw:frac:eq:energybounds}\\
 |U_\lambda(t)-U_0(t)|&\le|\tau|t/3,&
 |U_\lambda(t)-U_0(t)|&\le B|\tau|/\sqrt t\quad(t\ge1).
                                                        \label{xtw:frac:eq:Udiff}
\end{align}
All constants are absolute.
\end{lemma}
\begin{proof}
Integrate $w''\bar w$ and take real parts:
\begin{equation}\label{xtw:frac:eq:energyidentity}
 \Re(w'(t)\overline{w(t)})=
 \int_0^t\left(|w'(s)|^2+\frac s2|w(s)|^2\right)ds.
\end{equation}
A positive zero would contradict $w'(0)=1$. Cauchy--Schwarz gives
$\int_0^t|w'|^2\ge|w(t)|^2/t$, proving the first bound.
For $t\ge2$, restrict the integral to $[t/2,t]$ and use the trace inequality
\[
 \int_A^B(|v'|^2+q^2|v|^2)\,ds
       \ge q\tanh(q(B-A))|v(B)|^2.
\]
To prove it, subtract
$v(B)\cosh(q(s-A))/\cosh(q(B-A))$ and integrate the cross term by parts;
the remainder energy is nonnegative. Taking $q=\sqrt t/2$ proves the
second bound; $1\le t\le2$ follows from the first.
Write $U=R+\ii I$. Its Riccati equation gives
\[
 (R-U_0)'=-\tfrac12(R+U_0)(R-U_0)+\tfrac12I^2.
\]
Variation of constants from $\delta>0$ and then $\delta\downarrow0$
proves $R\ge U_0$: the initial discrepancy is $\OO_\tau(\delta^3)$
and its propagator is at most $(\delta/t)^2$.
Lastly $D=U_\lambda-U_0$ satisfies
\[
 D(t)=-\ii\tau\int_0^t
 \exp\left[-\frac12\int_s^t(U_\lambda(r)+U_0(r))dr\right]ds.
\]
The kernel modulus is at most $(s/t)^2$, giving $|\tau|t/3$.
For $t\ge2$, split at $t/2$ and apply the square-root damping to bound its
integral by $B/\sqrt t$. The compact remaining interval is immediate.
\end{proof}
Integration of $\Re U\ge2/t$ also gives $|w(t)|\ge t$.
For the discrete towers, positivity of their exponent coefficients implies
\begin{equation}\label{xtw:frac:eq:positivitycomplex}
 |q_j(x_0e^{\ii\theta})|\le q_j(x_0),\qquad
 \Re(1-q_j(x_0e^{\ii\theta}))\ge1-q_j(x_0),
 \quad x_0=(eh)^{-1}.
\end{equation}
Thus a complex tilt does not produce an exponential Gronwall loss.

\begin{lemma}[Summable inner error]\label{xtw:frac:lem:inner}
With \eqref{xtw:frac:eq:cutoffs}--\eqref{xtw:frac:eq:ch}, set
$a_k=\epsilon U_\lambda(\epsilon k)+\epsilon^2V_{c_h}(\epsilon k)$.
Uniformly on \eqref{xtw:frac:eq:window},
\begin{align}
 \sum_{k=K}^M|\bar u_k-a_k|
      &=\OO(\epsilon^{17/16}\operatorname{polylog}(1/\epsilon)),
                                                        \label{xtw:frac:eq:innersum}\\
 |\bar u_M-a_M|&=\OO(\epsilon^3T^{3/2}\operatorname{polylog}(1/\epsilon)),
                                                        \label{xtw:frac:eq:innerend}\\
 \log\Delta_M&=\log\frac{2\epsilon^2}{w_\lambda(T)^2}
                    +\epsilon D_{c_h}(T)
                    +\OO(\epsilon^{17/16}\operatorname{polylog}(1/\epsilon)),
                                                        \label{xtw:frac:eq:innersep}\\
 D_c(t)&=-I_c(t)-t^2/2+\lambda t+U_\lambda(t)/2.
                                                        \label{xtw:frac:eq:Ddefinition}
\end{align}
The value $T=\epsilon M$ includes the integer rounding exactly.
\end{lemma}
\begin{proof}
We give the residual and kernel estimates, since a fixed-cutoff $o(1)$ would
not suffice. On $0<t\le1$, put $B_0=-(2/3)\log t/t^2$ and
$d=U-2/t$. From Lemma~\ref{xtw:frac:lem:energy},
$|d|\le B(1+|\tau|)t$. Expanding $U^3-8/t^3$ shows
\[
 |G-(B_0'+UB_0)|\le B(1+|\tau|)^3(1+|\log t|)/t.
\]
The linear $V$ kernel has modulus at most $(s/t)^2$, so
$|V_0-B_0|\le B(1+|\tau|)^3(1+|\log t|)$.
Add $c_h/w^2$. On $t\ge1$ the same kernel decays at least as
$e^{-c(t^{3/2}-s^{3/2})}$ and $|G|\le P(1+|\tau|)t^{3/2}$.
It follows that
\begin{equation}\label{xtw:frac:eq:Vbound}
 |V_{c_h}(t)|\le P(1+|\tau|)
 \left[\frac{1+|c_h|+|\log t|}{t^2}+1+t\right].
\end{equation}
Here and below $P$ may change, but is a fixed polynomial. Finite
differentiation of the exact differential equations gives
\[
 |U'''|\le P(1+|\tau|)t^{-4},\quad
 |V''|\le P(1+|\tau|)(1+|c_h|+|\log t|)t^{-4}\quad(t\le1).
\]
For $t\ge1$, $U_0=\OO(\sqrt t)$ and $U_0'=\OO(1)$ follow from the
separate Airy expansions and compactness. The difference equation gives
$|U'|\le P(1+|\tau|)$, hence $|U'''|\le P(1+|\tau|)t$;
the $V$ equation gives $|V'|\le P(1+|\tau|)(1+|c_h|)t^{3/2}$ and
$|V''|\le P(1+|\tau|)(1+|c_h|)t^2$. The factor $1+|c_h|$
is absorbed in the fixed polylogarithmic factor in every residual estimate.
More precisely, the homogeneous term and its first two derivatives on
$t\ge1$ are bounded by $|c_h|P(1+|\tau|)P(t)e^{-ct^{3/2}}$;
this is also bounded by $|c_h|P(1+|\tau|)t^{-4}$ after enlarging
constants, as needed in the displayed residual bound.

Taylor expansion of $a_{k-1}$ through $\epsilon^3$ cancels the coefficients
of $\epsilon^2$ and $\epsilon^3$, because
\[
 U'=t-\lambda-U^2/2,\qquad
 V'+UV=U''/2-U^3/3=G.
\]
Thus the scalar recurrence residual at $t=\epsilon k$ is bounded by
\begin{equation}\label{xtw:frac:eq:innerresidual}
 P(1+|\tau|)\epsilon^4
 \left[\frac{(1+|c_h|+|\log t|)^2}{t^4}+1+t^2\right].
\end{equation}
The exact midpoint recurrence adds a factor
$\cosh(\Delta_{k-1}/2)$. Coefficient positivity and the more strongly
contracted critical-real pair give
$|\Delta_k(\ii\tau)|\le\Delta_k(0)\le v_k-v_{k+1}\le B/(k+1)^2$.
Its extra residual is $\OO(\epsilon^4/t^4)$ and is already included.
The same bound ensures that the sinh quotient defining the separation
logarithm has no zero along the chosen continuation.

For $K\le k\le\epsilon^{-1}$, the exact critical-real midpoint is at least
$v_{k+1}\ge3/(2(k+1))$ for sufficiently large $K$. The approximation has
the same lower real-part bound: \eqref{xtw:frac:eq:Vbound} gives
$\epsilon|V|/\Re U=\OO(\epsilon^{3/5}\operatorname{polylog})$ on
$[t_0,1]$. Future propagation kernels therefore have total mass $\OO(k)$.
For later indices use the explicit barrier $b_k=\frac14\sqrt{k/h}$.
For $0\le u\le1$ and $0\le s\le1$,
\[
 1-(1-s)e^{-u}\ge u+s-su-u^2/2.
\]
For $k\ge\ceil{h^{1/3}}$ and $k\le M$, $b_k\le1/10$ eventually,
so the right side at $u=b_{k-1},s=k/h$ exceeds
$b_{k-1}+(4/5)k/h$. Since
$b_k-b_{k-1}\le1/(4\sqrt{hk})\le k/(4h)$, this is a subsolution.
The autonomous starting bound proves it for both exact real branches.
Equation~\eqref{xtw:frac:eq:positivitycomplex} transfers it to their complex real
parts. For the approximation,
$\epsilon|V|/\Re U=\OO(\epsilon\sqrt T\operatorname{polylog})=o(1)$,
so its real part has the same contraction scale. Future kernel mass at
$t_k\ge1$ is $\OO((\epsilon\sqrt{t_k})^{-1})$.

The early lemma and the small-$t$ expansions
\begin{equation}\label{xtw:frac:eq:smallUV}
 U=2/t-\lambda t/3+t^2/4+\OO(P(1+|\lambda|)t^3),\qquad
 V_c=\frac{c-(2/3)\log t}{t^2}+\OO(P(1+|\lambda|)(1+|c|+|\log t|))
\end{equation}
give
\[
 K|\bar u_K-a_K|\le P(1+|\lambda|)
 [\log^2K/K^2+\epsilon^2K\log(1/\epsilon)
                      +\epsilon^3K^3+\epsilon^4K^4]
 =\OO(\epsilon^{6/5}\operatorname{polylog}).
\]
Sum \eqref{xtw:frac:eq:innerresidual} against the kernel masses. The total error is
at most this entrance error plus fixed polynomial factors times
\[
 \epsilon^2\int_{t_0}^1
 (1+|c_h|+|\log t|)^2t^{-3}\,dt
 +\epsilon^2\int_1^T(1+t^{3/2})dt.
\]
These are respectively $\OO(\epsilon^{6/5}\operatorname{polylog})$ and
$\OO(\epsilon^2T^{5/2}\operatorname{polylog})
=\OO(\epsilon^{17/16}\operatorname{polylog})$.
For the endpoint, forcing before $T/2$ is superpolynomially damped, and
the later residual divided by its local contraction gap is
$\OO(\epsilon^3T^{3/2}\operatorname{polylog})$.

Finally the exact multiplier is
\[
 \log(\Delta_k/\Delta_{k-1})=
 \lambda\epsilon^2+\log(1-k/h)-\bar u_{k-1}+r(\Delta_{k-1}).
\]
The $r$ tail is $\OO(K^{-3})$; the quadratic log remainder sums to
$\OO(\epsilon^3T^3)$, both $o(\epsilon)$. Replace the midpoint by $a_k$
using \eqref{xtw:frac:eq:innersum}. Euler--Maclaurin yields its endpoint correction
$+\epsilon U(T)/2$; the remainders are
$\OO(\epsilon^2t_0^{-2}\operatorname{polylog}+\epsilon^2T)$.
The leading integral is $2\log w$.
At $t_0$, the early separation matches \eqref{xtw:frac:eq:Ddefinition}: its singular
term is $[-(2/3)\log K-2c_*-2/3]/K$, and $w$ supplies
$\lambda\epsilon^2K^2/6$.
The omitted terms $t_0^3,\lambda^2t_0^4$, and
$\epsilon t_0\operatorname{polylog}$ are $o(\epsilon)$ with the stated
margin. This fixes the additive constant and proves \eqref{xtw:frac:eq:innersep}.
As a check, $D_c'=-V_c-(t-\lambda)/2-U^2/4$ exactly.
\end{proof}

\section{Second outer correction and overlap cancellation}\label{xtw:frac:sec:outer}
Let $\TT(z)e^{-\TT(z)}=z$ be the small tree branch and $v=1-\TT$.
Write $\eta=e^{\lambda\epsilon^2}/(eh)$, $z_j=j\eta$, and define
\begin{equation}\label{xtw:frac:eq:RS}
 R=\frac{\TT e^{\TT}}{v^2},\qquad
 S=\frac{\TT e^{2\TT}(5+v-v^2)}{2v^5},\qquad
 Q_{\rm out}=\TT+\eta R+\eta^2S.
\end{equation}
Direct differentiation verifies
\begin{equation}\label{xtw:frac:eq:Sidentity}
 S=\frac{\TT}{v}\left[\frac{\TT''}{2}+R'
                                      +\frac{(\TT'+R)^2}{2}\right].
\end{equation}
All derivatives here are in $z$. This is the second cancellation in
$Q_{\rm out}(z)=ze^{Q_{\rm out}(z+\eta)}$.

\begin{lemma}[Summed outer control]\label{xtw:frac:lem:outer}
For either exact branch and $j\le i=h-M$,
\begin{align}
 |Q_{{\rm out},j}-z_je^{Q_{{\rm out},j+1}}|
     &\le B|\eta|^3|\TT_j|/|v_j|^7,\label{xtw:frac:eq:outerresidual}\\
 \sum_{j\le i}|q_j-Q_{{\rm out},j}|
     &\le B\frac{|q_i-Q_{{\rm out},i}|}{|v_i|}
                 +B|\eta|^2|v_i|^{-6}=o(\epsilon).
                                                        \label{xtw:frac:eq:outercontrol}
\end{align}
The error is uniform on \eqref{xtw:frac:eq:window}.
\end{lemma}
\begin{proof}
Positivity of $\TT$ and its derivative bound show
\[
 |\TT(z_j)-\TT(|z_j|)|
       \le |\arg\eta|\,|z_j|\TT'(|z_j|).
\]
The relative perturbation of the real gap is $\OO(L_\epsilon/T)=o(1)$.
Thus all gaps lie in a fixed narrow cone and are comparable to the real
ones, while $|\eta|/|v_i|^3=\OO(T^{-3/2})$.
Analytic Taylor's theorem, \eqref{xtw:frac:eq:Sidentity}, and the derivative bounds
for rational functions of $\TT,v$ give \eqref{xtw:frac:eq:outerresidual}.
The actual $q_i$ has modulus at most its real value, at most
$1-c\epsilon\sqrt T$. The approximation lies in the same type of disk,
since $\eta R/v=\OO(T^{-3/2})$ and $\eta^2S/v=\OO(T^{-3})$.
Choose a radius $r$ between $\TT(|z_i|)$ and $1$ large enough for both.
Then $|z_i|e^r\le r$, so the disk is invariant under descent; each
mean-value derivative has modulus at most $1-c'|v_i|$.
Summing the residuals along the ray, then the propagation kernel, gives
\eqref{xtw:frac:eq:outercontrol}; the residual contribution is $\OO(T^{-3})$.
The change of variable $dz=e^{-\TT}v\,d\TT$ justifies the powers in this
integral estimate, and neighboring gaps are comparable because
$|\eta|/|v_i|^2=o(1)$.

For the incoming error put $s=T-\lambda$. The bounded connection ratio
$r(\lambda)$ and the separate sector expansions give, uniformly,
\begin{equation}\label{xtw:frac:eq:Ularge}
 U(T)=\sqrt{2s}-\frac1{2s}-\frac5{8\sqrt2\,s^{5/2}}+\OO(s^{-4}),\qquad
 V_{c_h}(T)=-2s/3+\OO(\operatorname{polylog}/\sqrt s).
\end{equation}
Expansion of $v_i-\eta R_i-\eta^2S_i$ gives exactly the three displayed
$\epsilon U$ terms and the leading $\epsilon^2V$ term.
Using \eqref{xtw:frac:eq:innerend}, the incoming discrepancy is bounded by
\[
 \operatorname{polylog}(1/\epsilon)
 [\epsilon T^{-4}+\epsilon^2T^{-1/2}+\epsilon^3T^{3/2}].
\]
Division by $|v_i|\asymp\epsilon\sqrt T$ gives
$\OO(\operatorname{polylog}[T^{-9/2}+\epsilon/T+\epsilon^2T])=o(\epsilon)$.
The residual $T^{-3}=\OO(\epsilon^{9/8})$ is also $o(\epsilon)$.
\end{proof}

\subsection{The equilibrium product at a moving endpoint}
Let $P=\prod_{j=1}^{i-1}\TT(z_j)$. The exact identity
$\log(q_j/\TT_j)=q_{j+1}-\TT_j$ and Lemma~\ref{xtw:frac:lem:outer} imply
\begin{equation}\label{xtw:frac:eq:ratioproduct}
 \log\prod_{j<i}\frac{q_j}{\TT_j}
   =-\log v_i+\eta\int_0^{z_i}S(z)dz+o(\epsilon).
\end{equation}
Indeed $\int R\,dz=-\TT-\log v$ cancels the endpoint $+\TT_i$;
the lower endpoint is $\OO(\eta)$, the $R$ discretization error is
$\OO(\epsilon/T)$, and the $S$ error is
$\OO(\epsilon T^{-5/2})$. The separation at $i$ is exponentially small
in $T^{3/2}$, so its divided-difference product has the same logarithm as
either individual branch product to the required accuracy.

An overlap term larger than the desired correction must be retained:
\begin{align}
 \eta\int_0^{z_i}S(z)dz
   &=\frac5{6(2s)^{3/2}}+\OO(\epsilon/T+T^{-3}),\label{xtw:frac:eq:Soverlap}\\
 \log w_\lambda(T)^{-2}
   &=\log\frac{a^{5/2}}{\pi\Ai(-a\lambda)^2}
       +\tfrac12\log s-\frac{2\sqrt2}{3}s^{3/2}
       -\frac5{6(2s)^{3/2}}+\OO(T^{-3}).\label{xtw:frac:eq:woverlap}
\end{align}
Each $s^{-3/2}$ term has size $\epsilon^{9/16}$, greater than
$\epsilon$. They cancel; discarding either would invalidate the proof.

Put $\widehat H=he^{-\lambda\epsilon^2}$ and $r=i-1$. Isolating
$\log\TT(z)=\log z+\TT(z)$ permits exact factorial summation at zero.
Stirling and trapezoidal Euler--Maclaurin yield
\begin{equation}\label{xtw:frac:eq:productEM}
 \log P=\widehat H\int_0^{r/\widehat H}\log\TT(v/e)dv
      +\tfrac12\log(2\pi r)+\tfrac12\TT(r/(e\widehat H))+o(\epsilon).
\end{equation}
For the smooth part the error is
$\OO(|\eta|\int|\TT''|\,|dz|)=\OO(|\eta|/|v_i|)
=\OO(\epsilon^2/\sqrt T)$, by gap-cone comparability.
This is a moving-endpoint bound, stronger than a fixed-endpoint $o(1)$.
Since
\[
 \int_0^1\log\TT(v/e)dv=-\alpha,\qquad
 \log\TT((1-y)/e)=-\sqrt{2y}-y/3+\OO(y^{3/2}),
\]
and
\[
 1-r/\widehat H=\epsilon^2s+\epsilon^3
      +\epsilon^4(\lambda T-\lambda^2/2)+\OO(\epsilon^5\operatorname{polylog}),
\]
expansion of the removed endpoint integral gives
\begin{align}
 \log P={}&-\alpha h+\alpha\lambda/\epsilon
    +\tfrac12\log(2\pi eh)+\frac{2\sqrt2}{3}s^{3/2}\notag\\
 &+\epsilon\left[-\frac{\alpha\lambda^2}{2}
                    +\frac{\sqrt{2s}}2+\frac{s^2}{6}\right]
    +\OO(\epsilon^2T^{5/2}\operatorname{polylog}),\label{xtw:frac:eq:productexp}\\
 -\log v_i={}&-\log\epsilon-\tfrac12\log(2s)
                    +\epsilon\sqrt{2s}/3
                    +\OO(\epsilon^2T\operatorname{polylog}).\label{xtw:frac:eq:gapexp}
\end{align}
The $\epsilon^3$ endpoint shift produces an essential order-$\epsilon$
term. The term $-\alpha\lambda^2/2$ comes from the second term of
$e^{-\lambda\epsilon^2}$ and is equally essential.

\subsection{Completion of the finite-height correction}
The difference at level $i$, its outer product, and the final outer
exponential give $H_h$. That final exponential changes the logarithm by
$\OO(h^{-1})$. Combining \eqref{xtw:frac:eq:innersep},
\eqref{xtw:frac:eq:ratioproduct}--\eqref{xtw:frac:eq:gapexp} cancels the two larger overlap
terms and gives leading factor $h^{1/6}\FF(\lambda)$.
The remaining order-$\epsilon$ coefficient is
\[
 -\frac{\alpha\lambda^2}{2}+\frac{\lambda^2}{2}
 -I_c(T)-\frac{s^2}{3}+\frac56\sqrt{2s}+\frac12U(T).
\]
The $V$ equation gives
\begin{equation}\label{xtw:frac:eq:Ilarge}
 I_c(T)=J_c(\lambda)-s^2/3+(4/3)\sqrt{2s}
                                  +\OO(\operatorname{polylog}/T).
\end{equation}
Together with \eqref{xtw:frac:eq:Ularge}, all growing terms cancel.
The Wronskian identity
\[
 \frac{d}{dt}\frac{\Ai(a(t-\lambda))}{w_\lambda(t)}
          =-\frac{\Ai(-a\lambda)}{w_\lambda(t)^2}
\]
fixes the homogeneous finite part and proves
\begin{equation}\label{xtw:frac:eq:gauge}
 J_c-J_0=ca\frac{\Ai'(-a\lambda)}{\Ai(-a\lambda)},\qquad
 -\FF(J_c-J_0)=-\frac c2\FF'(\lambda).
\end{equation}
Consequently
\begin{equation}\label{xtw:frac:eq:windowcorrection}
 \begin{split}
 h^{-1/6}e^{\alpha h-\alpha\lambda/\epsilon}H_h(e^{\lambda\epsilon^2}/(eh))
  ={}&\FF(\lambda)+\epsilon\left[
       \QQ(\lambda)-\frac\alpha2\lambda^2\FF(\lambda)
                          -\frac{c_h}{2}\FF'(\lambda)\right]\\
   &+\OO(\epsilon^{17/16}\operatorname{polylog}(1/\epsilon)).
 \end{split}
\end{equation}
The logarithmic correction after all cancellations is
$\epsilon\operatorname{polylog}$; exponentiation costs only
$\OO(\epsilon^2\operatorname{polylog})$. This proves an actual uniform
expansion on the growing window, not a formal matching rule.

\section{The Gamma factor on an offset grid}\label{xtw:frac:sec:offset}
We now extend the quantitative proof, not merely its leading equivalent,
to $y$ in a fixed compact $Y\subset(0,\infty)$. Set
\[
 \beta=y-1,\qquad H=h+\beta,\qquad \epsilon=H^{-1/3},\qquad
 x=e^{\lambda\epsilon^2}/(eH),\qquad B_y=e^\beta/\Gamma(y).
\]
At terminal distance $k$, $\ell_{h-k}/H=1-k/H$ exactly. Therefore the
shifted recurrence is exactly \eqref{xtw:frac:eq:recurrence}, with $H$ in place of
$h$, up to its final integer index $h-1=H-y$. Its initial pair is unchanged.
The inner cutoffs, ceilings of the real powers of $\epsilon$, lie far
before $H-y$, uniformly on $Y$. Every proof in
Sections~\ref{xtw:frac:sec:entrance}--\ref{xtw:frac:sec:complex} uses integer $k$ but never
integrality of $H$. Positivity is preserved by $\ell_j>0$. Thus all inner
bounds and the same $c_H$ hold uniformly on $Y$.

The outer grid $z_j=(j+\beta)\eta$, $\eta=e^{\lambda\epsilon^2}/(eH)$,
has spacing exactly $\eta$. At $i=h-M$, its endpoint weight is
$\ell_i=H-M$, so the incoming gap and overlap calculations are identical.
The residual retains its factor $\TT_j$ near the origin. Descending
contraction uses only $|z_j|\le|z_i|$, hence is uniform also when $0<y<1$.
The exact product identity is
\[
 \sum_{j<i}\log(q_j/\TT_j)
       =\sum_{j=2}^i(q_j-\TT_j)+\TT_i-\TT_1.
\]
Near zero, $\TT,R,S=\OO(z)$. The missing lower intervals in the integrals
have length $\OO_Y(|\eta|)$ and contribute $\OO_Y(|\eta|^2)$;
the term $-\TT_1$ and final outer exponential contribute
$\OO_Y(|\eta|)=\OO_Y(\epsilon^3)$, below first order.

The product of the linear factors gives the new constant exactly. Put
$r=i-1$, $R_*=r+\beta=H-M-1$, and $S_*=He^{-\lambda\epsilon^2}$.
Then
\begin{equation}\label{xtw:frac:eq:offsetexact}
 \log\prod_{j=1}^r\TT(\ell_j/(eS_*))
 =\log\Gamma(R_*+1)-\log\Gamma(y)-r\log(eS_*)
                          +\sum_{j=1}^r\TT(\ell_j/(eS_*)).
\end{equation}
Apply the trapezoidal formula on the actual positive grid beginning at
$y\eta$, not a fictitious grid continuing below zero. Its second-derivative
integral is $\OO(1/|v_i|)$, so its remainder is
$\OO(|\eta|/|v_i|)$. The lower-endpoint contribution, for example, is
\[
 \tfrac12\TT(y\eta)-\eta^{-1}\int_0^{y\eta}\TT(z)dz
       =\tfrac12\eta y(1-y)+\OO_Y(|\eta|^2).
\]
Stirling at the real $R_*$ and \eqref{xtw:frac:eq:offsetexact} therefore give
\begin{align}
 \log\prod_{j=1}^r\TT(\ell_j/(eS_*))
  ={}&S_*\int_0^{R_*/S_*}\log\TT(v/e)dv
       +\beta\log(eS_*)-\log\Gamma(y)\notag\\
    &+\tfrac12\log(2\pi R_*)+\tfrac12\TT(R_*/(eS_*))
             +\OO_Y(|\eta|/|v_i|+|\eta|).\label{xtw:frac:eq:offsetEM}
\end{align}
The endpoint $R_*=H-M-1$ is exactly the endpoint used in
\eqref{xtw:frac:eq:productexp}. The only new term is
\begin{equation}\label{xtw:frac:eq:offsetnew}
 \beta\log(eS_*)-\log\Gamma(y)
 =\beta(1+\log H)-\log\Gamma(y)-\beta\lambda\epsilon^2.
\end{equation}
The first two terms give $B_yH^\beta$, and the last is
$\OO_Y(\epsilon^2\operatorname{polylog})$. In particular it is not an
additional order-$\epsilon$ tilt term. The error
$|\eta|/|v_i|=\OO(\epsilon^2/\sqrt T)$ and all lower-endpoint errors
also fit below the $\epsilon^{17/16}$ budget.

We have proved the uniform shifted transform expansion
\begin{equation}\label{xtw:frac:eq:shiftwindow}
 \begin{split}
 B_y^{-1}H^{-\beta-1/6}e^{\alpha H-\alpha\lambda/\epsilon}E_{h,y}(x)
 ={}&\FF(\lambda)+\epsilon\left[
 \QQ(\lambda)-\frac\alpha2\lambda^2\FF(\lambda)
                         -\frac{c_H}{2}\FF'(\lambda)\right]\\
 &+\OO_Y(\epsilon^{17/16}\operatorname{polylog}(1/\epsilon)).
 \end{split}
\end{equation}
This proves the critical leading amplitude as well as its quantitative
correction uniformly on $Y$. No complex-$y$ statement is used.

\section{Fourier tails and the full first order remainder}\label{xtw:frac:sec:minor}
\subsection{Positive factorization on every minor arc}
Write the full and truncated shifted towers at a fixed height as $A_j,B_j$
and $\Delta_j=A_j-B_j$, and put
\[
 D_j(w)=\int_0^1\exp(\ell_jws\Delta_{j+1}(w))ds.
\]
Repeated divided differences give the exact entire identity
\begin{equation}\label{xtw:frac:eq:factorization}
 E_{h,y}(w)=(e^{Hw}-1)w^{h-1}
                  \prod_{j=1}^{h-1}\ell_j B_j(w)D_j(w).
\end{equation}
All nonmonomial factors have nonnegative Taylor coefficients. At
$x_0=(eH)^{-1}$ this implies
\begin{equation}\label{xtw:frac:eq:minorproduct}
 \left|\frac{E_{h,y}(x_0e^{\ii\theta})}{E_{h,y}(x_0)}\right|
 \le\prod_{j=1}^{h-1}\frac{|B_j(x_0e^{\ii\theta})|}{B_j(x_0)}.
\end{equation}
Let $T_0=0$, $T_{L+1}(z)=ze^{T_L(z)}$. The exponent
$\ell_jwB_{j+1}(w)$ dominates $T_{h-j}(\ell_jw)$ coefficientwise,
since every deeper weight is at least $\ell_j$. Hence its contribution
to negative logarithmic modulus is at least
\[
 D_L(z,\theta)=\sum_{k\ge1}[z^k]T_L(z)\,z^k(1-\cos(k\theta)),\qquad
 L=h-j,\quad z=\ell_j/(eH).
\]
For $\tau=\TT(z)$ and $\delta=1-\tau$, the deficit
$d_L=\tau-T_L(z)$ satisfies
\begin{equation}\label{xtw:frac:eq:treedeficit}
 d_L\le(6/L)e^{-\delta L/2},\qquad
 D_L(z,\theta)\ge D_\infty(z,\theta)-2d_L.
\end{equation}
Indeed $d_{j+1}=\tau(1-e^{-d_j})\le d_j-d_j^2/3$ first gives
$d_j\le3/j$; use this through half the steps, then
$d_{j+1}\le\tau d_j$ through the other half.

For $0<|\theta|\le1/8$, choose integer terminal distances
$k=h-j\in[H|\theta|,2H|\theta|]$. On them
$\ell_j/H=1-k/H$ exactly, also for real noninteger $H$.
The degree block $m\in[1/|\theta|,2/|\theta|]$ has
$1-\cos(m\theta)\ge1-\cos1$.
Using $[z^m]\TT(z)=m^{m-1}/m!$ and Stirling, the terms on this block
are bounded below by a constant times $m^{-3/2}$: the extra factor
$(1-k/H)^m$ is bounded below because $k/H\le2|\theta|$.
Thus $D_\infty\ge c_0\sqrt{|\theta|}$.
The gap identity
$k/H=\int_0^\delta v e^v\,dv$ gives
$\delta\asymp\sqrt{|\theta|}$ and
$\delta L\ge cH|\theta|^{3/2}$.
By \eqref{xtw:frac:eq:treedeficit}, the finite deficit is at most half this damping
once $H|\theta|^{3/2}$ exceeds a fixed constant. There are at least
$H|\theta|/2$ selected indices. Consequently
\begin{equation}\label{xtw:frac:eq:minorbound}
 \left|\frac{E_{h,y}(x_0e^{\ii\theta})}{E_{h,y}(x_0)}\right|
 \le e^{-cH|\theta|^{3/2}}\qquad
 A_0H^{-2/3}\le|\theta|\le1/8.
\end{equation}
For $1/8\le|\theta|\le\pi$, use levels with
$\ell_j/H\in[1/4,1/2]$. Their degree-one terms alone give $e^{-cH}$.
All constants are uniform on $Y$, because each chosen range contains the
claimed number of valid levels for large $h$. These far arcs exclude any
other lattice resonance.

\subsection{Polynomial growth of the perturbation finite part}
We justify \eqref{xtw:frac:eq:Jdefinition} and the global integrability required
in inversion. At $t\le1$, the singular subtraction used in
\eqref{xtw:frac:eq:Vbound} gives
$|I_0(1)|+|V_0(1)|\le B(1+|\tau|)^3$.
Set $S_0=\max(S_1,2(1+|\tau|))$ for a fixed sufficiently large $S_1$.
On $[1,S_0]$, stable propagation gives
$|V_0(t)|\le B(1+|\tau|)^3(1+t)$, hence
$|I_0(S_0)|\le B(1+|\tau|)^5$.
On $t\ge S_0$, the sector of $s=t-\ii\tau$ stays strictly inside
$|\arg s|<\pi/3$. The bounded connection ratio in \eqref{xtw:frac:eq:wairy}
ensures the $\Ai$ part is exponentially smaller than the $\Bi$ part,
with absolute constants. Thus \eqref{xtw:frac:eq:Ularge} holds uniformly.
Substitution of
\[
 V_{\rm as}(t)=-2s/3+(2\sqrt2/3)s^{-1/2}
\]
in the $V$ equation leaves residual $\OO(s^{-3/2})$ (leading magnitude
$17/(12\sqrt2)$). Variation of constants, split at $t/2$, gives
\[
 |V_0(t)-V_{\rm as}(t)|\le Bt^{-2}
 +B(1+|\tau|)^4 e^{-c(t^{3/2}-S_0^{3/2})}.
\]
Its integral converges. It proves \eqref{xtw:frac:eq:Ilarge} and, conservatively,
\begin{equation}\label{xtw:frac:eq:Jpolynomial}
 |J_0(\ii\tau)|\le B(1+|\tau|)^5.
\end{equation}
Accordingly $\QQ(\ii\tau)$ is absolutely integrable and has Airy
exponential decay up to a polynomial. The same holds for $\FF'$ and
$\lambda^2\FF$.

\subsection{Full line error and the local limit}
Extend the normalized centered transform
\[
 \Psi_{h,y}(t)=B_y^{-1}H^{-\beta-1/6}e^{\alpha H}
 e^{-\ii\alpha H^{1/3}t}E_{h,y}(e^{\ii tH^{-2/3}}/(eH))
\]
by zero outside $|t|\le\pi H^{2/3}$.
Its normalization at zero is bounded above and below uniformly on $Y$,
by \eqref{xtw:frac:eq:shiftwindow}. Thus \eqref{xtw:frac:eq:minorbound} gives the bound
$Be^{-c|t|^{3/2}}$ outside a fixed compact interval, plus the exponentially
small remote arc. Choose the fixed $A$ in \eqref{xtw:frac:eq:window} large enough
that $cA^{3/2}>1$ with a strict margin for all relevant decay constants.
The tails of the finite transform and $\FF$ then have integral
$o_Y(\epsilon)$. The tails of $\epsilon\QQ$ and
$\epsilon c_H\FF'$ are also $o_Y(\epsilon)$, including
$c_H=\OO(\log H)$.
On the growing central window, the integrated error in
\eqref{xtw:frac:eq:shiftwindow} is still
$\OO_Y(\epsilon^{17/16}\operatorname{polylog})=o_Y(\epsilon)$.
We have therefore proved
\begin{equation}\label{xtw:frac:eq:L1}
 \int_\R\left|\Psi_{h,y}(t)-\FF(\ii t)
 -\epsilon\left[\QQ(\ii t)-\tfrac\alpha2(\ii t)^2\FF(\ii t)
                              -\tfrac{c_H}2\FF'(\ii t)\right]\right|dt
 =o_Y(\epsilon).
\end{equation}
This is the first-order $L^1$ estimate; a fixed Fourier cutoff could not
prove it.

In particular, define the integer Boltzmann size $N_{h,y}$ by
\[
 \Prb(N_{h,y}=n)=c_{n,h;y}x_0^n/E_{h,y}(x_0).
\]
The leading term of \eqref{xtw:frac:eq:L1} implies the uniform local limit
\begin{equation}\label{xtw:frac:eq:llt}
 \sup_{y\in Y}\sup_{n\in\Z}
 \left|H^{2/3}\Prb(N_{h,y}=n)
       -p\left(\frac{n-\alpha H}{H^{2/3}}\right)\right|\longrightarrow0,
 \quad p=f/\FF(0).
\end{equation}
The limiting characteristic function is
$\Ai(0)^2/\Ai(-a\ii t)^2$, and $p(0)=\Ai(0)^2/a$.
The bound follows directly from the $L^1$ error in Fourier inversion;
characteristic convergence and continuity at zero identify the probability
law. For every height, after enlarging a compact-$Y$ constant to cover the
finite initial heights,
\begin{equation}\label{xtw:frac:eq:probmajorant}
 \Prb(N_{h,y}=n)\le B_YH^{-2/3},\qquad
 \frac{c_{n,h;y}}{n!}\le
 B_Y\frac{e^n}{n!}H^{n+\beta-1/2}e^{-\alpha H}.
\end{equation}
For the finitely many initial heights this follows by continuity in $y$:
$H$ is bounded above and away from zero on $Y$, and the finite entire
tower evaluation divided by its stated positive amplitude has a finite
maximum. Thus this is a genuinely all-height majorant.

\section{Height averaging and the universal relative correction}\label{xtw:frac:sec:transfer}\label{xtw:frac:sec:familytransfer}
Let $q$ be the inverse Fourier transform of $\QQ$, a continuous function,
and put $z=(n-\alpha H)/H^{2/3}$. Inverting \eqref{xtw:frac:eq:L1} gives uniformly
in integer $n$ and $y\in Y$
\begin{equation}\label{xtw:frac:eq:coeffwindow}
 \begin{split}
 \frac{c_{n,h;y}}{n!}
  ={}& B_y\frac{e^n}{n!}H^{n+\beta-1/2}e^{-\alpha H}\,
 \left\{ f(z)+H^{-1/3}\left[q(z)-\frac\alpha2f''(z)
                                      -\frac{c_H}{2}zf(z)\right]
                            +o_Y(H^{-1/3})\right\}.
 \end{split}
\end{equation}
Integration by parts sends $\FF'$ to $zf(z)$, while $\lambda^2\FF$
transforms to $f''(z)$; the Airy tails remove all boundary terms.

By \eqref{xtw:frac:eq:probmajorant}, heights outside
$|H-n/\alpha|\le\sqrt n\log n$ are negligible relative to the full sum
by more than any fixed inverse power. To check this uniformly, the logarithm
of the positive weight $H^{n+\beta-1/2}e^{-\alpha H}$ has its maximum at
$(n+\beta-1/2)/\alpha$ and second derivative of order $-1/n$ near that
point. Outside a fixed proportional central interval the exponent loses
$c n$, and in its central complement it loses
$c(H-n/\alpha)^2/n$. Integral comparison on each monotone tail transfers
these bounds to the shifted grid $H=y,y+1,\ldots$.
The retained heights satisfy $h\le n$ for large $n$ because $\alpha>1$.

The full shifted-grid weight sum is
\begin{equation}\label{xtw:frac:eq:gammasum}
 \sum_{h\ge1}(h+\beta)^{n+\beta-1/2}e^{-\alpha(h+\beta)}
 =\frac{\Gamma(n+\beta+1/2)}{\alpha^{n+\beta+1/2}}
                                      [1+\OO_Y(n^{-1/2})].
\end{equation}
Indeed a unimodal positive function differs in integral and grid sum by
at most twice its maximum; the omitted interval $[0,y]$ is harmless.
Uniform Stirling makes the maximum-to-integral ratio $\OO_Y(n^{-1/2})$.
This grid error is below the $n^{-1/3}$ scale.

Let $\E_w$ denote expectation under the normalized weight on the retained
window. Direct Gamma integration and the same sum--integral comparison give
\begin{align}
 \E_wz&=\OO_Y(n^{-2/3}),&
 \E_wz^2&=\alpha^{4/3}n^{-1/3}+o_Y(n^{-1/3}),\label{xtw:frac:eq:moments1}\\
 \E_w|z|^3&=\OO_Y(n^{-1/2}),&
 \E_wH^{-1/3}&=\alpha^{1/3}n^{-1/3}(1+o_Y(1)).\label{xtw:frac:eq:moments2}
\end{align}
For clarity, the continuous first two moments are exactly
\begin{align}
 &\alpha^{2/3}(1/6-\beta)
       \frac{\Gamma(n+\beta-1/6)}{\Gamma(n+\beta+1/2)},\label{xtw:frac:eq:gammafirst}\\
 &\alpha^{4/3}\left[n+\beta^2-\frac23\beta-\frac5{36}\right]
       \frac{\Gamma(n+\beta-5/6)}{\Gamma(n+\beta+1/2)}.\label{xtw:frac:eq:gammasecond}
\end{align}
For the signed first moment, split $z$ times the weight at its zero and
its finitely many extrema. Its maximum absolute value is
$\OO_Y(n^{-1/6})$ times the weight maximum, so its normalized lattice
error is $\OO_Y(n^{-2/3})$. The second-moment error is
$\OO_Y(n^{-5/6})$. Gaussian rescaling gives the absolute third moment;
deleting the already bounded tails changes none of these orders.
This explicitly rules out an order-$n^{-1/3}$ lattice-phase displacement.

Taylor expansion of $f$ through second order, using its bounded third
derivative, yields
\[
 \E_w f(z)=f(0)+\tfrac12\alpha^{4/3}f''(0)n^{-1/3}
                                      +o_Y(n^{-1/3}).
\]
On the shrinking $z$ window, continuity of $q,f''$ gives
\[
 \E_wH^{-1/3}[q(z)-\tfrac\alpha2f''(z)]
 =\alpha^{1/3}[q(0)-\tfrac\alpha2f''(0)]n^{-1/3}
                                      +o_Y(n^{-1/3}).
\]
The two $f''$ terms cancel exactly. The entrance contribution is bounded
in absolute value by $\OO_Y(n^{-1/2}\log n)=o_Y(n^{-1/3})$.
Thus the averaged braces in \eqref{xtw:frac:eq:coeffwindow} are
\begin{equation}\label{xtw:frac:eq:averaged}
 f(0)+\alpha^{1/3}q(0)n^{-1/3}+o_Y(n^{-1/3}).
\end{equation}
Finally, \eqref{xtw:frac:eq:gammasum}, \eqref{xtw:frac:eq:heightsum}, and uniform Stirling give
\[
 B_y\frac{e^n}{n!}
       \frac{\Gamma(n+\beta+1/2)}{\alpha^{n+\beta+1/2}}
 =B_y\alpha^{-\beta-1/2}\rho^{-n}n^{\beta-1/2}
                                                  [1+\OO_Y(1/n)].
\]
Since
\[
 B_y\alpha^{-\beta-1/2}f(0)=C_y,\qquad
 \frac{\alpha^{1/3}q(0)}{f(0)}=\frac{c_1}{C},
\]
this proves the asymptotic parts of Theorems~\ref{xtw:frac:thm:correction} and
\ref{xtw:frac:thm:family}. The sign of $c_1$ is established next.

\section{A real Airy integral for the correction}\label{xtw:frac:sec:signreduction}
The complex perturbation integral is useful for deriving the coefficient;
a different representation is better for certifying its sign.
For $t>0$ define
\begin{align}
 A&=\Ai(at),&B&=\Bi(at),&p&=(\pi/a)A(B+\ii A),\notag\\
 v&=2a\Ai'(at)/A,&
 u&=2a\frac{\Bi'(at)+\ii\Ai'(at)}{B+\ii A},&
 g_t(z)&=\tfrac12-tz/2-z^3/12.\label{xtw:frac:eq:realdefinitions}
\end{align}
All complex logarithms in the following formulas are principal:
\begin{align}
 \Psi&=\tfrac12g_t(v)p^2-(t+v^2/2)p-v\log p+2/(3p),\notag\\
 \Phi&=\tfrac12g_t(u)p^2+(t+u^2/2)p-u\log p-2/(3p),\label{xtw:frac:eq:psiphi}\\
 S(t)&=\frac{\Im\Psi}{\pi A^2},\notag\\
 E(t)&=\frac1\pi\Im\left\{
 \frac{\omega[4\overline\Phi-(4\pi\ii/3)\overline u-2t^2]}{(A+\ii B)^2}
                      \right\},\qquad \omega=e^{2\pi\ii/3},\notag\\
 K(t)&=E(t)-S(t)-\frac{2t}{3a}+\frac1{a\sqrt{2t}}.\label{xtw:frac:eq:realkernel}
\end{align}
Here $K(t)$ denotes a kernel and is unrelated to the entrance index $K$.

\begin{theorem}[Real integral identity]\label{xtw:frac:thm:realintegral}
The integral in the following formula converges absolutely:
\begin{equation}\label{xtw:frac:eq:realconstant}
 c_1=\alpha^{-1/6}\kappa\int_0^\infty K(t)dt.
\end{equation}
The substitution $t=z^2$ gives a continuous integrand at zero,
\begin{equation}\label{xtw:frac:eq:transformedkernel}
 2zK(z^2)=\frac{\sqrt2}{a}
           +2z\left[E(z^2)-S(z^2)-\frac{2z^2}{3a}\right],
 \qquad (2zK(z^2))_{z=0}=2^{5/6}.
\end{equation}
\end{theorem}
\begin{proof}
We include the branch and endpoint controls needed for the contour
operations. Write $s=t-\lambda$, $A_s=\Ai(as)$,
$v_s=2a\Ai'(as)/A_s$, and
\[
 P(t,\lambda)=\frac{A_sw_\lambda(t)}{\Ai(-a\lambda)}
       =\frac\pi a A_s[\Bi(as)-r(\lambda)A_s],\quad
 H(t,\lambda)=P(t,\lambda)g_s(U_\lambda(t)).
\]
The Wronskian gives
\begin{equation}\label{xtw:frac:eq:Greenalgebra}
 U=v_s+2/P,\qquad
 (\partial_t+\partial_\lambda)P=A_s^2/\Ai(-a\lambda)^2.
\end{equation}
Polynomial division of the cubic gives
\[
 H=P g_s(v_s)-s-v_s^2/2-v_s/P-2/(3P^2).
\]
Therefore
\begin{equation}\label{xtw:frac:eq:primitive}
 \Psi_s(P)=\tfrac12g_s(v_s)P^2-(s+v_s^2/2)P-v_s\log P+2/(3P)
\end{equation}
is a primitive of $H$ with respect to $P$, at fixed $s$.

\paragraph{The Green finite part.}
Let $Y=w^2V_0$ and $b(\lambda)=a\Ai'(-a\lambda)/\Ai(-a\lambda)$.
The Wronskian gives
\[
 L(t)=\frac{A_s}{\Ai(-a\lambda)w(t)}=\int_t^\infty w(r)^{-2}dr,
 \qquad L'=-w^{-2}.
\]
Since $Y'=w^2G$, integration by parts yields
$I_0(T)=\FP\int_0^T Hdt-L(T)Y(T)$.
At zero,
\[
 Y=-(2/3)\log t+\OO(t^2),\quad L=t^{-1}+b+\OO(t),\quad
 H=-2/(3t^2)-(2/3)b/t+14\lambda/9+\OO(t).
\]
These expansions define the finite part by adding
$2/(3\delta)-(2/3)b\log\delta$ to the integral from $\delta$.
At infinity, $L(T)Y(T)=P(T)V_0(T)=-\sqrt{2s}/3+o(1)$.
Consequently
\begin{equation}\label{xtw:frac:eq:GreenJ}
 J_0(\lambda)=\lim_{T\to\infty}
 \left[\FP\int_0^TH(t,\lambda)dt+(T-\lambda)^2/3
                                      -\sqrt{2(T-\lambda)}\right].
\end{equation}
In particular, every finite constant is tied to
\eqref{xtw:frac:eq:Idefinition}; none is absorbed into an unspecified finite part.

\paragraph{Principal branches from resolvents.}
Put $C_+(s)=\Bi(as)+\ii\Ai(as)=2e^{\ii\pi/6}\Ai(a\omega s)$
and $P_+(s)=(\pi/a)A_sC_+(s)$; define $P_-$ by conjugation.
For $\Im s<0$, consider $\mathcal L=-d^2/dx^2+x/2$ on $\R$ and
spectral parameter $z=-s/2$.
The right-decaying solution is $\Ai(a(x+s))$, and the left-decaying one
is $C_+(x+s)$. The connection formula and sector expansions show left
decay as a power times $e^{-c_s\sqrt{|x|}}$, $c_s>0$.
Their Wronskian, in the order left then right, is $-a/\pi$.
Thus the diagonal Green function is $P_+(s)$.
Integrating its Green equation against its conjugate gives
\begin{equation}\label{xtw:frac:eq:resolventpositive}
 \Im P_+(s)=\Im z\int_\R|G_s(x,0)|^2dx>0.
\end{equation}
The endpoint terms and weighted integrals converge by the stated decays.
Hence $P_+$ stays in the upper half-plane throughout the lower $s$
half-plane; its logarithm is principal, including both rotation sectors.
Nonvanishing alone would not establish this branch assertion.

Similarly, $P(t,\lambda)$ is the diagonal Green function at $t$ for
$\mathcal L$ on $[0,\infty)$ with Dirichlet boundary at zero and
spectral parameter $\lambda/2$: the Wronskian of $w,A_s$ is
$-\Ai(-a\lambda)$. For $\lambda=\ii\tau$, $\tau>0$, it has positive
imaginary part. The left and right energy identities also give
\begin{equation}\label{xtw:frac:eq:Pbounds}
 \Re(1/P)\ge1/t,\qquad\Im(1/P)<0,\qquad |P|\le t.
\end{equation}
Indeed $1/P=(U-v_s)/2$, the left energy gives $\Re U\ge2/t$,
and the right decaying energy gives $\Re(-v_s)\ge0$.
Thus the actual $P$ is in the first quadrant. All branches in the
primitive are compatible; the lower half-axis is obtained by conjugation.

\paragraph{Normalized Fourier primitive.}
The raw primitive need not decay at both Fourier ends. Normalize each half:
\[
 L_\pm(t,\lambda)=\frac\kappa{A_s^2}
                    [\Psi_s(P(t,\lambda))-\Psi_s(P_\pm(s))],
 \qquad \pm\Im\lambda>0.
\]
Because $(\partial_t+\partial_\lambda)s=0$, \eqref{xtw:frac:eq:Greenalgebra} gives
$(\partial_t+\partial_\lambda)L_\pm=\FF H$.
Let $\inv$ denote upward imaginary-axis inversion at coordinate zero,
$\inv(R)=(2\pi\ii)^{-1}\int_{-\ii\infty}^{\ii\infty}R(\lambda)d\lambda$,
and use $L$ piecewise. The jump at zero is
$L_-(t,0)-L_+(t,0)=2\ii\kappa\Im\Psi_t(P_+(t))/A_t^2$.
It follows that
\begin{equation}\label{xtw:frac:eq:jump}
 \inv(\FF H)=\partial_t\inv(L)+\kappa S(t).
\end{equation}
We next justify its integrated endpoints rather than taking pointwise
limits under the integral.

\paragraph{The origin endpoint uniformly in Fourier frequency.}
For $0<t\le1$, Lemma~\ref{xtw:frac:lem:energy} gives
$|U(t,\ii\tau)|\le B/t+B|\tau|t$, while
$|v_{t-\ii\tau}|\le B(1+|\tau|)^{1/2}$.
Together with \eqref{xtw:frac:eq:Pbounds} these give
\[
 |P|^{-1}\le B(t^{-1}+1+|\tau|),\quad
 |\log P|\le B+|\log t|+\log(1+|\tau|).
\]
Uniformly on $0\le t\le1$, outside a fixed compact frequency interval,
\[
 |A_{t-\ii\tau}|^{-2}\le
 B(1+|\tau|)^{1/2}e^{-c|\tau|^{3/2}}.
\]
Substitution in \eqref{xtw:frac:eq:primitive} bounds the raw primitive divided
by $A_s^2$ by
\begin{equation}\label{xtw:frac:eq:originenv}
 B(t^{-1}+1+|\log t|)(1+|\tau|)^m e^{-c|\tau|^{3/2}}
\end{equation}
for a fixed $m$. The reference primitive has the same kind of bound
without the $t^{-1}$ factor, since $P_+$ and its inverse are polynomially
bounded at infinity and nonzero on the remaining compact set.
Split at $|\tau|=t^{-1/4}$. In the inner part, Taylor expansion gives
\begin{equation}\label{xtw:frac:eq:originexp}
 \frac{\Psi_{t-\lambda}(P(t,\lambda))}{A_{t-\lambda}^2}
 =\frac1{\Ai(-a\lambda)^2}
     \left[\frac2{3t}-2b(1+\log t)\right]
 +\OO\left(\frac{t(1+|\tau|)^m(1+|\log t|)}{|\Ai(-a\ii\tau)|^2}\right).
\end{equation}
For example, use $P=t+bt^2+\OO((1+|\tau|)t^3)$,
$v_s=2b+\OO((1+|\tau|)t)$, and
$A_s^{-2}/A_{-\lambda}^{-2}=1-2bt+\OO((1+|\tau|)t^2)$.
The small parameters $t\sqrt{1+|\tau|}$ and $(1+|\tau|)t^2$
are uniformly small on this interval. The integrated error is
$\OO(t(1+|\log t|))$, using the Airy envelope. On the complementary
interval, \eqref{xtw:frac:eq:originenv} is smaller than every fixed power of $t$
after integration. The reference primitive converges by domination.
Since $\FF'=2b\FF$, $\inv(\FF b)=0$. With $f_0=f(0)=\kappa/a$ and
\[
 k_{\rm ref}=\inv\left(
       \frac{\kappa\Psi_{-\lambda}(P_\pm(-\lambda))}{\Ai(-a\lambda)^2}
                         \right),
\]
we conclude
\begin{equation}\label{xtw:frac:eq:originlimit}
 \inv(L(t,\cdot))=2f_0/(3t)-k_{\rm ref}+o(1).
\end{equation}

\paragraph{The infinity endpoint including remote frequencies.}
Take $\tau\ge0$, $s=T-\ii\tau$, $e(\tau)=r(\ii\tau)+\ii$,
and $\zeta(s)=(\sqrt2/3)s^{3/2}$. Connection identities and the separate
sector expansions give $|e(\tau)|\le Be^{-c\tau^{3/2}}$ and, uniformly
on $-\pi/2\le\arg s\le0$,
\[
\begin{aligned}
 P_+(s)&=(2s)^{-1/2}[1+\OO(|s|^{-3})],\\
 H_+(s)&=P_+(s)g_s(2C_+'(s)/C_+(s))
       =-2s/3+(2s)^{-1/2}+\OO(|s|^{-2}).
\end{aligned}
\]
The exact relation $P=P_+(1-d)$ has
$d=e(\tau)\Ai(as)/C_+(s)$.
The global exponent difference satisfies
\[
 \begin{aligned}
 2\Re[\zeta(T-\ii\tau)-\zeta(-\ii\tau)]
 &=\sqrt2\int_0^T\Re\sqrt{x-\ii\tau}\,dx\\
 &\ge\int_0^T\sqrt{x+\tau}\,dx
 \ge\frac{T}{2\sqrt2}\sqrt{T+\tau}.
 \end{aligned}
\]
Enlarging constants on bounded frequencies therefore gives
$|d|\le Be^{-cT\sqrt{T+\tau}}<1/2$ for all sufficiently large $T$,
uniformly in $\tau$. The principal branches imply
$\log P-\log P_+=\log(1-d)$ with no $2\pi\ii$ term.
Taylor's theorem and
$\Psi_s''(P)=g_s(v_s)+v_s/P^2+4/(3P^3)=\OO(|s|^{3/2})$
give the uniform integrable expansion
\begin{equation}\label{xtw:frac:eq:infinityL}
 L_+(T,\ii\tau)=-\frac{\kappa\pi}{a}e(\tau)H_+(T-\ii\tau)
 +\OO\bigl((T+\tau)|e(\tau)|e^{-cT\sqrt{T+\tau}}\bigr).
\end{equation}
The error integrates to $\OO((T+1)e^{-cT^{3/2}})$.
The remaining nonpolynomial part of $H_++2s/3$ integrates to
$\OO(T^{-1/2})$ against $e(\tau)$. This covers $\tau\gg T$ as well as
bounded $\tau$ without separating exponentially large factors.
For $m_j=\inv(\lambda^j\FF)$, integration by parts on each half-axis gives
\[
 \inv(r-r_\pm)=\frac a{\pi\kappa}m_1,\qquad
 \inv(\lambda(r-r_\pm))=\frac a{2\pi\kappa}m_2,
 \qquad r_+=-\ii,\ r_-=\ii.
\]
The origin boundary terms vanish because they are multiplied by $\lambda$
or $\lambda^2$. Hence
\begin{equation}\label{xtw:frac:eq:infinitylimit}
 \inv(L(T,\cdot))=\frac23Tm_1-\frac13m_2+\OO(T^{-1/2}).
\end{equation}

\paragraph{Fubini and the two finite parts.}
For each fixed $0<\delta<T<\infty$, the energy bounds make $H$ and the
differentiated normalized primitive polynomial in $|\tau|$, while $\FF$
supplies an exponential envelope. The integrations in
\eqref{xtw:frac:eq:jump} may therefore be exchanged before taking limits.
For the upper endpoint in \eqref{xtw:frac:eq:GreenJ}, define its truncated expression
$J_T$. Splitting its $t$ integral at $2(1+|\tau|)$, the energy bounds on
the initial part and the just-proved sector bounds on the remaining part
give $|J_T(\ii\tau)|\le B(1+|\tau|)^m$ uniformly for $T\ge1$.
Indeed on the latter part
$H+2s/3-(2s)^{-1/2}=\OO(|s|^{-2})$ plus a damped connection term;
if $T$ is smaller than the splitting point, its endpoint terms are themselves
polynomial in $|\tau|$. Thus dominated convergence against $\FF$ is valid.
Apply \eqref{xtw:frac:eq:jump}, \eqref{xtw:frac:eq:originlimit}, and
\eqref{xtw:frac:eq:infinitylimit} in \eqref{xtw:frac:eq:GreenJ}. The elementary endpoint
inversion is
\[
 f_0T^2/3-2Tm_1/3+m_2/3-f_0\sqrt{2T}+o(1).
\]
The $Tm_1,m_2$ terms cancel exactly, and $f_0=\kappa/a$ yields
\begin{equation}\label{xtw:frac:eq:invertedJ}
 \inv(\FF J_0)=k_{\rm ref}
 +\kappa\int_0^\infty\left[S(t)+\frac{2t}{3a}
                                  -\frac1{a\sqrt{2t}}\right]dt.
\end{equation}
This computation tracks both finite-part constants explicitly.

\paragraph{Rotation to positive real Airy arguments.}
In the lower $s$ half-plane,
\[
 k_{\rm ref}/\kappa=-\frac1\pi\Im
          \int_0^{-\ii\infty}\frac{\Psi_s(P_+(s))}{\Ai(as)^2}ds.
\]
Rotate from $\arg s=-\pi/2$ to $-2\pi/3$. The sector expansions give
$P_+(s)\asymp s^{-1/2}$, $v_s=\OO(|s|^{1/2})$, and
$\Ai(as)^{-2}=\OO(|s|^{1/2}e^{-c|s|^{3/2}})$ throughout this closed
wedge. The primitive quotient therefore has an exponential envelope times
a fixed polynomial and logarithm. There are no poles: $\Ai$ has no zeros
there and the rotated $\Ai(a\omega s)$ has argument in $[0,\pi/6]$.
There is no singularity at zero. Both connecting arcs vanish.
Equation~\eqref{xtw:frac:eq:resolventpositive} supplies the principal logarithm
throughout, so no winding correction is omitted.
At $s=\omega^2t$, the connection identities give
\begin{align*}
 \Ai(a\omega^2t)&=\tfrac12e^{-\ii\pi/3}(A+\ii B),&
 C_+(\omega^2t)&=2e^{\ii\pi/6}A,\\
 P_+(\omega^2t)&=e^{\ii\pi/3}\overline p,&
 v_{\omega^2t}&=\omega\overline u,\\
 \log P_+(\omega^2t)&=\overline{\log p}+\ii\pi/3,&
 \Psi_{\omega^2t}(P_+)&=\omega[\overline\Phi-(\ii\pi/3)\overline u].
\end{align*}
Here $0<\arg p\le\pi/6$, which also checks the displayed logarithm
directly. Including $ds=\omega^2dt$ gives
\[
 k_{\rm ref}/\kappa=-\frac1\pi\Im\int_0^\infty
 \frac{4\omega[\overline\Phi-(\ii\pi/3)\overline u]}{(A+\ii B)^2}dt.
\]
The same rotation for $m_2/2$ gives
$m_2/(2\kappa)=-(1/\pi)\Im\int_0^\infty
2\omega t^2/(A+\ii B)^2dt$.
Their difference is precisely $\int E$.
Since $q(0)=m_2/2-\inv(\FF J_0)$, \eqref{xtw:frac:eq:invertedJ} proves
\eqref{xtw:frac:eq:realconstant}. The following tail bounds prove its absolute
convergence at infinity; at zero, $E,S$ are finite because their Airy
denominators and $p$ are nonzero, proving \eqref{xtw:frac:eq:transformedkernel}.
\end{proof}

\section{An explicit tail bound and a certified sign}\label{xtw:frac:sec:certificate}
\subsection{A positive real tail inequality}
The following deliberately broad bound suffices to establish a strict sign:
\begin{equation}\label{xtw:frac:eq:tail}
 \int_{10}^\infty|K(t)|dt<0.141.
\end{equation}
Here is a full derivation from bounded Airy remainders. Write
$r=t^{-3/2}$, $\xi=(\sqrt2/3)t^{3/2}$ and define scaled functions by
\[
 \Ai(at)=\frac{e^{-\xi}}{2\sqrt\pi(at)^{1/4}}x,\quad
 \Bi(at)=\frac{e^\xi}{\sqrt\pi(at)^{1/4}}y,\quad
 \Bi'(at)=\frac{(at)^{1/4}e^\xi}{\sqrt\pi}z.
\]
The real remainder inequalities of \cite{dlmf7} imply
\begin{align}
 x&=1-5r/(24\sqrt2)+e_A r^2,& |e_A|&\le385/2304<0.1672,\notag\\
 y&=1+5r/(24\sqrt2)+e_B r^2,& |e_B|&\le4(385/2304)<0.6685,\notag\\
 z&=1-7r/(24\sqrt2)+e_D r^2,& |e_D|&\le3(455/2304)<0.5925.
                                                        \label{xtw:frac:eq:scaledairy}
\end{align}
For $\Bi$ the relevant multiplier is
$\chi(13/6)+1<\chi(3)+1=3\pi/4+1<4$; for $\Bi'$ it is
$\chi(2)+1=3$, where
$\chi(x)=\sqrt\pi\Gamma(x/2+1)/\Gamma(x/2+1/2)$ is increasing.
The $\Ai$ and $\Ai'$ remainders use the first omitted term.
For $t\ge10$, $r<0.032$, and these three scaled functions and the
similarly scaled $-\Ai'$ all belong to $[0.99,1.01]$.

Set $p_0=(\pi/a)\Ai(at)\Bi(at)$ and
$U_B=2a\Bi'(at)/\Bi(at)$. Then
$p_0=xy/\sqrt{2t}$, $U_B=\sqrt{2t}\,z/y$, and
$H_\infty=p_0g_t(U_B)$ satisfies the exact identity
\[
 \frac{H_\infty+2t/3-(2t)^{-1/2}}t
 =h(r,x,y,z)=\frac23-\frac{xz}2-\frac{xz^3}{6y^2}
                                  +\frac{r(xy-2)}{2\sqrt2}.
\]
On the stated box direct differentiation gives
$|h_x|<0.70$, $|h_y|<0.37$, $|h_z|<1.04$.
For the baselines $x_0,y_0,z_0$ obtained by deleting the remainders
in \eqref{xtw:frac:eq:scaledairy}, exact simplification gives
\[
 \frac{h(r,x_0,y_0,z_0)}{r^2}
 =-\frac{1875\sqrt2r^3+44680r^2+128256\sqrt2r+1594368}
             {6912(5\sqrt2r+48)^2}.
\]
Its magnitude is below $0.101$ on $[0,0.032]$: bound the numerator
at $0.032$ and the denominator below by $6912\cdot48^2$.
The mean-value bound is therefore
\[
 |h|/r^2\le0.101+0.70(0.1672)+0.37(0.6685)+1.04(0.5925)
                 =1.081585<1.1.
\]
Thus
\begin{equation}\label{xtw:frac:eq:Htail}
 |H_\infty+2t/3-(2t)^{-1/2}|\le1.1t^{-2}\qquad(t\ge10).
\end{equation}
This is an explicit inequality, not a remainder guessed from a formal series.

To bound the difference from $S$, put $R=\Ai(at)/\Bi(at)>0$.
The real formula for $S$, obtained by taking the imaginary part in
\eqref{xtw:frac:eq:psiphi}, is
\begin{equation}\label{xtw:frac:eq:realS}
 S=\frac1a\left[g_t(v)p_0-t-v^2/2
       -\frac v{p_0}\frac{\arctan R}{R}
       -\frac2{3p_0^2(1+R^2)}\right].
\end{equation}
The Wronskian gives $U_B=v+2/p_0$, and hence exactly
\[
 S-H_\infty/a=\frac1a\left[
 \frac v{p_0}\left(1-\frac{\arctan R}{R}\right)
                  +\frac2{3p_0^2}\frac{R^2}{1+R^2}\right].
\]
Use $|1-\arctan R/R|\le R^2/3$ and the scaled bounds
$p_0\ge0.69/\sqrt t$, $|v|\le1.5\sqrt t$,
$R\le0.511e^{-2\xi}$. The resulting bracket is less than $3t$,
and $3(0.511)^2<1$, so
\begin{equation}\label{xtw:frac:eq:Stail}
 |S-H_\infty/a|<t e^{-(4\sqrt2/3)t^{3/2}},
 \qquad\int_{10}^\infty|S-H_\infty/a|dt<10^{-20}.
\end{equation}
For $E$, the same scaled bounds give
\[
 |p|\le0.74/\sqrt t,\quad |p|^{-1}\le1.45\sqrt t,\quad
 |u|\le1.5\sqrt t,\quad |\log p|\le\log t.
\]
For the last inequality, use
$|\log p|\le-\log|p|+|\arg p|
\le\frac12\log t+\log1.45+10^{-12}<\log t$.
Therefore $|g_t(u)|\le1.05t^{3/2}$ and
$|\Phi|\le(2.83+1.5\log t)\sqrt t$. In particular
\[
 4|\Phi-(\ii\pi/3)u|+2t^2
 \le(17.64+6\log t)\sqrt t+2t^2\le3t^2.
\]
The final inequality follows at $t=10$ from
$17.64+6\log10<10^{3/2}$, and its gap is increasing.
Also $1/(\pi|A+\ii B|^2)\le1.02\sqrt t\,e^{-2\xi}$.
Consequently
\begin{equation}\label{xtw:frac:eq:Etail}
 |E(t)|\le3.1t^{5/2}e^{-(2\sqrt2/3)t^{3/2}},
 \qquad\int_{10}^\infty|E(t)|dt<10^{-9}.
\end{equation}
The logarithmic derivative of this majorant is at most $-4$ for $t\ge10$,
so its integral is at most one fourth of its value at $10$.
Combining \eqref{xtw:frac:eq:Htail}--\eqref{xtw:frac:eq:Etail} gives
\[
 \int_{10}^\infty|K(t)|dt
 \le\frac{1.1}{10a}+10^{-9}+10^{-20}<0.141,
\]
proving \eqref{xtw:frac:eq:tail}.

\subsection{Directed finite interval quadrature}
The supplementary code evaluates \eqref{xtw:frac:eq:transformedkernel} on
$0\le z\le3.2$, so $0\le t\le10.24$. It uses only the Python standard
library. Its arithmetic assumptions are stated explicitly: the documented
correctly rounded Decimal operations \cite{decimal} and the inspected
program are trusted. This is not a proof-assistant formalization, an
independent hardware verification, or a machine proof of the whole
positive-real tail argument.

Intervals have directed lower and upper contexts at 70 decimal digits.
Inputs are exact integers, rational numbers, or finite decimal strings;
ordinary binary floats are rejected. Exp, logarithm, and square root are
rounded to nearest and widened by adjacent representable numbers. General
Decimal power is not used. Integer powers use directed multiplication.
Reciprocals reject intervals containing zero; elementary functions reject
inputs outside their proved domains. The positive arctangent helper rejects
negative arguments, rather than silently invoking an unproved branch.
Constructors and powers reject unsupported usages, including nonfinite
values. Guards use explicit exceptions, so optimization with Python
\texttt{-O} cannot remove them.

The Airy evaluator uses the convergent fundamental series for $f''=xf$,
with starting terms $1$ and $x$. If the current term has degree $n$,
its next term is multiplied by $x^3/((n+3)(n+2))$.
Both value and derivative tails are bounded by geometric series once their
explicit ratios are below one. It uses 64 steps on $0\le x\le9$, covering
$az^2\le a(3.2)^2<9$. The initial values
$\Ai(0)=1/(3^{2/3}\Gamma(2/3))$ and
$\Ai'(0)=-1/(3^{1/3}\Gamma(1/3))$ are enclosed using positive-real
log-Gamma Stirling remainders \cite{dlmfgamma}; the $\Bi$ values follow
from the exact initial-value relations.
The arithmetic starts with Machin's identity
$\pi=16\arctan(1/5)-4\arctan(1/239)$ and exact rational alternating
series remainders. For $x=1/3,2/3$, it evaluates $\log\Gamma(x+100)$
with the eleven terms $B_{2k}/[2k(2k-1)(x+100)^{2k-1}]$,
$1\le k\le11$, following the elementary Stirling main terms.
The error lies between zero and the signed twelfth term. Subtracting
$\sum_{j=0}^{99}\log(x+j)$ gives $\log\Gamma(x)$ before exponentiation.
All Bernoulli numbers are generated as exact rationals.
Endpoint monotonicity on $[0,\infty)$ encloses the Airy functions and
first derivatives over each interval. Second derivatives and their chain
rules follow directly from the Airy differential equation.

Second-order interval jets carry the value, first derivative, and second
derivative of the transformed kernel through all operations.
For an interval $[l,h]$ with exact midpoint $m$, its integral is enclosed by
\begin{equation}\label{xtw:frac:eq:midpointcertificate}
 (h-l)\,[2mK(m^2)]
   +\left[-\frac{(h-l)^3}{24}M_2,\frac{(h-l)^3}{24}M_2\right],
 \qquad M_2\ge\sup_{l\le z\le h}|(2zK(z^2))''|.
\end{equation}
This is Taylor's theorem integrated about the midpoint, not an estimate
from successive floating quadratures. The real formula
\eqref{xtw:frac:eq:realS} avoids extracting an exponentially small imaginary part.
Principal logarithms of first-quadrant numbers are evaluated by
$\frac12\log(x^2+y^2)+\ii\arctan(y/x)$; arctangent is reduced by the
half-angle identity before its alternating series remainder is enclosed.

The deterministic adaptive calculation starts with the 64 intervals of
width $1/20$ covering $[0,16/5]$ and bisects the interval of greatest
enclosure width, with a fixed tie rule, until total width is at most
$0.08$. It produces 562 terminal panels. The mandatory replay regenerates
them, checks their exact rational coverage and shared endpoints, recomputes
every panel, and adds their directed enclosures independently of the
adaptive accumulator. Thus the certificate is reproducible arithmetic,
not a pasted numerical receipt.
The resulting outward-rounded enclosures are
\begin{align}
 \int_0^{10.24}K(t)dt&\in[-0.491128,-0.411312],\label{xtw:frac:eq:corebracket}\\
 \alpha^{-1/6}\kappa&\in[0.9540289883031181,0.9540289883031182].
                                                        \label{xtw:frac:eq:prefactorbracket}
\end{align}
The tail from $10.24$ is bounded by the larger tail from $10$ in
\eqref{xtw:frac:eq:tail}. Using the full endpoints recorded in the replay gives
\begin{equation}\label{xtw:frac:eq:certifiedc1}
 c_1\in[-0.603067859400578,-0.257885944975658]
                    \subset(-0.61,-0.25).
\end{equation}
The enclosure in \eqref{xtw:frac:eq:certifiedc1} is intentionally rounded outward.
It completes the sign part of Theorem~\ref{xtw:frac:thm:correction}. It also proves
$d=c_1/C<0$ for the entire compact-positive-shift family.
A sign certificate for this constant is not an effective error bound for
approximating any particular finite coefficient.

\section{Ancestor choices and local Gamma laws}\label{xtw:frac:sec:ancestor}
Let $\mathcal T_n$ be the labeled trees on $\{0,1,\ldots,n\}$ rooted at
$0$, and let $d_T(v)$ be the distance from the root. The exact identity is
\begin{equation}\label{xtw:frac:eq:treeidentity}
 P_n(y)=\sum_{T\in\mathcal T_n}\prod_{v\ne0}(y+d_T(v)-1).
\end{equation}
Indeed the root has an unordered labeled set of branches. A branch root
contributes $xy$, and its descendants have generating function $F_{y+1}$;
the labeled exponential formula is precisely \eqref{xtw:frac:eq:family}.
The finite tower restricts maximum depth, explaining the exact-height
coefficients in \eqref{xtw:frac:eq:heightsum}. Formula~\eqref{xtw:frac:eq:treeidentity}
also proves that each $P_n$ has nonnegative integer coefficients.

Decorate each nonroot vertex $v$ by choosing one of its $d_T(v)$ strict
ancestors. Give its choice of root $0$ weight $y$, and each other ancestor
choice weight $1$. Summing decorations gives exactly the product in
\eqref{xtw:frac:eq:treeidentity}. Choose uniformly from all such decorated trees at
$y=1$, and let $K_n$ count vertices that chose root $0$. Its probability
generating function is $P_n(y)/P_n(1)$.

The universal relative correction cancels when the two coefficient
expansions are divided. Thus
\begin{equation}\label{xtw:frac:eq:rowpgf}
 \E y^{K_n}=\frac{(n/\rho)^{y-1}}{\Gamma(y)}
                                      [1+o_Y(n^{-1/3})]
\end{equation}
uniformly for compact positive real $y$ sets. For real fixed $t$ take
$y=e^{t/\sqrt{L_n}}$, $L_n=\log(n/\rho)$. Then
\[
 \log\E e^{t(K_n-L_n)/\sqrt{L_n}}
 =-t\sqrt{L_n}+L_n(e^{t/\sqrt{L_n}}-1)
                         -\log\Gamma(e^{t/\sqrt{L_n}})+o(1)
 \longrightarrow t^2/2.
\]
The real MGF continuity theorem proves the normal limit in
\eqref{xtw:frac:eq:pgfclt}. No differentiation of an error term is needed for it.

\subsection{Local complex control and fixed cumulants}\label{xtw:frac:sec:cumulants}
A rate on a real interval plus positivity gives a further local result.
Let $\gamma$ be Euler's constant and let $\zeta$ denote the Riemann zeta
function.
\begin{theorem}[Local Gamma correction to the Poisson transform]\label{xtw:frac:thm:modpoisson}
Uniformly for complex $|z-1|\le1/16$,
\begin{equation}\label{xtw:frac:eq:modpoisson}
 \E z^{K_n}=\frac{e^{L_n(z-1)}}{\Gamma(z)}
                                      [1+\OO(n^{-5/24})].
\end{equation}
For every fixed positive integer $r$, the ordinary cumulant satisfies
\begin{equation}\label{xtw:frac:eq:cumulants}
 \operatorname{cum}_r(K_n)=L_n+\gamma-
 \sum_{k=2}^r\left\{\!\begin{matrix}r\\k\end{matrix}\!\right\}
                 (-1)^k(k-1)!\zeta(k)+\OO_r(n^{-5/24}),
\end{equation}
where the braces are Stirling numbers of the second kind. In particular,
\begin{equation}\label{xtw:frac:eq:meanvariance}
 \E K_n=L_n+\gamma+\OO(n^{-5/24}),\qquad
 \operatorname{Var}K_n=L_n+\gamma-\zeta(2)+\OO(n^{-5/24}).
\end{equation}
\end{theorem}
\begin{proof}
On $|z-1|\le1/4$ define the analytic residual
\[
 R_n(z)=\Gamma(z)e^{L_n(1-z)}P_n(z)/P_n(1)-1.
\]
There are no Gamma poles on this disk. On its real diameter,
\eqref{xtw:frac:eq:rowpgf} gives $|R_n(y)|\le C_0n^{-1/3}$ for sufficiently
large $n$, uniformly in $y\in[3/4,5/4]$.
For the outer complex bound, positivity gives
$|P_n(z)|\le P_n(|z|)$. The leading real family estimate and bounded
compact Gamma factors therefore give
\begin{equation}\label{xtw:frac:eq:outerresidualbound}
 |R_n(z)|\le C_1n^{|z|-\Re z}+1\le C_2n^{1/24},
 \qquad |z-1|\le1/4,
\end{equation}
since $\Re z\ge3/4$, $|\Im z|\le1/4$, and
$|z|-\Re z=(\Im z)^2/(|z|+\Re z)\le1/24$.
In particular a bounded outer normal family is not assumed.

Set $w=4(z-1)$ and work in the upper half unit disk. The harmonic measure
of its semicircular arc is
\[
 v(w)=\frac2\pi\Im\log\frac{1+w}{1-w}.
\]
The quotient has positive real and imaginary parts in the half disk,
so the logarithm is unambiguous. Its boundary values are $0$ on the
open diameter and $1$ on the open arc. At fixed $|w|$, the argument is
maximal at $w=\ii|w|$, because its tangent is
$2\Im w/(1-|w|^2)$. Hence for $|w|\le1/4$,
\[
 0\le v(w)\le(4/\pi)\arctan(1/4)<1/3.
\]
Apply the subharmonic maximum principle to $\log|R_n|$, using the
real-diameter bound and \eqref{xtw:frac:eq:outerresidualbound}. With
$C=\max(C_0,C_2)$ it gives
\[
 \log|R_n(z)|\le\log C+
     \left[-\frac13(1-v(w))+\frac1{24}v(w)\right]\log n
 \le\log C-\frac5{24}\log n.
\]
Zeros of the residual are allowed for this subharmonic function.
At the two corners, continuity of $R_n$ and the stronger endpoint real
bound give nonpositive upper limits for its difference from the harmonic
majorant; alternatively their harmonic measure is zero. Reflection gives
the lower half disk, and the diameter already has the stronger bound.
This proves \eqref{xtw:frac:eq:modpoisson}.

For large $n$, $|R_n|<1/2$ on the smaller disk, so
$\log(1+R_n)$ is analytic and $\OO(n^{-5/24})$ there, normalized to
zero at $z=1$. Gamma also has a normalized analytic logarithm there.
For $|t|\le1/32$,
$|e^t-1|\le e^{1/32}-1\le1/31<1/16$. Therefore
\[
 \log\E e^{tK_n}=L_n(e^t-1)-\log\Gamma(e^t)
                       +\log(1+R_n(e^t))
\]
is the analytic cumulant-generating function, with an analytic uniformly
$\OO(n^{-5/24})$ error. Cauchy's formula bounds its $r$th error derivative
by $r!32^r$ times this supremum for each fixed $r$.
Finally, the Euler-operator identity gives
\[
 \left.\frac{d^r}{dt^r}\log\Gamma(e^t)\right|_{t=0}
 =\sum_{k=1}^r\left\{\!\begin{matrix}r\\k\end{matrix}\!\right\}
                         (\log\Gamma)^{(k)}(1).
\]
The convergent Gamma derivative identities \cite{dlmfgamma} are
$(\log\Gamma)'(1)=-\gamma$ and
$(\log\Gamma)^{(k)}(1)=(-1)^k(k-1)!\zeta(k)$ for $k\ge2$.
Subtracting this contribution proves \eqref{xtw:frac:eq:cumulants} and
\eqref{xtw:frac:eq:meanvariance}.
\end{proof}
For example the third constant is $\gamma-3\zeta(2)+2\zeta(3)$ and
the fourth is $\gamma-7\zeta(2)+12\zeta(3)-6\zeta(4)$.
The exponent $5/24$ is deliberately nonoptimal. This theorem is local near
$z=1$; it is not a global complex-$y$ estimate, a lattice local limit for
$K_n$, control of cumulants with growing order, or an explicit finite-$n$
error certificate. The finite-height local limit in \eqref{xtw:frac:eq:llt}
concerns a different random variable.

\section{The positive real Borel boundary}\label{xtw:frac:sec:borel}
Let $B(z)=\sum_{n\ge0}b_nz^n$. The coefficient theorem implies that its
radius is $\rho$ and, as real $z\uparrow\rho$,
\begin{equation}\label{xtw:frac:eq:borel}
 B(z)=\sqrt{\frac2\rho}(1-z/\rho)^{-1/2}
   +c_1\Gamma(1/6)(1-z/\rho)^{-1/6}
   +o((1-z/\rho)^{-1/6}).
\end{equation}
To prove this, put $x=z/\rho$ and $s=-\log x$. For $0<\beta<1$,
integral comparison for the positive decreasing function
$t^{-\beta}e^{-st}$ gives
\[
 \sum_{n\ge1}n^{-\beta}e^{-sn}
                 =\Gamma(1-\beta)s^{\beta-1}+\OO(1).
\]
Apply this at $\beta=1/2$ and $5/6$. Replacing
$s=(1-x)+\OO((1-x)^2)$ by $1-x$ changes neither displayed term at the
remainder scale. For a sequence $r_n=o(n^{-5/6})$, split the generating
sum at a fixed $N$; its finite prefix is $\OO(1)$ and its tail is bounded
by $\eta\sum n^{-5/6}x^n$ for all sufficiently large $N$. Letting
$\eta\downarrow0$ gives the asserted little-$o$. Finally
$C\Gamma(1/2)=\sqrt{2/\rho}$.
This is a one-sided real Abelian conclusion. No complex sector or
$\Delta$-domain Borel singular expansion is asserted.

\section{Threshold inversion at the fractional scale}\label{xtw:frac:sec:inverse}
Write
\[
 \nu_y=y-\tfrac12,\quad C_y=C\rho^{1-y}/\Gamma(y),\quad d=c_1/C,
 \qquad
 M_y(t)=2t\log\frac{t}{e\sqrt\rho}+\nu_y\log t+\log(2\pi C_y).
\]
The squared Stirling formula and Theorem~\ref{xtw:frac:thm:family} give
\begin{equation}\label{xtw:frac:eq:logfamily}
 \log P_n(y)=M_y(n)+d n^{-1/3}+o_Y(n^{-1/3})
\end{equation}
uniformly on compact positive $Y$. There is an exact useful lower bound:
restrict the root to a single child in \eqref{xtw:frac:eq:treeidentity}, then use
coefficient positivity in the shift to obtain
\begin{equation}\label{xtw:frac:eq:monotone}
 P_n(y)\ge nyP_{n-1}(y+1)\ge nyP_{n-1}(y).
\end{equation}
Thus $P_n(y)$ is strictly increasing once $ny>1$, with a common eventual
index for $y\in Y$.
Let $g_y$ be the piecewise-linear interpolation of $\log P_n(y)$ on that
eventually increasing part. For sufficiently large $L$ define
$z_y(L)=g_y^{-1}(L)$. The finitely many earlier values are bounded
uniformly on $Y$, so they cannot create an earlier crossing of a large
threshold. Hence exactly
\begin{equation}\label{xtw:frac:eq:exactceiling}
 N_y(Y_0):=\min\{n\ge0:P_n(y)\ge Y_0\}=\ceil{z_y(\log Y_0)}
\end{equation}
for sufficiently large $Y_0$.

Let $x_y=x_y(L)$ be the unique large positive solution $M_y(x_y)=L$.
Its derivative is
$M_y'(t)=2\log(t/\sqrt\rho)+\nu_y/t$, positive uniformly on $Y$ for
large $t$. Interpolation of $M_y$ contributes $\OO_Y(1/t)$ because
$M_y''(t)=2/t-\nu_y/t^2$; interpolation of $dt^{-1/3}$ and the uniform
little-$o$ sequence preserves the error in \eqref{xtw:frac:eq:logfamily}.
First bounding $z_y-x_y$ by $\OO_Y(x_y^{-1/3}/\log x_y)$ and then
Taylor expanding therefore proves
\begin{equation}\label{xtw:frac:eq:exactinverse}
 z_y(L)=x_y-\frac{d x_y^{-1/3}}{2\log(x_y/\sqrt\rho)}
                    +o_Y(x_y^{-1/3}/\log x_y).
\end{equation}
Replacing $M_y'$ by the displayed denominator changes only a lower-order
term. Since $d<0$, the fractional term increases the inverse relative to
the smooth leading-log inverse.

\subsection{An explicit common Lambert base}
For the positive real Lambert function $W_0$, put
\begin{equation}\label{xtw:frac:eq:lambert}
 X=\frac{L}{2W_0(L/(2e\sqrt\rho))},\qquad
 D=2\log(X/\sqrt\rho),\qquad
 d_y=\frac{\nu_y\log X+\log(2\pi C_y)}{D},\qquad r_y=X-d_y.
\end{equation}
Here $d_y$ is a Newton displacement, distinct from the universal $d$.
The number $X$ solves $f(X)=L$ for
$f(t)=2t\log(t/(e\sqrt\rho))$; $f'(X)=D$ and $f''(t)=2/t$.
Uniformly on $Y$, $d_y=\OO_Y(1)$, and Taylor's theorem gives explicitly
\[
 M_y(r_y)-L=(d_y^2-\nu_y d_y)/X+\OO_Y(X^{-2}).
\]
The linear cancellation is exact because
$Dd_y=\nu_y\log X+\log(2\pi C_y)$.
It follows that $x_y-r_y=\OO_Y(1/(X\log X))$, which is smaller than
$X^{-1/3}/\log X$. Replacing $x_y$ by $X$ in the fractional correction
also changes only a smaller term. Consequently
\begin{equation}\label{xtw:frac:eq:lambertinverse}
 z_y(L)=r_y-\frac{dX^{-1/3}}{2\log(X/\sqrt\rho)}
                          +o_Y(X^{-1/3}/\log X).
\end{equation}
This accuracy is established by the controlled Newton residual, not
inferred from a coarser leading inverse.
At $y=1$ the expression simplifies to
\begin{equation}\label{xtw:frac:eq:y1inverse}
 z_1(L)=X-\frac14-\frac{\log(8\pi/\sqrt\rho)}{4\log(X/\sqrt\rho)}
       -\frac{(c_1/C)X^{-1/3}}{2\log(X/\sqrt\rho)}
                   +o(X^{-1/3}/\log X).
\end{equation}

Let $R_y(L)$ denote the displayed approximation in
\eqref{xtw:frac:eq:lambertinverse}. For each compact $Y$ there is a positive
$\eta(L)\to0$ such that, uniformly for $y\in Y$,
\begin{equation}\label{xtw:frac:eq:ceilingsandwich}
 \ceil{R_y(L)-\eta(L)X^{-1/3}/\log X}
 \le N_y(e^L)\le
 \ceil{R_y(L)+\eta(L)X^{-1/3}/\log X}.
\end{equation}
It follows from \eqref{xtw:frac:eq:exactceiling} and the uniform real inverse error.
This is a shrinking-width asymptotic sandwich, not a computable error
certificate. A universal single ceiling of $R_y$ is not justified:
thresholds can lie arbitrarily close to integer crossings.
\begin{writenote}[7 October 2026]
\eqref{xtw:frac:eq:exactceiling} is an instance of item~(1) of the transseries volume's \file{p0:thm:staircase}, and \eqref{xtw:frac:eq:ceilingsandwich} follows from it as item~(2) does; $X$ is an instance of \file{p0:prop:factorial-core}; the Newton and fractional corrections are this Part's own (front matter, ``The inverses and the transseries volume''). At $y=1$, \eqref{xtw:frac:eq:y1inverse} is Part~I's $r(Y)$ of \eqref{xtw:airy:eq:Lambertintro} with the fractional term added.
\end{writenote}

\section{Reproducibility and source scope}\label{xtw:frac:sec:reproducibility}
The supplementary archive contains the editable TeX, the PDF, exact
finite checks, the directed interval certificate, regression fixtures,
deterministic receipts, and an offline builder. The mathematical proofs are
in this article. Finite checks support identities and normalizations;
they do not prove an infinite asymptotic statement by testing examples.

The mandatory replay checks polynomial rows through degree 30; independently
enumerates 18,248 Pr\"ufer trees through six nonroot vertices; checks six
positive rational shifts against independent finite-tower expansions;
and checks 17 rational inequalities used after the analytic tail estimates.
A separate arithmetic test suite performs 20,008 deterministic inclusion
and normalization checks. Domain and file-integrity guards have 267 tests.
Every test uses explicit failure conditions that remain active with
Python \texttt{-O}. The certificate regenerates and independently
reaggregates all 562 terminal panels in both modes. The panel verifier
is independent of adaptive scheduling and running-total bookkeeping, but
reuses the interval kernel and panel implementation; it is not a second
implementation of the entire numerical method.

The archive documents exact commands. A code-only run needs Python 3.12
or later and its standard library; it requires no network, special-function
package, or TeX. The full builder additionally needs an installed pdfLaTeX
stack. It checks the source inventory, runs the mandatory replay, compiles
with shell escape disabled and fixed metadata, and publishes only into a
fresh output directory. The ZIP has a sorted explicit inventory and fixed
entry metadata. Byte identity is intended for the same source and installed
Python/TeX stack; cross-version PDF identity is not promised. Existing
outputs, symlinked paths, unexpected files, and altered fixtures are
rejected. No private research paths, timing logs, caches, or downloaded
paper bodies are included.

Hashes establish integrity relative to the received bundle; they do not
authenticate an adversary's replacement of both code and hashes. The
Python implementation, operating system, and source directory are trusted.
The mathematical and arithmetic trust boundaries are distinct: the sign
certificate depends on the proved real integral and tail inequality, and
the interval code does not machine-formalize those proofs. Optional ordinary
floating diagnostics are not part of the mandatory certificate or its inputs.
\begin{writenote}[7 October 2026]
In this report the package's programs are \file{code/193-frac-*.py}, its fixtures and provenance \file{data/193-frac-code-*}, its generated outputs \file{data/193-frac-generated-*}, and its \file{README_REPRODUCIBILITY.md}, \file{DATA_SOURCES.md}, \file{code/README.md} and \file{code/fixtures/README.md} are \file{193-frac-*.md}. \file{Report193.tex} (its text is this Part), the delivered PDF, the package \file{README.md}, \file{SHA256SUMS.json} and the three files \file{generated/certificate/coefficient.json}, \file{core.json} and \file{panels.json}, byte copies of the fixtures \file{code/fixtures/*_reference.json}, are not shipped. The programs find their fixtures and \file{PROVENANCE.json} by their delivered paths; the report's README rebuilds that layout on a copy.
\end{writenote}

\subsection{Attribution and the limit of the literature comparison}
Hanna's OEIS entries explicitly record both the unshifted sequence and the
entire shifted equation \eqref{xtw:frac:eq:family}. The broad continued-exponential
and weighted-tree framework is prior work of Bender and Vinson.
Their 1996 paper and Poland's 1998 continuation remain relevant; full-text
theorem overlap has not been resolved, so no search failure is used as
evidence of novelty. The present claims concern the displayed theorems and
proofs, with those attributions retained.

The retrieved A096542 entry has an indexing inconsistency in one formula
line: it states $[y]P_n(y)=na_n$, whereas its rows and the formal equation
give $[y]P_n(y)=na_{n-1}$. Indeed taking the coefficient of $y$ at $y=0$
in $F_y=e^{xyF_{y+1}}$ gives $[y]F_y=xF_1$.
For example $[y]P_3(y)=15=3a_2$. The reproducible checks use the defining
formal equation and this independently derived indexing, not the
inconsistent line.
\begin{writenote}[7 October 2026]
Part~I records the same mismatch (Section~\ref{xtw:airy:sec:verification}); Remark~\ref{xtw:airy:rem:oeis542} states the corrected formula with its proof and checks it against the entry's rows.
\end{writenote}

The primary analytic references are the standard Airy Wronskians,
connection formulas, sector expansions, real remainder inequalities, and
Gamma identities in the NIST DLMF. The discrete entrance, outer error,
Fourier transfer, Green finite-part reduction, explicit tail estimates,
and local complex continuation argument have been supplied here.
There is no all-orders coefficient theorem, global complex row-PGF theorem,
row-statistic local limit, effective asymptotic onset, or complex Borel-sector
expansion in this report. Such statements require additional analysis.

\subsection{Directions for further analysis}
The estimates suggest several concrete next questions, which remain open here.
\begin{enumerate}
\item \textbf{The next coefficient scale.} Derive and justify the next
modified Riccati equation and outer corrector, retaining any logarithmic
entrance terms until after height averaging. This requires a new overlap
budget and a second-order growing-window Fourier remainder. The present
$\epsilon^{17/16}$ error does not justify a further coefficient.
\item \textbf{Sharper certified constants and finite-index control.}
A finer finite quadrature and additional rigorously bounded Airy tail terms
could narrow the interval for $c_1$. An effective onset for approach from
below would separately require explicit constants throughout the discrete
remainder proof; improving the constant's decimal enclosure alone is
insufficient.
\item \textbf{Varying shifts and broader probability laws.}
The regimes $y\downarrow0$ and $y\to\infty$ require uniform treatment of the
lower endpoint and Gamma factor. A local limit for $K_n$ would require
control of its full Fourier circle, beyond the local complex disk proved
here. The existing finite-height minor arcs concern $N_{h,y}$ and cannot
be transferred to $K_n$ without an additional argument.
\item \textbf{Complex Borel continuation and effective inversion.}
A sectorial Borel expansion needs analytic continuation and complex bounds,
not just the positive-real Abelian estimate. A certified integer-threshold
algorithm would need an effective coefficient remainder and explicit
crossing tests before replacing a shrinking ceiling sandwich by a definite
integer output.
\item \textbf{The full prior-work comparison.}
A page-level comparison with authorized full texts of Bender--Vinson and
Poland would clarify the overlap in unbounded linear weights, critical
finite depth and coefficient asymptotics. Until that comparison is complete,
priority remains open.
\end{enumerate}
\begin{writenote}[7 October 2026]
These directions are gathered with Part~I's questions in Section~\ref{xtw:sec:further}.
\end{writenote}

\section{Further questions and research}\label{xtw:sec:further}

[Added 7 October 2026, batch~110.] This section gathers, under Vladimir's
standing rule of 4~October 2026, the open questions of both manuscripts:
Part~I's five ``Questions not settled by the present estimates''
(Section~\ref{xtw:airy:sec:literature}) and Part~II's five ``Directions for
further analysis'' (Section~\ref{xtw:frac:sec:reproducibility}), with what
Part~II settles of Part~I's list. The delivered lists stay where they are.
No claim of either manuscript was found to be false; the two refuted
statements are OEIS formula lines (Remarks~\ref{xtw:airy:rem:oeis537}
and~\ref{xtw:airy:rem:oeis542}).

\begin{enumerate}
\item\label{xtw:q:rates} \emph{Rates and further coefficients} (Part~I,
item~1; Part~II, item~1). Part~II proves the first correction,
$b_n\rho^n\sqrt n=C+c_1n^{-1/3}+o(n^{-1/3})$
(Theorem~\ref{xtw:frac:thm:correction}), which settles the first-coefficient
part of Part~I's item. Open: an effective rate, and the next coefficient
scale, which needs a further modified Riccati equation and outer corrector, a
new overlap budget and a second-order growing-window Fourier remainder; Part~II
says its $\epsilon^{17/16}$ error ``does not justify a further coefficient''.
\item\label{xtw:q:shifts} \emph{Fixed positive shifts} (Part~I, item~2;
Part~II, item~3). Part~II proves the Gamma prefactor and the same relative
correction uniformly on compact sets of shifts $y>0$
(Theorem~\ref{xtw:frac:thm:family}), which answers Part~I's item for every
fixed positive shift. Open: uniformity as $y\downarrow0$ or $y\to\infty$, and
a local limit for the root-ancestor statistic $K_n$, which would need control
of its full Fourier circle.
\item\label{xtw:q:tilts} \emph{Larger tilt domains and rare heights}
(Part~I, item~3). The real Airy window stops at the convenient bound
$\lambda\le1$; extending it toward the Airy zeros, and large-deviation
equivalents for the height away from the central range, need further
analysis. Open.
\item\label{xtw:q:effective} \emph{Effective inversion, sharper certified
constants and onsets} (Part~I, item~4; Part~II, items~2 and~4). Both inverse
sandwiches are asymptotic existence statements. A certified threshold
algorithm needs explicit constants and onsets for the coefficient error and a
crossing-aware integer comparison; a narrower enclosure of $c_1$ needs finer
quadrature and more rigorously bounded Airy tail terms, and an effective onset
of the approach from below needs explicit constants throughout the discrete
remainder proof. Open.
\item\label{xtw:q:borel} \emph{Complex Borel continuation} (Part~II,
item~4). Part~II's Borel boundary \eqref{xtw:frac:eq:borel} is a one-sided
real Abelian statement; a sectorial expansion needs analytic continuation and
complex bounds. Open.
\item\label{xtw:q:sources} \emph{The full prior-work comparison} (Part~I,
item~5; Part~II, item~5). A page-level comparison with the full texts of
Bender--Vinson and Poland, on unbounded linear weights, critical finite depth
and coefficient asymptotics. Until then neither Part claims priority; the
write did not read those papers either. Open.
\end{enumerate}

\clearpage
\begin{thebibliography}{99}
\bibitem{oeis537}
P. D. Hanna, sequence A096537, \emph{The On-Line Encyclopedia of Integer
Sequences}, created 24 June 2004; includes the amplitude conjecture attributed
to V. Kote\v{s}ovec, 16 December 2019.
\url{https://oeis.org/A096537}. Retrieved 4 October 2026.

\bibitem{oeis542}
P. D. Hanna, triangle A096542, \emph{The On-Line Encyclopedia of Integer
Sequences}, created 25 June 2004.
\url{https://oeis.org/A096542}. Retrieved 4 October 2026.

\bibitem{bv}
C. M. Bender and J. P. Vinson, Summation of power series by continued
exponentials, \emph{Journal of Mathematical Physics} \textbf{37}(8) (1996),
4103--4119. \url{https://doi.org/10.1063/1.531619}.
Primary metadata and abstract inspected; full-text theorem scope unresolved.

\bibitem{poland}
D. Poland, Summation of series in statistical mechanics by continued
exponentials, \emph{Physica A} \textbf{250} (1998), 394--422.
\url{https://doi.org/10.1016/S0378-4371(97)00533-5}.
Primary metadata and abstract inspected; full-text theorem scope unresolved.

\bibitem{chapoton}
F. Chapoton, Free hyperplane arrangements associated to labeled rooted trees
(2003). Preprint arXiv math/0301372, version 2; see Propositions 6 and 7.
\url{https://arxiv.org/abs/math/0301372}.

\bibitem{dlmf2}
NIST Digital Library of Mathematical Functions, Section 9.2,
Airy differential equation, initial values, Wronskian and connection formulas;
especially equations 9.2.3, 9.2.7, 9.2.10 and 9.2.12.
\url{https://dlmf.nist.gov/9.2}.

\bibitem{dlmf7}
NIST Digital Library of Mathematical Functions, Section 9.7,
Airy asymptotic expansions; especially equations 9.7.5--9.7.8, including
the separate function and derivative expansions.
\url{https://dlmf.nist.gov/9.7}.
Report~193 cites the same section for the real and complex error bounds.

\bibitem{dlmf9}
NIST Digital Library of Mathematical Functions, Section 9.9(i),
location of the zeros of the Airy function.
\url{https://dlmf.nist.gov/9.9}.

\bibitem{report192}
\emph{Report192: Airy Limits and the Corrected Asymptotic of A096537},
4 October 2026. Companion leading-result report.
The present report states its leading facts explicitly and supplies the
quantitative estimates required for the sequel.
[Printed as Part~I of this report.]

\bibitem{dlmfgamma}
NIST Digital Library of Mathematical Functions, Sections 5.2, 5.7 and 5.11,
Gamma function, convergent psi series, and real Stirling remainders.
\url{https://dlmf.nist.gov/5.2};
\url{https://dlmf.nist.gov/5.7.E4};
\url{https://dlmf.nist.gov/5.11.ii}.

\bibitem{decimal}
Python Software Foundation, \emph{Decimal fixed-point and floating-point
arithmetic}, Python standard-library documentation.
\url{https://docs.python.org/3/library/decimal.html}.

\bibitem{tsvol} ProveIt, \emph{Transseries: the polynomial-logarithmic calculus, series reversal at infinity, and inversion}, the transseries volume, \file{Analysis/Transseries/docs/series-and-transseries/Transseries_And_Inversion/transseries_and_inversion.tex}; Part~0, \file{p0:thm:staircase}, \file{p0:prop:factorial-core} and \file{p0:prop:operator-form}. Cited by the merge (batch 110).
\end{thebibliography}
\end{document}
